% SAME-FAMILY SOURCE DRAFT. Author: audit_finite57_algorithms, 25 October 2026.
% Higher-dimensional tail: complete exact20 and exact19 cover/proper writing.
% Exact20/19 additions await different-author/root review; separate wholeD19/build pending.
% No workers/imports/enumeration/tests/Lean/TeX jobs. Canonical editions unchanged.
\section{The independent higher-dimensional tail}
\label{app:highertail}
The exact-eighteen argument is separate from the exclusion of higher
ranks. We start this second argument at the published ceiling twenty-nine.
Every remaining dimension is handled by its actual finite forest; no
Coxeter structure is inherited by an arbitrary face. The common geometry,
ordinary and typed costs, actual $A$-chain pressures and generic prefix
are the lemmas already proved above. Here we supply the additional
large-component caps and then join the finite forest cases.

The argument separates at a useful point. An $A$ cap is proved by a
central curve and a plateau transfer. A $B/D$ cap is proved by a long
prefix and the genuinely different short tails. These two geometric
arguments apply to whole components with arbitrary finite spectators.
They reunite only in the finite rank and energy accounting. The
statements below do not replace an actual component by a convenient
abstract completion: every added cap is already an actual wall.

\subsection{Common large-component caps}
\begin{proposition}[Large-component caps]\label{app:tail:prop:caps}
At a farthest vertex, with arbitrary finite spectators, the following
components are impossible:
\[
 A_r\ (21\le r\le29),\qquad B_r\ (23\le r\le29),\qquad
 D_r\ (24\le r\le29).
\]
\end{proposition}

We give their finite continuation and energy proofs below. All references to old directions concern the
chosen farthest vertex, and all added walls are already actual.

\subsection{Clearing spectators and the finite continuation list}
\label{app:tail:sec:capinventory}

Put $c_0=61/100$. For a long index, write $D_i=\omega_i^2$,
$q_i=\sqrt{D_i}$ and $h_i=x\cdot\omega_i$, and call it weak if
$h_i\le q_i-c_0$. An actual edge with
\[
 \frac{\cosh L_i}{q_i}>T(q_i),\qquad
 T(q)=\frac{1+(1-c_0/q)^2}{q[1-(1-c_0/q)^2]}
      =\frac1{c_0}-\frac{q-c_0}{q(2q-c_0)},
\]
is strong. We use $T(q)<100/61$ and, on $5/2\le q\le31/8$,
$T(q)<761/500$. For the latter, clearing denominators gives the
positive concave quadratic
$-715800q^2+3268319q-1860500$; check its two endpoints.

First consider an isolated own contact $\gamma$ with full foreign
cost $Q$. Write $v=\gamma\omega_i+v_S$,
$u=(1-Q)/\gamma^2>0$, so $A^2=\gamma^2(D_i-u)$.
If $u<1$, the positive actual combination
\[
 \xi=c+v_S+\frac{1-Q}{\gamma}n_i
\]
has products zero with $c,S$, product $\gamma(u-1)<0$ with $n_i$,
and square $-\gamma^2u(1-u)$. It is future and feasible, giving
the pure component pressure
\begin{equation}\label{app:tail:eq:clearforeign}
 h_i>\sqrt{D_i-u}-\sqrt{u(1-u)}.
\end{equation}
At $u=1$ the same nonzero feasible vector is null, impossible.
For $u>1$, the actual positive set $S+n_i+c$ has residual
$\gamma^2(u-1)$; the finite-projection formula in \eqref{eq:corepressure} and $x\cdot n_i>0$ give
\begin{equation}\label{app:tail:eq:clearfinite}
 h_i>\sqrt{D_i-u}-\sqrt{u-1}.
\end{equation}
In $A$ the sharper form is $h_i>\sqrt{D_i-1}-\sqrt{u-1}$:
with $B=\sqrt{D_i-1}$, $t=\sqrt{u-1}$ the exact finite pressure
is $\sqrt{B^2-t^2}/(1+t/B)>B-t$.
No normalized projected mirror is asserted native.

The own labels $3,4,6$ and parallel contact are excluded by the
local affine fork/star, $(6,3)$ path, and null sum. Thus
$\gamma\ge\ph/2$. If $Q>0$, $u\le3/\ph^2=6-3\ph<7/6$.
For $A$, $D_i\ge85/11$, $q_i>11/4$, and
$\sqrt{D_i-1}>5/2$. Formula \eqref{app:tail:eq:clearfinite} has loss
$<4/21+5/12=17/28<c_0$, since
$u-1\le5-3\ph<73/500<(5/12)^2$.
Formula \eqref{app:tail:eq:clearforeign} has loss at most
$2/21+\sqrt{457}/42<127/210<c_0$.
For $B/D$, $D_i\ge10$; the respective losses are
$<1/5+2/5$ and
$u/6+\sqrt{u(1-u)}\le(1+\sqrt{37})/12<3/5$.
Thus every foreign isolated cap is strong.

At $Q=0$ the isolated estimates are also strong, uniformly for
$D_i\ge85/11$. For own $5$, the actual old $H_4$ segment has
$\delta_5=(7-3\sqrt5)/16<1/36$ and, by concavity of the old
heights, pressure at most $\gamma h_i$ after its correction.
Its loss is $<25/76+5/24<c_0$; here
$q_i+\sqrt{D_i-25/16}>19/4$ and $\gamma>4/5$.
For own $7$, the adjacent negative $A_2$ core gives loss
$<400/1701+9/25<c_0$, using $.9<\gamma<10/11$.
For finite own $m\ge8$, the old singleton projection gives
$h_i>\sqrt{D_i-1}-\tan(\pi/m)$ and loss
$<4/21+5/12<c_0$.
For dotted own contact, $c+n_i/\gamma$ gives
\eqref{app:tail:eq:clearforeign} with $u=\gamma^{-2}$, and the same
$127/210$ bound. One left $H_4$ segment suffices: its displacement
is $\gamma(3\omega_{i+1}/4+\omega_{i-3}/4)$ and concavity
bounds its scalar product by $\gamma h_i$.

Here is the complete finite continuation table at a central release.
A retained left path $A_{i-1}$ admits endpoint $3$, endpoint $4$,
next-to-end $3$, or an admissible $E_6,E_7,E_8$ position $3$.
Their Schur residuals are
\begin{equation}\label{app:tail:eq:leftresiduals}
 (i+1)/(2i),\quad1/i,\quad2/i,\quad1/4,1/7,1/16.
\end{equation}
For a left $3$ at position $j$ from the far end,
$\sigma=\gamma+j/(2i)$; for endpoint $4$,
$\sigma=\gamma+j/(\sqrt2 i)$. In each case
$\cosh L_i=q_i\sigma/\sqrt{\delta-Q}$.
The complete right table is
\[
\begin{array}{c|l}
\text{retained right part}&\text{finite attachment}\\ \hline
B_s,\ s\ge4&\text{long-leaf }3\\
B_3&\text{long-leaf }3\text{ or short-leaf }3\ (F_4)\\
B_2&3\text{ at either end}\\
D_s,\ s\ge5&\text{long-leaf }3;\ \text{fork }3\text{ only for }s=5,6,7\\
D_4&3\text{ at any leaf}\\
D_3=A_3&\text{middle }3,\ \text{end }3,4,5 .
\end{array}
\]
Two contacts to a retained component create a cycle. A completion
joining both long retained pieces has connected rank at least $21$,
so finite classification forces the corresponding simple endpoint
$A/B/D$ bridge. An additional foreign branch at its internal new
node is not finite. The exceptional fork extensions in the table
are retained and estimated; an $E_7$ piece is not declared sealed
merely because it is exceptional.

Zero own contact needs a separate check. Internal $3$, internal $4$,
right long $3$, and two-piece bridges give the usual affine fork,
star or cycle; a right $D$-fork contact gives affine $E_8$.
A zero-own far $3@1$ in $B/D$ satisfies
$A^2=Q-1/2<1/(2i)\le1/10$; any foreign component then has
derivative square at least $1/4>A^2$, contradicting
Proposition~\ref{prop:shadows}. With $Q=0$ its height is negative.
For $A$, a zero-own far $3$ has height $<1/(2m)$, a next-to-end
$3$ has $Q<2/m\le2/5$ and height $Q-2/N<Q$, and an endpoint
$4$ has $Q<1/m\le1/5$ and negative height. A positive $Q$ in
the first two cases again forces an orthogonal closure. Interior
$A$ contacts retain affine $E_8$ arms $(1,2,5)$ or $E_7$ arms
$(1,3,3)$. This proves positivity of the own contact without
assuming the whole star has no spectators.

For completeness, all other near-position estimates follow directly
from \eqref{app:tail:eq:leftresiduals}. Own $3/4/6$ is affine except the
near $4$/simple-$3$ role. For own $\gamma>101/125$, near simple
$3$ has $\delta\le3/5$, $\sigma\ge\gamma+2/5$;
the next-to-end row has $\delta=2/i$, $\sigma\ge\gamma+1/i$;
the $E$ rows have $\delta\le1/4$.
For $B/D$, the first row exceeds $100/61$ from $i=8$ onward,
and the next-to-end row from $i=7$ onward. At $i=5$ use
$T(q)<38/25$ on $[5/2,10/3]$, and at $i=6,7$ use
$T(q)<8/5$ on $[3,4]$; these are positive concave quadratics
after clearing denominators. White $4$ is stronger, except far
$4$/own $3$, which retains affine $(3,4,3,3)$.
The allowed $B$ near $4$/own $3$ has
$\cosh L_i=(i-1)+i/\sqrt2$ and is strong for $i\ge5$;
in $D$ an own $4$ with no right contact is already affine on the
entire retained right fork. The $A$ near-$4$ role is kept below.

If a nonisolated far $3$ packet has $Q>0$, the actual data in $B/D$ are
\[
 \cosh L_i=\frac{2i\gamma+1}{\sqrt{i+1-2iQ}},\quad
 \frac{\cosh L_i}{q_i}>R_0(i)
 :=\frac{(809/500)i+1}{\sqrt{i(i+2)}}.
\]
$R_0$ increases for $i\ge5$, with derivative sign
$(309/500)i-1$; $T$ also increases for $q\ge\sqrt{10}$.
The three exact comparison rows are
\[
\begin{array}{c|c|c|c}
i&q\text{ upper}&R_0(i)\text{ lower}&T(q)\text{ upper}\\ \hline
5\text{--}8&4&1.536&1.530\\
9\text{--}12&5&1.563&1.550\\
13\text{--}29&7.62&1.577&1.577.
\end{array}
\]
The lower squared margins are $.052740,.321913,.546001$;
the $(T\text{ numerator},T\text{ denominator})$ pairs are
$(27.4921,18.0316)$, $(44.2721,28.6395)$,
$(107.2045,68.003166)$, proving each upper comparison by
multiplication. In $A$, a helper-side parameter $m\ge5$ instead
gives ratio $>(2fm+1)/\sqrt{m(m+2)}\ge227/(25\sqrt{35})>
761/500$, $f=101/125$; the derivative sign is
$(2f-1)m-1>0$, with final squared margin $342365$.
The foreign near-$4$ role has ratio squared $>121/35$.
Thus every remaining weak packet has $Q=0$.

\subsection{The \texorpdfstring{$A_{21}$}{A21} certificate and its plateau transfer}
\label{app:tail:sec:Acap}

Put $N=r+1\in[22,30]$. Then
$D_i=2i(N-i)/N$, $C_{ij}=2i(N-j)/N$ for $i<j$, and
$\|x_A\|^2=\frac12\sum_{i=0}^{N-1}(h_{i+1}-h_i)^2$,
with $h_0=h_N=0$ strictly concave.
For $5\le i\le N-5$, $85/11\le D_i\le15$.
The preceding inventory leaves a near-$4$/simple-$3$ packet or
\begin{equation}\label{app:tail:eq:Aweakpacket}
 v_i=\gamma_i\omega_i+\omega_{e_i}/2,\quad
 e_i\in\{1,N-1\},\quad Q=0,\quad
 f:=101/125<\gamma_i<g:=27/25,\quad A_i>2.
\end{equation}
Allowed own contacts are $5$, finite $m\ge7$, or dotted.
The dotted upper bound follows because at $\gamma=g$,
$\sigma^2/\delta\ge4g-2=58/25>(761/500)^2$;
the difference numerator is $[2m(g-1)-1]^2$.
The near-$4$ edge ratio increases with $m$ and at $m=8$ is
$2\sqrt2/3+7/12>761/500$, so $m\ge8$ is strong.
At $m=7$, $N=22,23$ use
$D_7\le224/23<(25/8)^2$ and $T(q)<3/2$ on $[3,25/8]$,
whose cleared quadratic is $-3400q^2+13237q-7442>0$.
Every weak near-$4$ therefore has a mate at index at most $5$
for $N=22,23$, at most $6$ otherwise, or the reversal.
That mate is actual by its positive $B$ replacement. Its data are
$\delta=1/j$ and
$\cosh L_j=(1/2+(j-1)/(j\sqrt2))\sqrt{jD_j}$.
By reversal take $6\le j\le N-4$; then
$jD_j\ge72(N-6)/N\ge216/5$ and the coefficient is $>13/12$.
Thus $\cosh^2L_j>507/10>49$, giving normalized height
$>\sqrt3/2>43/50$.

We next exclude adjacent weak literal packets. Write $j=i+1$,
$A=\gamma_i$, $B=\gamma_j$, and $R=c_i\cdot c_j$.
Weighted AM--GM of their time coefficients gives
\[
\begin{aligned}
\text{same helper}:\quad R&\ge-AB(1-1/N)\ \pm(A-B)/N\\
 &\qquad +(N-1)/(2N)+(N+1)(A/B+B/A)/(4N),\\
\text{outward helpers}:\quad R&\ge1/2-AB+(A-1)(B-1)/N,\\
\text{reversed helpers}:\quad R&\ge1/2-AB+(A+1)(B+1)/N.
\end{aligned}
\]
The first is $>1-g^2>-1/2$; the others are
$>1/2-g^2-(1-f)(g-1)/22>-7/10>-1/\sqrt2$.
Hence the native gap gives $R=0$ for one helper and
$R\in\{0,-1/2\}$ for opposite helpers.
For outward helpers, $R=0$ is impossible because
$\eta_\ell=\delta_\ell/\sigma_\ell^2<4/5$ and
\[
 (D_i-4/5)(D_j-4/5)-C_{ij}^2
       =(2/5)C_{ij}-24/25+8/(5N)>0.
\]
For reversed helpers, put $l=i$, $s=N-j$, $l+s=N-1$.
The same determinant, multiplied by $N$, is
\[
 4lsA^2B^2-(l+1)(s+1)(A^2+B^2)-2AB
 -(l+1)A-(s+1)B+(N+2)/4.
\]
It increases in $A,B\ge4/5$: the derivative uses
$4lsB^2-(l+1)(s+1)\ge(39/25)ls-N>0$.
If either contact is at least $.9$, concavity in $l$ reduces the
minimum to $l=5$ or $s=5$; at $(4/5,9/10)$ it is at least
$(509N-9974)/500\ge306/125>0$.
The last possibility has both contacts $\ph/2$.
In raw norm-two coordinates their difference squared is
$\ph^2(N-1)/N<4$ for the same helper, so their wall product is
positive. For opposite helpers, $R=0$ would have rational coefficient
$N(N+6)$ or $-N^2$ in the expansion of
$N^2[(v_i\cdot v_j)^2-A_i^2A_j^2]$, according to orientation.
Neither vanishes. This proves $R=-1/2$ for opposite helpers.

The retained two-cap Schur matrix is
$H=I-t\mathbf1\mathbf1^T$, $t=(l-1)/(2l)$, for one helper;
for opposite helpers it has diagonal residuals
$\delta_l,\delta_s\in(1/2,3/5]$ and off-diagonal $-1/2$.
It is positive with positive inverse. Put
$\lambda=H^{-1}(A_i,A_j)^T>0$ and
$\mathcal D=A_i\lambda_i+A_j\lambda_j$.
Its actual mixed point has positive displacement
$Z=W_i\omega_i+W_j\omega_j$, and
\[
 Z^2=\mathcal D(1+\mathcal D),\qquad
 x\cdot Z\ge1+\mathcal D-\sqrt{1+\mathcal D}.
\]
For one helper, $W>(4/5)l(A_i+A_j)$ and
$\mathcal D=[l(A_i+A_j)^2+(A_i-A_j)^2]/2$;
thus $\sqrt{1+\mathcal D}/W<7/16$.
For opposite helpers, $5/6<\lambda_i/\lambda_j<6/5$,
$\mathcal D/(\lambda_i+\lambda_j)^2\le3/55$,
$\lambda_i+\lambda_j>40$, and $W>32$;
the corresponding bound is $<\sqrt{15/176+1/1024}<3/10$.
In either case
\[
 W_iq_i+W_jq_j-\|Z\|
 <\frac{W}{2(\sqrt{85/11}+\sqrt{80/11})}<3W/32.
\]
The weighted mean loss of two weak indices is therefore below
$17/32$ or $63/160$, both below $.61$. Adjacency is impossible.
Every weak near-$4$ can now be removed: its mate has normalized
height $>43/50$; an opposite mate or a strong anchor among
$N-6,N-5$ has height $>39/50$ and correlation at most $1/3$.
For $N\ge24$ the latter follows from $6/(N-6)\le1/3$;
for $N=22,23$ the mate index is at most five and the opposite
anchor at least $16$, giving correlation at most
$30/\sqrt{8160}<1/3$.
Their sum contradicts $(41/25)^2>8/3$.
Thus all weak central indices have \eqref{app:tail:eq:Aweakpacket} and
are nonadjacent.

It remains to exclude weak endpoint $5$. We give the complete fixed
$N=22$ certificate before its exact lift. Nonadjacency and concavity
give strong $6$, a middle floor $h_{11}\ge b$, and a strong $16$
or $17$, where
\[
\begin{gathered}
 a=\sqrt{96/11}-.61,\quad b=\sqrt{120/11}-.61,\quad
 d=\sqrt{85/11}-.61,\\
 2.344<a<2.3443,\qquad 2.6928<b<2.693,\qquad 2.169<d<2.17.
\end{gathered}
\]
If $16$ is strong, the minimum reference $x_0$ through heights
$a,b,a$ at $6,11,16$ has
$E=x_0^2=a^2/6+(b-a)^2/5>47/50$ and positive coefficients
$c_6=c_{16}=(11a-6b)/60>4/25$, $c_{11}=(b-a)/5>0$.
Put $K=5a/6$, $L=a/12$. For a far helper at $21$, each pressure
vector in the table has projection
$(5\gamma/6)\omega_6+\omega_{16}/12$, scalar product
$\gamma K+L$ with $x_0$, and height
$A^2=(85/11)\gamma^2+(5/11)\gamma-23/44$.
\begin{center}\small
\begin{tabular}{ccl}
\toprule
own&pure pressure vector $U$&target $T$\\ \midrule
$5$&$\gamma(\omega_4+\omega_6)/2+\omega_{19}/6$
 &$H(A,2/3-\gamma^2)$\\
$7,8$&$\gamma(\omega_4+\omega_6)/2+\omega_{20}/4$
 &$A-\sqrt{\gamma^2-3/4}$\\
$\ge9$&$\gamma(\omega_4+\omega_6)/2+\omega_{21}/2$
 &$H(A,1-\gamma^2)$\\
dotted&$(\gamma-\gamma^{-1})\omega_5+
 (\omega_4+\omega_6)/(2\gamma)+\omega_{21}/2$
 &$A-\sqrt{1-\gamma^{-2}}$\\ \bottomrule
\end{tabular}\end{center}
Here $x\cdot U\ge T$, $S=\|U-\operatorname{proj}U\|^2$,
and $\Delta=T-\gamma K-L$. The residual variances in row order are
$2\gamma^2/3+1/12$, $2\gamma^2/3+1/6$,
$2\gamma^2/3+5/12$, and $5\gamma^2/3+\gamma^{-2}-19/12$.
The exact bounds and chosen positive multipliers are
\[
\begin{array}{c|ccccc}
\text{own}&5&7&8&\ge9&\text{dotted}\\ \hline
\Delta>&3/10&1/4&9/40&13/50&3/10\\
S<&3/4&3/4&3/4&13/12&5/4\\
\mu&1/5&1/5&1/5&24/125&1/6.
\end{array}
\]
The actual sets are respectively
$J=\{5,21,20\}$, $\{5,21\}$, $\{5\}$, and the dotted pair;
they use the finite-projection formula in \eqref{eq:corepressure} or the future-core formula in \eqref{eq:corepressure}.
The following seals cover the full native ranges, not samples.
For $5$, $\gamma\in(.809,.81)$ gives $A>2.21$,
$\sqrt{2/3-\gamma^2}<1/9$, $\gamma K+L<1.78$.
For $7$, $\gamma\in(.9,.902)$ gives $A>2.47$, radical $<.253$,
and projection $<1.96$.
For $8$, $\gamma\in(.923,.9239)$ gives $A>2.5475$, radical
$<.322$, projection $<2.0003$, so excess $>.2252$.
For finite $m\ge9$, $\gamma>.939$ by the cubic
$8\gamma^3-6\gamma-1$; its deficit increases, since
$A'>11/4$, $H_A>20/23$, $H_\delta<0$, and $K<2$.
At $.939$, $A>2.5917$, $\sqrt{A^2+\delta}>2.6144$,
$\sqrt\delta<.344$, giving $H>2.29$ and projection $<2.03$.
The dotted deficit decreases on $(1,27/25)$ because
$A'<3$, $K>1.9$, and the radical derivative exceeds $15/8$.
At $27/25$, $A>2.99$, radical $<.38$, projection $<2.306$.
Its variance increases to a value below $5/4$.

The vector $x_0+\mu(U-\operatorname{proj}U)$ has all coefficients
nonnegative: the largest old-$6$ removals are $.154,4/25,3/20$,
and the old-$16$ removal is at most $1/60$.
Its Cauchy margin is
$(E+\mu\Delta)^2-E-\mu^2S>0$, increasing in $E$.
The two tight rows have exact lower margins
\[
\begin{aligned}
 (47/50+9/200)^2-47/50-3/100&=9/40000,\\
 (47/50+(24/125)(13/50))^2-47/50-(24/125)^2(13/12)
 &=219/39062500.
\end{aligned}
\]
If $17$ is strong, the reference through $a,b,d$ has energy
$a^2/12+(b-a)^2/10+(b-d)^2/12+d^2/10>24/25$.
Its old coefficients exceed $4/25$; replace the helper projection
by $\omega_{17}/10$. The deficit decreases by less than $1/40$
and variance decreases by $1/60$. The same multipliers still give
positive margins (the tight finite-$8$ margin is already $.01$).
For a helper at $1$, reverse the helper term in every vector.
The projection joins the $\omega_6$ coefficient; its value on $x_0$
is unchanged, while $A^2$ gains $12\gamma/11$.
Since the two time coefficients sum to less than $25/4$,
$A_{\rm left}-A_{\rm right}>192/1375$.
Targets increase by more than $.1$, using $H_A>4/5$.
The variance gains $\gamma/3$. All finite rows then have
$\Delta>13/40,S<3/2$ and work with $\mu=1/6$;
dotted rows have $\Delta>2/5,S<5/3$ and work with $\mu=4/25$.
Their old-$6$ removals are at most $11/72,59/375<4/25$.
This treats both helpers and both strong anchors and excludes weak $5$.

For $N>22$, express each reference or dual vector by its potential
$t_i=W\cdot\omega_i$ and insert $N-22$ copies of $t_{11}$
between its two old halves. Its norm
$\frac12\sum(t_{i+1}-t_i)^2$ is unchanged.
The only split coefficient is the positive peak at $11$, replaced
by nonnegative drops at $11,N-11$. Residual potentials are supported
in the unchanged first and last six indices, so all projection
coefficients and variances above are exactly preserved.
Concavity and nonadjacency give the required middle heights at
$11,N-11$; the old $6,16/17$ floors can only increase.
For the two helpers the time coefficients are
\[
 A_N^2=10\gamma^2-1/2-(10\gamma-1)^2/(2N),\quad
 A_N^2=10\gamma^2+2\gamma-1/2-(10\gamma+1)^2/(2N).
\]
They increase with $N$, as do $H(A,\delta)$ and $A-\sqrt{-\delta}$.
Every lifted positive certificate therefore keeps its strict margin.
Reversal also excludes weak $N-5$.

Finally the $N=22$ full central reference has
\[
 E_{\rm full}=(\sqrt{85/11}-.61)^2/5+
       \sum_{i=5}^{10}(\sqrt{D_{i+1}}-\sqrt{D_i})^2>1007/1000.
\]
Lower/upper pairs for $\sqrt{D_5},\ldots,\sqrt{D_{11}}$ are
\[
\begin{gathered}
(2.7797,2.7799),\quad(2.9541,2.9543),\quad(3.0895,3.0897),\\
(3.1908,3.1910),\quad(3.2612,3.2614),\quad(3.3028,3.3030),\\
(3.3166,3.3167);
\end{gathered}
\]
squaring them gives lower energy $1.007192468$.
For $N\ge23$ lift this same curve with a constant plateau between
$11,N-11$; these two actual heights are at least $\sqrt{11}-.61$
even if weak, since their strong neighbors have dual square at least
$D_{10}(23)=260/23>11$.
For $N\ge23$, only $6,\ldots,10$ and their reversals may now be omitted,
with at most three weak indices per flank. For $N=22$, the weak indices
may instead be any nonadjacent subset of $6,\ldots,16$, including the
middle index $11$; this subset has size at most six.
Removing one isolated floor, including that middle floor, loses
less than $1/1600$, since $q''(t)=-1/q(t)^3$ and $q>11/4$ on
the original central interval. Segment interiors are disjoint;
the surviving potential stays concave with nonnegative duals.
Consequently $\|x_A\|^2>E_{\rm full}-6/1600>4013/4000>1$.
This proves the $A$ cap, with all spectator pressure already cleared.

\subsection{Classical packets and the two short tails}
\label{app:tail:sec:BDcaps}

For $B/D$ put $D_i=2i$ and $S_i=h_i/\sqrt2$.
The continuation table in Section~\ref{app:tail:sec:capinventory} makes every
bulk direction strong except
\begin{equation}\label{app:tail:eq:capclassicalpacket}
 v_i=\gamma_i\omega_i+\omega_1/2,\qquad
 A_i^2=2i\gamma_i^2+2\gamma_i-1/2,\qquad Q=0,
 \quad101/125<\gamma_i<6/5.
\end{equation}
The upper bound follows from
$\sigma^2/\delta\ge4(6/5)-2=14/5>(100/61)^2$ at $\gamma=6/5$;
the difference numerator is $[2i(\gamma-1)-1]^2$.
For a right long $3$ the data are
$\delta=1/2-Q$, $\sigma=\gamma+1/2$; own $0/3$ is affine and
the next contact gives $\sigma/\sqrt\delta\ge1+1/\sqrt2>100/61$.
For a right fork $3$ in $D_s$, $s=5,6,7$,
$\delta=1-s/8-Q$, $\sigma=\gamma+1/4$.
Own $0/3/4/6$ retains affine $E_8$/cycle/$(3,3,4,3)$ path/$(6,3)$;
the next contact gives ratio at least
$(\cos(\pi/5)+1/4)/\sqrt{3/8}>1.72>100/61$.
This includes every possible foreign extension of the finite
$E_6,E_7,E_8$ piece. A two-piece white bridge has residual
$1/(2i)$ and numerator at least $\gamma+1/2+1/(2i)$, hence is
strong; zero own instead gives an affine cycle.

Adjacent packets \eqref{app:tail:eq:capclassicalpacket} are impossible.
Weighted AM--GM gives
\[
 c_i\cdot c_{i+1}\ge1/2-ab+(a/b+b/a)/4
             \ge1-ab>-11/25>-1/2,
\]
so their actual product is zero. Their retained Schur matrix is the
same $I-t\mathbf1\mathbf1^T$, $t=(i-1)/(2i)$, as in
Section~\ref{app:tail:sec:Acap}. Its positive actual mixed point has
$\sqrt{1+\mathcal D}/W<7/16$ and
$W_i\sqrt{2i}+W_{i+1}\sqrt{2i+2}-\|Z\|<W/(4\sqrt{10})$.
The derivation uses $A_i>2$, $\gamma_i>4/5$ and $i\ge5$ only:
the same explicit inverse and time-coefficient formulas apply.
Their sum is below $.61$, proving nonadjacency.

The index-$5$ pressure table and its native scalar seals are given
in Appendix~\ref{app:prefixinventory}. Weak $5$ forces $x_1>.5965$, strong
$6$, and one of $16,17$ strong. The positive prefix interpolation
through $1,6,16$ then has energy
\begin{equation}\label{app:tail:eq:capbootstrap}
 .5965^2+(2.01815-.5965)^2/5+1.55051^2/10
       =1.00043812051>1.
\end{equation}
Endpoint $17$ increases the last term. Thus $5$ is strong using
bulk indices alone, before either short-tail estimate below.

In $B$, the bulk reaches $r-4$. At $i=r-3$, the extra right $B_3$
short-end $3$ gives $F_4$, forbidding every white left or foreign
extension. Own $0$ is affine; otherwise
$\delta=1/4$, $\sigma\ge1/2+1/(2\sqrt2)$, strong.
At $i=r-2$, the extra short-end $3$ has
$\delta=1/2-Q$, $\sigma=\gamma+1/(2\sqrt2)$.
Own $0/4/6$ is affine, while own $\ge5$ is strong.
Own $3$ with $Q>0$ is strong since $\delta\le1/4$.
The only new weak role is therefore
\begin{equation}\label{app:tail:eq:Bterminal}
 Q=0,\qquad v=(\omega_{r-2}+\omega_r)/2.
\end{equation}
If the short-end contact also joins the left path, the full finite
white diagram is $B_r$ and has no foreign extension at $c$.
Its residual $1/(2i)$ and numerator
$\ge\gamma+1/(2i)+1/(2\sqrt2)$ are strong even at own $0$
for $i\ge21$. This is a retained zero-own branch, not an assumed
affine exclusion.
For \eqref{app:tail:eq:Bterminal}, deleting $r-1$ and adding the already
actual $c$ gives positive $A_r$, so the same wall actually closes
$r-1$. Its exact cosh is
\[
 C=\frac{(1+\sqrt2)i+1}{\sqrt{i+3}},\qquad i=r-2.
\]
For $21\le i\le23$, $C/\sqrt{2i+2}>79/50$, $C>9$, and
$\sqrt{2i+2}<7$. The first comparison reduces to
$984i^2-19014i-17498>0$, increasing from $i=21$.
Using $\sqrt{C^2-1}>C-1/C$, its normalized loss is
$<1/(79/50+4/63)<.61$. Thus this mate is strong.

In $D$, the bulk reaches $r-5$. At $i=r-4$ the extra fork $3$
has $\delta=1/2-Q$, $\sigma=\gamma+1/4$.
Its own $0/3/4/6$ is affine as above; $Q>0$ is strong.
At $Q=0$, own finite $\ge7$ or dotted has ratio $>1.61$,
above $T(\sqrt{2i})<1.577$ for $19\le i\le29$.
Only own $5$ needs the averaged projection below.
At $i=r-3$ the full right $D_3=A_3$ list is middle $3$, end $3/4/5$.
Middle $3$ has $\delta=1/2-Q$, $\sigma=\gamma+1/2$ and is
affine or strong. End $4$ has $\delta=1/4-Q$,
$\sigma=\gamma+1/(2\sqrt2)$; own $0$ is affine, otherwise it is
strong. End $5$ is finite $H_4$, forbidding further white contacts,
but own $0$ is retained: already
$\delta=1-3\ph^2/8$ and $\sigma\ge\ph/4$ give ratio $>100/61$.
End $3$ has $\delta=5/8-Q$, $\sigma=\gamma+1/4$.
A left bridge is finite $A_r$, with no further foreign extension;
its residual $1/8+1/(2i)$ and numerator
$\ge\gamma+1/4+1/(2i)$ are strong for $i\ge21$ unless own $0$,
which is an affine cycle. Without that bridge, own $0/3/4/6$ is
affine; $Q>0$ and $\gamma\ge1$ are strong. Precisely own $5$ or
finite own $\ge7$, $\gamma<1$, $Q=0$, remains.

For own $5$ in either short $D$ layer, put $\gamma=\ph/2$ and
$v=\gamma\omega_i+\omega_f/2$ at the contacted fork.
The heights are $A^2=i(\ph+5/8)-1/2$ or
$i(\ph+5/8)-5/8$. Select old paths $i,i-1,i-2$ and
$i,i+1,i+2$; at the $D_3$ layer the last node is the other,
uncontacted fork. Each actual $H_4$ correction is
$w_\pm=\gamma(3n_i/2+n_{\rm side}+n_{\rm second}/2)$,
with $\delta_5=1-3\gamma^2/2$.
The average $Y=v-(w_++w_-)/2$ satisfies
\[
 x\cdot Y>A-\sqrt{\delta_5},\quad
 Y^2=A^2+1-15\gamma^2/8<A^2,\quad
 Y\cdot u_5=\sqrt{10}(\gamma+1/4).
\]
The cross product is $w_+\cdot w_-=3\gamma^2/4$.
For $i\ge19$, $A>6$, $\|Y\|>6.4$,
$\sqrt{\delta_5}<1/7$ (use $\ph>809/500$), and
$\widehat Y\cdot u_5<3.3602/6.4<.55$.
Thus $x\cdot\widehat Y>41/42>.97$; strong $5$ gives
$x\cdot u_5>.8$. Their sum contradicts $1.77^2>3.1$.

For the finite-$\ge7$ $D_3$ role, $i\ge21$ and
$.9<\gamma<1$. Now
$A^2=i(2\gamma^2+\gamma+1/8)-5/8$ and the two actual negative
$A_2$ corrections are $w_\pm=\gamma(4n_i/3+2n_{\rm side}/3)$,
with norm residual $-d$, $d=4\gamma^2/3-1>0$.
Their cross product is $8\gamma^2/9$; the average has
\[
 x\cdot Y\ge A-\sqrt d,\quad
 Y^2=A^2+1-14\gamma^2/9<A^2,
 \quad Y\cdot u_5=\sqrt{10}(\gamma+1/4).
\]
At $i=21$,
$Y^2-40(\gamma+1/4)^2=4\gamma^2/9+\gamma+1/2>0$;
larger $i$ only increases it. Hence correlation is $<1/2$.
The ratio $d/A^2$ increases in $\gamma$ and decreases in $i$;
its derivative numerator is
$[i(4\gamma^2+13\gamma+3)-5\gamma]/3>0$.
Its maximum is $1/195$, giving
$x\cdot\widehat Y>1-1/\sqrt{195}>.928$.
Strong $5$ gives $x\cdot u_5>.807$, so their sum contradicts
$1.735^2=3.010225>3$.
Every correction is positive and feasible against all unknown walls.
Thus the weak inventory \eqref{app:tail:eq:capclassicalpacket} extends
through $r-3$ in $D$, with the single $B$ terminal role retained.

We finish the prefix closure explicitly. Let
$f_j=\sqrt j-.61/\sqrt2$ and
$E_M=f_5^2/5+\sum_{j=5}^{M-1}(\sqrt{j+1}-\sqrt j)^2$.
The concave reference has nonnegative duals. The exact seals are
$E_{21}>1.009$, $E_{20}>.99752$, $f_5/5<.361$:
$f_5^2/5>.6514$ and the increment lower bounds for $j=5,\ldots,20$,
in units $10^{-5}$, are
\[
4545,3846,3333,2941,2631,2380,2173,2000,
1851,1724,1612,1515,1428,1351,1282,1219.
\]
They follow from $(\sqrt{j+1}-\sqrt j)^2>1/(4j+2)$.
One isolated missing floor loses $<1/3200$ by
$|(\sqrt t)''|<1/40$; two consecutive missing floors lose
$<1/800$, by the variance of their three segment slopes.
Disjoint segment interiors make these losses additive.

For a far weak packet at any $5\le i\le21$, the actual pressure
table in Appendix~\ref{app:prefixinventory}, with $q=\sqrt{2i}$, gives
$x_1>\sqrt2[T-\gamma(q-.61)]>.424$.
For finite own $5,7,8,\ge9$, use $A-\gamma q>.079$ and the
four bracket lower bounds
$ .61(.809)+.079-.111$,
$ .61(.900)+.079-.253$,
$ .61(.923)+.079-.323$,
$ .61(.939)+.079-.344$; each times $7/5$ exceeds $.424$.
For dotted $\gamma\in(1,6/5)$ the bracket decreases in both
$\gamma,q$: $A'-q<(q^2+2)/(4q^3)<.1$, while the radical
derivative exceeds $1$. Its endpoint bound is
\[
 \sqrt2[\sqrt{3119/50}-(6/5)\sqrt{42}+.732-\sqrt{11}/6]>.424,
\]
sealed by $7.8981<\sqrt{3119/50}$, $\sqrt{42}<6.480741$,
$\sqrt{11}/6<.552771$, and $\sqrt2>1.4142$.

In $D_r$, $r\ge24$, index $21$ is now in range. If strong,
$E_{21}-8/3200>1$. Otherwise $20$ is strong and $x_1>.424$;
the reference through $20$, with this first-coordinate constraint,
has energy
\begin{equation}\label{app:tail:eq:caprepaired}
 E_{20}+\frac54(.424-.361)^2-7/3200
 >.99752+.00496125-.0021875=1.00029375>1.
\end{equation}
Its slopes still decrease, so all dual coefficients are nonnegative.
For $B_r$, $r\ge26$, one of $21,22$ is strong and the same
$E_{21}$ argument applies. At $r=25$ it still applies through
$r-3$. At $r=24$, choose strong $21$ or $22$, or the strong mate
$23$ of \eqref{app:tail:eq:Bterminal}. At $r=23$, a terminal role at $21$
likewise supplies strong $22$; a far weak $21$ instead uses
\eqref{app:tail:eq:caprepaired}. In these terminal cases there is at most
one two-hole cluster; the excessive loss $9/3200+1/800<.009$
still leaves energy above one. This proves both classical caps and
completes Proposition~\ref{app:tail:prop:caps}.

\subsection{The same central geometry at ranks nineteen and twenty}
\label{app:tail:A19A20}
\begin{lemma}[The actual $N=20,21$ central geometry]
\label{app:tail:lem:Awindow2021}
Let $N\in\{20,21\}$, with an actual $A_{N-1}$ component and
arbitrary finite spectators. On $5\le i\le N-5$, define weak by
$x_A.\omega_i\le\sqrt{D_i}-.61$. Before any contrary energy target,
its complete weak list is a pure near-four cap
$v=\omega_i/\sqrt2+\omega_{i-1}/2$, $i=5,6$ or its reversal,
or a pure far-three helper
$v=\gamma\omega_i+\omega_e/2$, $e=1,N-1$,
$101/125<\gamma<27/25$, with own five, finite at least seven or
dotted. Every weak cap has zero spectator cost. Literal helper
weak indices are nonadjacent, for both helper assignments.
Each near-four has an actual positive finite mate at four/five or
its reversal, with normalized height $>43/50$.
\end{lemma}
\begin{proof}
The central interval is still $5\le i\le N-5$, but its smallest
dual square is now $15/2$, rather than $85/11$. We check the
changes to the inventory before using the preceding two-cap argument.
The exact data are
\[
 D_i\ge15/2,\quad q_i>27/10,\quad q_i<31/8,\quad
 C_{i,i+1}\ge7.
\]
Both retained paths have rank at least four. Thus finite classification
gives precisely the same no-contact, simple-three, endpoint-four and
two-end-three possibilities. A joined finite completion has rank at
least nineteen, so it is an ordinary $A$ bridge, with no extra foreign
branch. The zero-own repair uses only the helper-side length at least
five and the actual affine principals. Its far-three and next-to-far
rows have time below their positive foreign cost and force an
orthogonal closure; the end-four row has no foreign cost and negative
time. Every central own contact is therefore positive.

The isolated-own pressures and selected finite or future points in
Subsection~\ref{app:tail:sec:capinventory} are unchanged. Here
\[
 q_i+\sqrt{D_i-1}>26/5.
\]
For $u>1$ with foreign cost, the loss is less than
$5/26+5/12=95/156<.61$. The same bound holds for pure native
labels at least eight, by $\tan(\pi/8)<5/12$.
For $0<u<1$, the loss is at most
\[
 (5/26)u+\sqrt{u(1-u)}
 \le(5+\sqrt{701})/52<63/104<.61.
\]
The native-seven negative-core loss is less than
$20/81+9/25=1229/2025<.61$: use
$\sqrt{D_i-100/81}>5/2$, whose squared margin is $5/324$.
The native-five $H_4$ loss remains less than
$25/76+5/24<.61$. The affine owns and the null boundary are
unchanged. Thus every isolated cap is strong at these ranks as well.

For the other contact roles the target remains
$T(q)<761/500$ on $[5/2,31/8]$. All ordinary near-three,
next-to-end, exceptional-position and full-bridge estimates above
use only helper-side length at least five. Their foreign-cost
versions only increase the ratios. The same foreign far-three
and near-four comparisons therefore remove every foreign weak role.
The complete weak list is a pure near-four/simple-three cap or
the literal packet~\eqref{app:tail:eq:Aweakpacket}, with
$101/125<\gamma<27/25$. The dotted upper bound uses the
same identity with residual square
$[2m(\gamma-1)-1]^2$ and only $m\ge5$.
Near-four is strong for helper-side $m\ge8$. At $m=7$ the
two actual values of $q_7$ lie in $(3,25/8)$, where the
sharper positive-quadratic seal is $T(q)<3/2$. Thus a weak
near-four has its actual mate at four or five, or a reversal.

We next justify the two-cap argument at these ranks, before removing
near-four. Its mutual-product formulas are identities in $N$.
They still force product zero for a common helper and product zero
or $-1/2$ for opposite helpers. In the outward case the zero-product
determinant is positive because
\[
 (2/5)C_{i,i+1}-24/25+8/(5N)>0.
\]
In the reversed case the displayed determinant polynomial in
Subsection~\ref{app:tail:sec:Acap} has positive derivatives in the
two own parameters. If one own is at least $.9$, its endpoint
minimum is
\[
 (509N-9974)/500\ge206/500>0.
\]
Indeed $l,s\ge5$, $l+s=N-1$ and $ls\ge70$ make the positive
derivative bracket $(39/25)ls-N$ sufficient even after subtracting
$2(27/25)+(l+1)$. If both owns are five, the raw difference
and the nonzero rational coefficients $N(N+6),-N^2$ in the
preceding calculation apply literally. Hence opposite helpers have
product exactly $-1/2$.

The positive two-cap Schur matrices and their actual displacements
are the same as above. Their time coefficients exceed two. To check
this without separating incompatible minima, reverse so $i\le N/2$.
The smaller, far-helper time is exactly
\[
 A_N^2=2i\gamma^2-\tfrac12-
                     {(2i\gamma-1)^2\over2N}.
\]
It increases in $N$ and in $i$ on this half interval, and in
$\gamma\ge.8$. Its minimum at $N=20,i=5,\gamma=.8$ is
$187/40>4$. The near-helper time is larger by
$2\gamma(1-2i/N)\ge0$. We may therefore use the already
printed inverse calculation with $A_i,A_j>2$.
It gives $\sqrt{1+\mathcal D}/W<7/16$ for a common helper
and $<3/10$ for opposite helpers. The spatial rationalization
now gives
\[
 W_iq_i+W_jq_j-\norm Z
 <{W\over2(\sqrt{15/2}+\sqrt7)}<W/10.
\]
Thus the weighted loss of two adjacent weak packets is less than
$43/80$ or $2/5$, both below $.61$. Adjacent literal packets
are impossible, retaining the entire spectator block and every
unknown wall.

The near-four mate is an actual positive finite replacement, with
the same cosine formula printed in the cap proof. Reversing its
numbering if needed gives $6\le j\le N-4$ and
\[
 jD_j\ge72(N-6)/N>216/5,\qquad
 \tfrac12+{j-1\over j\sqrt2}>13/12.
\]
Its cosine is therefore greater than seven, and its normalized
height exceeds $43/50$. None of these steps used an energy target
stronger than the original $e_A<1$.
\end{proof}

\begin{proposition}
\label{app:tail:prop:A19A20}
An actual $A_{19}$ or $A_{20}$ component has projection energy
greater than $6/7$, with arbitrary finite spectators.
\end{proposition}
\begin{proof}
Suppose $e_A\le6/7$ and apply
Lemma~\ref{app:tail:lem:Awindow2021}. At the other end there is another such
mate, or, by nonadjacency, a strong anchor in $\{N-6,N-5\}$.
That anchor has normalized height $>.77$. The correlation is at most
\[
 \sqrt{{5\cdot6\over15\cdot14}}=1/\sqrt7<1/2.
\]
The sum of the two heights is $>1.63$, while the sum direction
has square less than three. This contradicts the assumed energy,
since $1.63^2-3(6/7)=5983/70000>0$. Reversal removes the other
near-four orientation. Only nonadjacent literal helpers remain.

There is now a strong $i\in\{5,6\}$ and a strong
$j\in\{N-6,N-5\}$. Their initial and final coordinate intervals
are disjoint. The exact height-energy identity and Cauchy give
\[
 e_A\ge{h_i^2\over2i}+{h_j^2\over2(N-j)}.
\]
At $N=20$, the endpoint floors are already $q_5-.61>2.1$ and
$q_6-.61>2.28$, by the positive squares
$15/2>2.71^2$ and $42/5>2.89^2$. Increasing $N$ improves them.
Each term is therefore greater than
\[
 \min\{2.1^2/10,2.28^2/12\}=1083/2500.
\]
Consequently $e_A>1083/1250>6/7$, with final difference
$81/8750>0$. This proves the proposition on the same original
projection, with no spectator or face-inheritance restriction.
\end{proof}

\subsection{The one short spectator beside $D_{23}$}
\label{app:tail:D23I}
\begin{lemma}
\label{app:tail:lem:D23I}
The entire forest $D_{23}+I$ cannot occur in dimension twenty-four.
\end{lemma}
\begin{proof}
The bulk inventory and its two-cap proof above apply through eighteen:
the retained right $D$ rank is at least five. This is an explicit
actual-rank application of the inventory, not an extension of the
range $M\le16$ in Proposition~\ref{prop:prefix}. Weak five would
give strong six and a strong sixteen/seventeen. The positive
interpolation in~\eqref{app:tail:eq:capbootstrap} exceeds one,
so five is strong before using any head or short spectator.

At nineteen the right block is $D_4$. Its long-three case is bulk;
both fork-three cases have $(\Delta,\sigma)=(1/2-Q,g+1/4)$.
The own-zero/three/four/six and parallel principals are actual.
Foreign cost and finite own at least seven are directly strong.
The only extra role is pure own five. The two actual $H_4$
projections and their full cross, already printed in
Subsection~\ref{app:tail:sec:BDcaps}, apply at $i=19$ and exclude
it using strong five. Full bridges are finite $D_{23}$ and
foreign-sealed; own zero has its cycle, every positive own is strong.
Thus literal-helper coverage and nonadjacency reach nineteen.
No $D_3$ or two-fork tail is used here.

One of eighteen/nineteen is strong. The exact increments in the cap
proof give $E_{18}>.97119,E_{19}>.98470$. At most six/seven holes
give a legally paid reference $w$ with
\[
 q=\norm w^2>.969315.
\]
It is flat on the first five coordinates, with mean $f_5/5<.361$.
Its remaining individual $P_5$ coefficient is greater than
$f_5/5-(\sqrt6-\sqrt5)>.146$; omitting six only increases it.
Its actual $P_5$ demand is $f_5$. No $I$ energy has been used.

Deleting long node one retains actual $D_{22}$, so finite
classification permits only isolation or a simple three at its long
leaf two. There is only the old $I$ outside this whole component.
The actual head whitelist, selected long initial roots and old $I$
are exactly the local hypotheses of
Lemma~\ref{lem:I-head-flat-reference}. Both signed D forks and
all unknown walls retain nonpositive products under its positive
corrections. That lemma is independent of a $B$ short root and
of the baseline value $q$. Its gain $>.07665$ costs less than
the available $P_5$ price and includes $I$ once. Hence
\[
 \norm{x_D}^2+y^2>.969315+.07665>1.
\]
This proves exactly the whole-forest statement needed here; it makes
no claim about a $D_{23}$ with unspecified further spectators.
\end{proof}

\subsection{The finite accounting in dimensions twenty-four through twenty-nine}
\label{app:tail:24to29}
\begin{proposition}
\label{app:tail:prop:24to29}
No compact hyperbolic Coxeter polytope has dimension at least twenty-four.
\end{proposition}
\begin{proof}
The published ceiling twenty-nine leaves only $24\le n\le29$.
Call $A_r$, $r\ge7$, or $B_r,D_r$, $r\ge9$, dominant.
If no component is dominant, each ordinary rank pays its strict
rank$/14$ cost. The common defect theorem leaves at most one
discarded short rank, or one typed rank-two defect. The respective
total costs exceed $(n-1)/14$ or $.09+(n-2)/14$, both above one.
Thus a dominant block exists. Every dominant block costs more than
$1/2$, so there is exactly one. Write its rank as $r$ and its
spectator rank as $S=n-r$. Only now may every other block be treated
by the small ordinary rates and the single permitted debit.
Their complete three conservative costs are
\[
 S/14,\qquad(S-1)/14,\qquad .09+(S-2)/14.
\]
They represent, respectively, no debit, one short $I$ rank discarded,
and one typed rank-two defect. A $J$ profile pays both ranks and
does not create a worse case. Empty ordinary sums are zero.

For dominant $A$, its half-cost first gives $S\le7$, hence
$r\ge17$. The shared uniform $A_{\ge13}>5/7$ cost then gives
$S\le4$: at $S=5$ the weakest typed case already costs
$5/7+.09+3/14=713/700>1$, while the short case exceeds one.
Thus $r\ge n-4$. For $n\ge25$ this contradicts the cap $r\le20$.
At $n=24$ only $r=20,S=4$ remains. Proposition
\ref{app:tail:prop:A19A20} and the three spectator costs close
it, the weakest short case already giving $6/7+3/14>1$.

For dominant $B/D$, the shared $9/14$ floor first gives $S\le5$.
For $n\ge25$ we have $r\ge20$. If $r=20,21$, the rank-valid
generic prefix gives energy $>.90$, and $S\ge4$; even the
discarded-short cost $3/14$ is too large. If $r=22,23$, its
energy is $>.94$, with $S\ge2$. One remaining ordinary rank
costs $>1/14>.06$; a sole typed rank-two block costs $>.09$.
The caps exclude every larger rank. Thus $25\le n\le29$ fail.

At $n=24$, $19\le r\le23$. The exact generic-prefix rows and
their full original rank hypotheses give
\[
\begin{array}{c|c|c|c}
r&19&20&21\ \hline
e_{B/D}>&.88864&.90715&.9240775
\end{array}
\qquad e_{B/D,\,r=22}>.9401975.
\]
Here $M=13,14,15,16$ respectively, so the $D$ hypothesis
$r\ge M+6$ and the $B$ hypothesis $r\ge M+5$ both hold.
The same increment and omission proofs give these rows; they do
not apply a prefix below its permitted rank. Together with the
three spectator costs at $S=5,4,3,2$, they exceed one.
For $r=23$ the $B$ cap already applies. The remaining $D_{23}$
has one spectator rank: an ordinary $A_1$ costs $>.1$, too much
beside $>.94$; a $J$ or typed defect needs two ranks. The sole
remaining case is the actual isolated $I$ and is exactly
Lemma~\ref{app:tail:lem:D23I}. This exhausts every finite forest,
without asserting that an abstract forest is realized. Dimensions
above twenty-nine use only the named published starting bound.
\end{proof}

\subsection{One isolated-$I$ argument for all the proper large classical rows}
\label{app:tail:classicalI}
\begin{proposition}
\label{app:tail:prop:classicalI}
An entire farthest-vertex forest $B_r+I$, $r\ge20$, or $D_r+I$,
$r\ge21$, is impossible. No further spectators are allowed in this
statement.
\end{proposition}
\begin{proof}
Use $M=15$ in the actual prefix inventory. Its hypotheses are
$B_r$ with $r\ge M+5=20$, or $D_r$ with $r\ge M+6=21$;
its geometric window is exactly $5\le i\le M+1=16$.
Deleting each of those actual long nodes leaves a left path of
rank at least four and a right $B$ of rank at least four, or $D$
of rank at least five. The complete white/foreign list and its
actual affine or point exclusions leave only the pure far-three
helper weak. The actual two-cap argument excludes adjacent weak
helpers. These are precisely the stated hypotheses and domain of
Appendix~\ref{app:prefixinventory}; no $M=17$ or unproved short-tail
continuation is being invoked.

First suppose five weak. Then six is strong, the actual weak-five
point gives $x_1>a_0=.5965$, and one of fifteen/sixteen is strong.
Put $b_0=2.01815<f_6$ and interpolate through
\[
 (1,a_0),\quad(6,b_0),\quad
 (K,b_0+\sqrt K-\sqrt6),\qquad K=15\text{ or }16.
\]
The actual endpoint floor is larger than the displayed surrogate.
The slopes are $a_0$, $.28433$ repeated five times, and
$c_K=1/(\sqrt K+\sqrt6)$ repeated $K-6$ times, then zero.
They decrease strictly. Thus this is a paid nonnegative-anchor
reference $W_0$, with no $I$ coordinate or previously credited
$I$ energy. Its first coordinate is $a_0$ and its retained
$P_6$ demand and pairing both equal $b_0$. Explicitly
\[
\begin{array}{c|c|c|c}
K&\norm{W_0}^2>&c_K<&
 \text{remaining individual }P_6\text{ coefficient}>\\\hline
15&.9851486&.159&.12533\\
16&1.000435&.156&.12833 .
\end{array}
\]
Indeed the common first-six energy is $.7600299945$.
The two tail contributions are greater than $1.4234^2/9$ and
$1.5505^2/10$, using $\sqrt{15}>3.8729$ and $\sqrt6<2.4495$.
The coefficient is $.28433-c_K$, greater than $1/60$ in both rows.
Lemma~\ref{lem:I-capped-initial-graft} applies to this exact paid
reference and adds $>.015$ using the single unused actual $I$ height.
Already the weaker fifteen total exceeds $1.0001486$; the sixteen
row exceeds one even before that graft. Consequently five is strong.

Choose an actual strong $K\in\{15,16\}$ and now use the concave
reference through the actual strong floors, omitting nonadjacent
weak knots. The common exact rational increment bounds give
\[
\begin{array}{c|c|c}
K&E_K\text{ lower}&\text{energy after all permitted holes}>\\\hline
15&.92564&.9240775\\16&.94176&.9401975 .
\end{array}
\]
At most five isolated internal holes cost $5/3200$. The reference
is constant on its first five coordinates, with mean
$b=f_5/5<.361$ and actual $P_5$ demand $f_5$.
Its individual $P_5$ coefficient is at least
$b-(\sqrt6-\sqrt5)>.36-.214=.146$; omitting six only increases
that coefficient. All other anchor coefficients remain nonnegative.
No $I$ energy has been used by this reference either.

The head deletion retains $B_{r-1}$ or $D_{r-1}$ of rank at least
nineteen or twenty. Finite classification permits exactly isolation
or a simple three at the retained long leaf two. A short-end
exceptional attachment is too small at these ranks. The entire
forest hypothesis leaves just the old $I$ as a possible foreign
factor. In the nonisolated row, own zero has zero time or the
positive actual null $c+\Omega_2/2$; own three has its triangle.
The other own ranges and every isolated foreign/native boundary are
exactly the local head inventory of
Lemmas~\ref{lem:I-finite-head-flat-reference}
and~\ref{lem:I-head-flat-reference}.

Their selected long roots $n_1,n_2,n_3$ and actual old $I$ have
the same physical Gram and coordinates here. Their finite projections
or future cores have positive actual-wall coefficients, retain all
unknown walls, and retain both signed forks in type $D$. Their
joint point covers every actual head/$I$ mutual parameter; it does
not presume a finite pair Gram. They do not use a terminal $B$ short
normal, a positive final $D$ coordinate or a tail conclusion.
The uniform gain is therefore $>.07665$ on this paid flat-five
reference, including $I$ exactly once. Its price is at most $.1$,
below the available $.146$. Hence
\[
 \norm{x_{B/D}}^2+y^2>.9240775+.07665=1.0007275>1.
\]
This proves the proposition with precisely its entire-forest scope.
In particular it does not exclude $D_{20}+I$, whose $M=15$ rank
hypothesis fails, or supply a theorem with arbitrary extra spectators.
\end{proof}

\subsection{One smaller-loss window for whole $B_{22},D_{22},D_{23}$}
\label{app:tail:whole58}
\begin{lemma}[A local whole-diagram $.58$ window]
\label{app:tail:lem:whole58local}
Let an entire actual farthest forest be $B_r$ or $D_r$, $r\ge19$.
On $5\le i\le\min\{20,r-4\}$ in $B$, or
$5\le i\le\min\{20,r-5\}$ in $D$, use the loss $.58$.
The complete weak list is the pure far-three helper, with own five,
finite at least seven or dotted and $\phi/2\le g<6/5$, except
that isolated dotted own at five is also permitted.
Literal helper weak indices are nonadjacent. Their actual first
coordinate bounds are $>.5455$ at five, $>.378$ through twenty,
and $>.45$ through nine. No short-tail continuation or whole-rank
exclusion is part of this local statement.
\end{lemma}
\begin{proof}
This time the whole-diagram hypothesis is essential. There is no
foreign inverse-Gram cost. Set $c_0=29/50$, $F_i=\sqrt i-c_0/\sqrt2$,
and call a long index $c_0$-strong when $S_i>F_i$. The old $.61$
floors are still denoted by $f_i$. We use actual unit roots, the
genuine short root in $B$, and both signed forks in $D$.

\paragraph{The new inventory is checked locally.}
The sufficient ratio for $c_0$-strength is
\[
 T_0(q)={50\over29}-{q-29/50\over q(2q-29/50)}.
\]
It increases for $q\ge\sqrt{10}$. Positive rational multiplication
at the right endpoints gives
\[
 T_0(q)<\begin{cases}
 83/50&q<7,\\801/500&q<15/4,\\159/100&q<16/5,\\33/20&q<32/5.
 \end{cases}
\]
The local bulk conditions are exactly a right $B$ of rank at least
four, or right $D$ of rank at least five. At the stated ranks this
gives $5..18$ in $B_{22}$, $5..17$ in
$D_{22}$, and $5..18$ in $D_{23}$. Each deletion leaves $A_{i-1}$
and a right $B_s$, $s\ge4$, or $D_s$, $s\ge5$. The complete
one-node finite-white list is the one proved above: left endpoints,
next-to-end and exceptional positions, right long leaf, the finite
$D$ fork extensions, and full simple endpoint bridges. No generic
prefix theorem is being invoked with $M>16$.

For clarity, at the additional $D_{23}$ index eighteen the retained
blocks are exactly $A_{17}+D_5$. A right long-leaf three is the
ordinary $D$ extension. A right fork three gives finite $E_6$,
with $(\Delta,\sigma)=(3/8,g+1/4)$. Own zero has the actual
affine $E_8$ with arms $1,2,5$; own three has its all-three cycle;
own four has the actual $3,3,4,3$ principal on the cap, the
contacted fork and the three old nodes ending at eighteen;
own six has the $6,3$ path; parallel has the pair null.
Every other own has ratio $>(.8+.25)/\sqrt{3/8}>83/50$.
A long-left/fork bridge with those $E_6$ arms cannot be finite at
rank twenty-three. The permitted full bridge instead meets the
long leaf and is finite $D_{23}$, with the same endpoint ratios.
This supplies the local endpoint argument at eighteen explicitly.

The complete $.58$ scalar checks are as follows. Near-three has
least ratio $[i(\phi+1)-1]/\sqrt{2i(i+1)}$. It increases with $i$.
At $i=6$ it exceeds $801/500$, handling six/seven; at eight it
exceeds $83/50$, handling every later index. At five, own at least
seven gives ratio $>5/3$. Thus only own-five/near-three at five is
new. Next-to-end positions have ratio at least
$(i\phi/2+1)/\sqrt{2i}$, above $159/100$ at five and $83/50$
from six onward. Exceptional positions have residual at most $1/4$,
$j\ge3,i\le8$, hence ratio at least $2g+3/8>83/50$.
The far-four/own-three principal is affine. Near-four/own-three
has actual cosine $(i-1)+i/\sqrt2$; its other owns are stronger.
Right long-three and the short-three $B_3$ exception, when used
later as a separate short-layer check, have ratio
at least $1+1/\sqrt2>83/50$. A $D_s$ fork, $s\ge5$, has the
same affine own list just checked. The only finite fork extensions
have $s=5,6,7$, with residual $1-s/8\le3/8$.
Own zero uses actual affine arms $1,2,5$ by extending the retained
long arm with at most the old $i,i-1,i-2$; all exist because $i\ge5$.
Own three closes the all-three cycle, own four uses the
$3,3,4,3$ path on $i-2,i-1,i,c,f$, and own six/parallel use
the local path/null. Every other own has remaining ratio
$>(.8+.25)/\sqrt{3/8}>83/50$. Full bridges have smaller residuals
and larger numerators, with their own-zero cycles retained.

The isolated own-five $H_4$ loss is $<25/76+5/24<.58$.
For seven, the old negative core gives loss
$<1000/4779+9/25<.57$: here $q+\sqrt{D-100/81}>3+2.9$.
For finite owns at least eight, the loss is largest at $D=10,m=8$,
where it is $\sqrt{10}+\sqrt2-4<.578$. The monotonic singleton
projection proves the whole remaining finite interval. The old
owns three/four/six and parallel retain their actual affine/null
principals. For isolated dotted own, $u=g^{-2}\in(0,1)$, the
actual core gives loss at most
\[
 {u\over\sqrt{2i}+\sqrt{2i-1}}+\sqrt{u(1-u)}.
\]
At $i\ge6$ the first coefficient is $<37/250$, so this is
$<(37/250+1.011)/2=.5795$. Only isolated dotted five remains.
Thus, apart from that row and near-five, every weak bulk cap is
the pure helper $v=g\omega_i+\omega_1/2$, with own five, finite
at least seven or dotted.

\paragraph{The new near-five role.}
For $v=(\phi/2)\omega_5+\omega_4/2$, put $g=\phi/2$ and use
the actual right $H_4$ correction
\[
 w=g(3n_5/2+n_6+n_7/2),\qquad
 \delta=(5-3\phi)/8,\qquad Y=v-w.
\]
Its extra old-four product is negative. Its normalized finite
projection retains every original wall. Direct substitution gives
\[
 A^2=(7+13\phi)/2,\quad Y^2=(33+49\phi)/8,\quad
 \sum Y_j=\sqrt2(2+5\phi/2).
\]
From $1.618<\phi<1.62$, $A>3.74$ and $\sqrt\delta<.136$,
we have $x.Y>3.6,Y^2<14.05,\sum Y_j<8.6$.
The independently valid old $.61$ inventory at $M=11$ supplies
a strong $K=11$ or twelve, before any new strong-five conclusion.
For $Z=Y+(3/4)P_K$ the two demand/square bounds are
\[
\begin{array}{c|c|c|c}
K&x.Z>&Z^2<&\text{squared-demand excess}>\\\hline
11&5.76&33.1375&.0401\\12&5.8725&33.7&.78625625 .
\end{array}
\]
Here $f_{11}>2.88,f_{12}>3.03$. Cauchy excludes both, so the
new near-five role is impossible without a $.58$ strong-five premise.

\paragraph{Helper nonadjacency and three actual consequences.}
Every weak helper has $\phi/2\le g<6/5$. The upper bound is
reproved at this threshold: at $g\ge6/5$ the squared ratio is at
least $4g-2\ge14/5>(83/50)^2$. Adjacent helper caps are distinct
actual walls, and their product is at least
$1-ab>-11/25>-1/2$. It must be zero. Their positive Schur inverse
and actual mixed point give loss
$<7/16+1/(4\sqrt{10})<.58$, using only $i\ge5,g>.8,A>2$.
Thus new weak helpers are nonadjacent; we have not presumed they
were already old $.61$-weak.
This last two-cap calculation uses only the two displayed pure-helper
profiles, their actual long duals and the retained positive old blocks.
A later short-layer classification may supply those same profiles;
after checking that physical interface it may reuse this calculation,
without extending the present classification domain.

The actual helper point table gives
$x_1>\sqrt2[T-g(\sqrt{2i}-.58)]$. Its complete native ranges imply
\[
\begin{array}{c|c}
i=5\text{ weak helper}&x_1>.5455\\
i\le20\text{ weak helper}&x_1>.378\\
i\le9\text{ weak helper}&x_1>.45 .
\end{array}
\]
For the first subtract at most $.03\sqrt2(6/5)$ from the old
scalar expression $>.5965$. For the second, the four finite lower
brackets are $.58(.809)+.079-.111$,
$.58(.900)+.079-.253$, $.58(.923)+.079-.323$ and
$.58(.939)+.079-.344$. Their products with $\sqrt2$ exceed $.378$.
The dotted difference decreases in $g,q$ and has enlarged endpoint
\[
 \sqrt{119/2}-(6/5)\sqrt{40}+.696-\sqrt{11}/6.
\]
The bounds $\sqrt{119/2}>7.71362,\sqrt{40}<6.32456$,
$\sqrt{11}/6<.55278,\sqrt2>1.4142$ give $>.3781118$.
For $i\le9$ the finite $A-gq$ term is $>.15$, increasing in $g$
and decreasing in $q$; the same brackets then exceed $.45$.
For dotted use $\sqrt{1391/50}>5.274,\sqrt{18}<4.243$,
$\sqrt{11}/6<.553,\sqrt2>1.4$, giving $>.45556$.
All these pressures use actual positive finite projections or future
cores. At a last $D$ long node, $\omega_{21}+\omega_{22}=\sqrt2P_{21}$
in $D_{22}$ is only a spatial identity; no extra normal is introduced.
\end{proof}

\begin{lemma}[A conditional early paid-reference contract]
\label{app:tail:lem:whole58reference}
Keep the whole-diagram and threshold conventions of
Lemma~\ref{app:tail:lem:whole58local}. Suppose the actual helper-or-strong
window, with only its permitted isolated dotted-five exception, has
been completed through sixteen. For either actual strong
$K\in\{15,16\}$ there is a paid vector $W$ supported through $K$,
with $\sum W_j=F_K$, remaining individual endpoint-$K$ coefficient
$>1/9$, and effective energy $L_W>Q_K$, where
$Q_{15}=.94006,Q_{16}=.95618$. The three initial-five branches,
their later conditional extensions to an actual strong nineteen/twenty,
and the complete flat-pressure table below have the stated positive-dual
and same-vector scope. This lemma supplies no missing short-layer
classification or strong endpoint.
\end{lemma}
\begin{proof}
Use $F_i=\sqrt i-.58/\sqrt2$ throughout this proof.

\paragraph{Early references in all three five branches.}
Put $\mathcal E_K=F_5^2/5+\sum_{i=5}^{K-1}d_i^2$.
Changing the initial loss from $.61$ to $.58$ adds
$ .03\sqrt{2/5}-.00357>.0154$.
One omitted internal floor $j$ loses $<1/[32(j-1)^3]$.
All holes $j\ge6$ therefore cost less than
$\frac1{32}\int_4^\infty t^{-3}dt=1/1024$.

For either actual strong $K=15$ or sixteen we need a paid vector
$W$, supported through $K$, with $\sum W_j=F_K$, endpoint
coefficient $>1/9$ and effective energy
\[
 L_W=2D_W-\norm W^2>Q_K,\qquad Q_{15}=.94006,\quad Q_{16}=.95618.
\]
Here $D_W$ is its pressure demand. Effective energy need not equal
the square of the vector.
If five strong, the concave reference and the full hole bound give
these values; its endpoint coefficient is at least $d_{K-1}>1/8$.
If five a weak helper, interpolate through $(1,.5455),(6,F_6),(K,F_K)$.
Its first-six energy exceeds $.74385$, and its tail energies at
fifteen/sixteen/nineteen exceed $.225,.2404,.2804$.
The slopes decrease and the endpoint coefficient
$1/(\sqrt K+\sqrt6)>1/9$; the first two energies exceed the required
$Q_K$.

If five is isolated dotted, its actual core gives
\[
 U'=P_5+(u/2)(e_6-e_5),\qquad
 x.U'\ge T'={\sqrt{10-u}-\sqrt{u(1-u)}\over\sqrt2}.
\]
Its loss is $<.588$, since with $a=\sqrt{10}-3<13/80$,
$au+\sqrt{u(1-u)}<(a+\sqrt{1+a^2})/2<.588$.
Weakness forces $u>.4$: at $u\le.4$ the coarser loss
$au+.5<.565$ is too small. In particular $T'>1.8185$.
Choose actual strong $J=6$ or seven and $K=15$ or sixteen.
Let $W_0$ interpolate through $(J,F_J),(K,F_K)$, and put
\[
 V=U'-(5/J)P_J,\qquad W=W_0+\mu V,
 \qquad \mu=.21\ (J=6),\quad .2\ (J=7).
\]
The residual has sum zero, support through $J$, and $W_0.V=0$.
Its squared norms are bounded by $77/150,194/175$, since
$U'^2=5-u+u^2/2$ decreases for $.4<u<1$.
The deficits exceed $.119,.221$. The $J$ payments are $.175,1/7$,
below the available $.1808,.165$; those capacities only increase
when $K$ increases. The endpoint-$K$ coefficient stays unchanged.
The complete energy data are
\[
\begin{array}{c|c|c|c|c}
J&F_J^2/J+2\mu\Delta-\mu^2V^2>&\text{tail}_{15}>&
 \text{tail}_{16}>&\text{tail}_{19}>\\\hline
6&.720462&.225&.2404&.2804\\7&.757957&.188&.2037&.2445 .
\end{array}
\]
Thus both early columns supply $L_W>Q_K$. Every source coefficient
is nonnegative, and the construction remains legal at a later actual
strong $K=19$ or twenty. No strong-five or terminal premise has been
used in constructing the early references.

For an actual normalized pressure $x.Z>G$, $Z_1=\cdots=Z_{16}=1$
and $Z^2\le R$, pay $1/9$ at the remaining endpoint $K=15$ or
sixteen. The residual $Z-P_K$ has zero product with $W$ because
$\sum W_j=F_K$, regardless of whether $L_W=\norm W^2$. One square gives
\[
 e>L_W+2(G-F_K)/9-(R-K)/81.
\]
This uses a positive-dual payment, not the invalid subtraction of two
lower pressure bounds. With $F_{15}<3.463,F_{16}<3.590$, each row
below is sufficient in both early references:
\[
\begin{array}{c|c|c}
G&R&\text{minimum excess above one}>\\\hline
3.956&19&.00023\\4&19.7&.0013\\4.05&20&.0087\\
4.1&20&.0198\\4.15&21&.0186\\4.4&22&.0618 .
\end{array}
\]
The remaining task is to verify actual point pressures fitting this
table; a diagram count or an abstract Gram completion would not suffice.
\end{proof}

\begin{proposition}
\label{app:tail:prop:whole58}
None of the entire finite diagrams $B_{22},D_{22},D_{23}$ occurs at
an actual farthest vertex.
\end{proposition}
\begin{proof}
The local windows of Lemma~\ref{app:tail:lem:whole58local} reach
eighteen/seventeen/eighteen respectively, and hence already contain
sixteen. Actual helper nonadjacency supplies a strong fifteen/sixteen,
so Lemma~\ref{app:tail:lem:whole58reference} applies with all its
individual price conditions. We now complete the genuinely different
short layers and fund their actual pressures by that table.

\paragraph{The actual $B_{22}$ short tails.}
At nineteen deletion leaves $A_{18}+B_3$. The new short-three $F_4$
role has residual $1/4$ and numerator $g+1/(2\sqrt2)$.
Own zero has its actual affine $F_4$ kernel, and every positive own
has ratio at least $1+1/\sqrt2>83/50$. The long-three and all full
bridges are strong. At twenty the right $B_2$ allows three at either
end. The long end is the same local row. The short end has residual
$1/2$ and numerator $g+1/(2\sqrt2)$; own zero/four/six/parallel
have their actual affine/null principals, while owns at least seven
are directly strong. In particular the own-zero short kernel on
$c,n_{22},n_{21},n_{20},n_{19}$ is
$(\sqrt2,2\sqrt2,3,2,1)$. Own four/six use respectively
$\sqrt2c+n_{19}+2n_{20}+n_{21}$ and
$\sqrt3c+2n_{20}+n_{19}$. All four left/right endpoint bridges
have residual $1/40$ and minimum ratio
$\sqrt{40}(1/40+1/(2\sqrt2))>83/50$, including own zero.

Own five uses the actual left $H_4$ point, with
$\alpha=\sqrt2g+1/2$, $A^2=20\alpha^2-1/2$, $R<20$.
The positive squares $A/\alpha>4.45$ and
$\sqrt{\delta_5}/\alpha<.083$ give $G>4.3$.
Own three is the actual cap $v=(\omega_{20}+\omega_{22})/2$.
Releasing old twenty-one gives a positive finite $A_{22}$ endpoint,
with its omitted-wall product strict and every unknown wall retained.
Its cosine is $(21+20\sqrt2)/\sqrt{23}>245/24$.
It gives $x.P_{21}>4.15$, $R=21$: the last comparison is
$1511(245/24)>15289$. Both points fit the paid table. Thus through
twenty only helpers, strong indices and possibly isolated dotted five
remain in $B_{22}$.

\paragraph{The actual $D_4$ and $D_3$ tails.}
For $D_{22}$ these are eighteen and nineteen; for $D_{23}$ they
are nineteen and twenty. The $D_4$ finite list allows three at any
leaf. The fork-three residual is $1/2$, numerator $g+1/4$.
Own zero/three/four/six/parallel have their exact principals.
Own at least eight is strong because
$\sqrt2(\cos(\pi/8)+1/4)>33/20>T_0(q)$.
Retain own five and the newly possible seven. Own five uses its
actual left $H_4$ point, with $A^2=i\alpha^2-1/2$,
$\alpha^2=\phi+5/8$ and residual ratio $<.01$.
At $i=18$ its first sixteen coordinates are not all flat.
Use instead the independently valid old $.61$ strong-five reference
and an old strong eleven/twelve: the first twelve are flat and the
old terminal price exceeds $1/7$. The pressure $G>4.1,R<20$ gives
energy $>1.0102$ in both rows. The old strong-five bootstrap here
uses only the complete $M=16$ window through seventeen.
Own seven uses
\[
 \xi=c+g^{-1}n_i+(2g)^{-1}n_{i-1},\qquad k=1-3/(4g^2)>0.
\]
Its cap and $i-1$ products vanish; the $i$ product is negative,
its square is $-k$ and its time is positive. For $.9<g<.902$,
at $i=18$ we have $A/\alpha>4.22$ and
$\sqrt{k}/\alpha<.18$, so $G>4.04,R<18$. Both quantities
only strengthen in the $D_{23}$ layer $i=19$; its $R<19$ is
still covered by the first paid row. All selected displacements
are flat through sixteen in this seven case. Bridges and long leaf
have their strong or own-zero cycle dispositions. Both forks remain.

At the right $D_3=A_3$ layer, middle-three is direct or affine.
End-four retains its actual affine own-zero kernel and is otherwise
strong. End-five is sealed $H_4$ and directly strong even at own
zero. Full left bridges retain their own-zero cycle and direct
positive-own bounds. Thus only no-left end-three remains, with
residual $5/8$, numerator $g+1/4$. Own zero/three/four/six/parallel
have the actual $E_8$/cycle/star/path/null principals. Own five is
the actual left $H_4$ point, $A^2=i\alpha^2-5/8$, giving
$G>4.1,R\le20$. Every $g\ge53/50$ is directly strong, since
$(131/100)^2(8/5)>(33/20)^2$.

For the whole interval $.9<g<53/50$ use both actual $A_2$ sides,
with corrections $g(4n_i/3+2n_{i-1}/3)$ and
$g(4n_i/3+2n_{i+1}/3)$. Their positive future cores have square
$-d$, $d=4g^2/3-1>0$, and all selected products are zero or
negative, including the additional contacted fork. Average their
pressures retaining the full correction cross $8g^2/9$.
The displacement is flat through sixteen, and
\[
 x.Y/\alpha\ge{A-\sqrt d\over\alpha},\quad
 A^2=i\alpha^2-5/8,\quad
 Y^2=A^2+1-14g^2/9<i\alpha^2.
\]
At $i=19$, splitting at $g=26/25$ gives respectively
$(A/\alpha,\sqrt d/\alpha)>(4.33,0)$ with radical $<.365$,
and $A/\alpha>4.337$ with radical $<.381$. Thus $G>3.956,R<19$.
At the actual $D_{23}$ index twenty the same positive squares give
$A/\alpha>4.44$ in the first interval and $>4.45$ in the second;
the radical bounds are unchanged. Hence $G>4.069>4.05,R<20$,
a separate sufficient table row. The increasing ratio $d/\alpha^2$
justifies both full intervals, not only the endpoints. This proves
the complete window through twenty for whole $D_{23}$ without any
claim at its later centre twenty-one.

\paragraph{The remaining actual $D_{22}$ centre twenty.}
Here the retained white diagram is $A_{19}+2A_1$. Its complete
finite list is one left endpoint-three/four or next-to-end-three;
with one fork only a left endpoint-three and fork-three/four;
with both forks the all-three endpoint $D_{22}$ bridge.
Their residuals are $1/40$ or $11/40$; the directly strong
branches include zero own where appropriate, and the left/fork-three
zero/three rows retain the cycle/triangle. Internal branches cannot
be finite at this rank. With no fork, the locally rechecked bulk
inventory applies. The isolated-five correction uses only actual
long roots and the spatial sum $\omega_{21}+\omega_{22}$.

With no left contact, one finite fork label or the pairs
$(3,3),(3,4),(3,5)$ exhaust the white possibilities. At own zero,
fork three/four/six have their affine principals. Fork five uses
the actual left $H_4$ point: $R=19$, $A^2/\alpha^2=6+8\phi$,
$\sqrt{\delta_5}/\alpha<15/56$, giving $G>4.08$.
The last bound uses $\phi>79/49$, because
$5\cdot49^2-109^2=124>0$; it is not inferred just from $\phi>8/5$.
For fork at least seven put $b=\delta^2\in(.81,1)$. The actual core
\[
 c+\delta^{-1}n_f+2(\delta^{-1}-\delta)n_{20}
\]
has zero cap/fork products, negative centre product, positive time
and negative square. Its normalized data are
\[
 R(b)=30-16/b+6/b^2,\quad
 G(b)=\sqrt{22-2/b}-\sqrt{-8+14/b-6/b^2}.
\]
For $b\le.9$, $R<19.7$ and $G>4.41-.41=4$.
For $b\ge.9$, $R\le20$ and $G>4.44-.39=4.05$.
The derivatives have numerators $16b-12$ and $12-14b$;
the latter radical square has maximum $1/6$, or $4/27$ on the
second interval. These are whole-interval positive-square bounds.

At own three, fork three has its triangle. Fork four/five instead
gives the actual finite $D_{22}$ contacted-fork endpoint, with
cosine $(20+22\delta)/2>17$, signed direction
$\sqrt2\omega_f=P_{21}\mathbin\pm e_{22}$, $R=22,G>4.4$.
The released-wall product is strict, so this is an actual endpoint.
Fork at least six is direct. Own four has its centre/left/other-fork
affine star. For own at least five and fork at least four, direct
ratios leave only own-five/fork-four; it has the same left-$H_4$
data as the $B_2$ own-five case, $G>4.3,R<20$.
It must not be discarded by the old $.61$ ratio.
For fork-three/own-five, the left-$H_4$ point has
$A^2=20\alpha^2-3/4$, $G>4.1,R<20$.
For other finite or dotted owns, $g\ge6/5$ is direct. On
$.9<g<6/5$, the actual core above with $n_{20},n_{19}$ has
$k=1-3/(4g^2)$, $A^2=20\alpha^2-3/4$ and $R<20$.
The full derivative bound $k/\alpha^2\le575/5046<.34^2$
and $A/\alpha>4.44$ give $G>4.1$.

For two forks, own three has its triangle and owns at least four
are direct. At own zero, pair-three-three has its square.
Pair-three-four has an actual finite $B_{22}$ contacted-fork
endpoint, cosine $(40+22\sqrt2)/4>17$, again $G>4.4,R=22$.
This adjacent vertex is not declared farthest. Pair-three-five uses
the actual finite $H_3$ projection, with displacement $\alpha P_{20}$,
$q=2\phi^{-6}<1/8$ and pressure
\[
 \sqrt{20-q}/(1+\sqrt{q/20})>4.45/1.08>4.1,\quad R=20.
\]
Every listed point fits the paid table. This exhausts the actual
centre classification and completes the window through twenty in
$D_{22}$ as well. All signs of the final coordinate were retained.

\paragraph{The final reference, without an assumed strong five.}
Nineteen/twenty are now helper-or-strong, so one is strong.
If five a weak helper, its earlier interpolation has energy
$>.74385+.2804>1$. If five isolated dotted, its same paid
reference has effective energy $>.720462+.2804>1$ or
$>.757957+.2445>1$. The endpoint-$J$ payment remains legal
because its tail coefficient decreases. Thus five is strong,
proved only after the terminal classifications.

If twenty strong, the reference energy exceeds
$\mathcal E_{20}-1/1024>.99752+.0154-1/1024>1$.
Otherwise twenty weak gives strong nineteen,
$\mathcal E_{19}>.9847+.0154=1.00010$, and $F_5/5<.366$.
If a weak helper occurs at six through nine, its actual $x_1>.45$
adds $>(5/4)(.45-.366)^2=.00882$, more than the full hole loss.
If none occurs there, holes start at ten and cost less than
$\frac1{32}\int_8^\infty t^{-3}dt=1/4096$.
The weak-twenty point gives $x_1>.378$ and adds
$>(5/4)(.378-.366)^2=.00018$. The total is
\[
 e>1.00010+.00018-1/4096>1.0000358>1.
\]
Each initial adjustment leaves positive decreasing slopes at both
junctions. Effective-reference, tail and initial gains above belong
to the stated disjoint branches; none is added twice. All three
whole diagrams are therefore excluded on the original vector.
\end{proof}

\subsection{Two small remainders and one extra prefix gain}
\label{app:tail:smallremainders}
\begin{lemma}
\label{app:tail:lem:A20ranktwo}
An actual $A_{20}$ component cannot have spectator rank two.
\end{lemma}
\begin{proof}
Its energy is $>6/7$. Every ordinary irreducible rank-two remainder
costs $>1/7$, and two ordinary $A_1$'s cost $>.2$. A $J$ pair
costs $>.34$. The common joint inequality $I+A_1>1/5$ also
excludes one actual $I$ plus its ordinary $A_1$; it counts $I$ once.
The complete remaining possibilities are the two isolated typed
rank-two exceptions $N_7,O_7$.

Either has energy $>.09$, so $e_A<.91$. The actual typed
propagation proof in Appendix~\ref{app:typedpropagation} forces
every other edge half-tangent square $>1/5$ for $N_7$.
For $O_7$, no opposite of an $A_{20}$ edge can contact the
retained old $I_2(7)$: its least white cost is
$\tfrac14\csc^2(\pi/7)>1$. The actual native gap then gives
$A^2>3/4$, as in~\eqref{app:eq:O7height}. In either case every
normalized $A_{20}$ edge height square exceeds $2/15$, since
$4(2/15)/(1+2/15)^2=120/289<3/7$.

The $N=21$ central inventory, foreign clearing and literal-helper
nonadjacency above do not use the contrary target $6/7$ until
the near-four comparison. Recheck that single step at $e_A<.91$.
Its actual mate has height $>.86$ and an opposite mate or actual
strong anchor has height $>.77$. Their correlation is at most
$1/\sqrt7<2/5$. Cauchy would give
\[
 1.63^2<.91(14/5),
\]
whose left-minus-right is $.1089>0$. Thus near-four is excluded
under this target too. There is a strong $i\in\{5,6\}$ and
strong $j\in\{15,16\}$, independently of the end-height payment.

In zero-sum coordinates, the two endpoint heights give
\[
 x_1>\sqrt{8/63}>a=89/250,\qquad -x_{21}>a.
\]
The positive squared margin is $8(62500)-63(7921)=977$.
The two possible actual strong floors give
\[
 S_5>151/100,\qquad S_6>131/80.
\]
For the first use $D_5=160/21>(11/4)^2$ and $\sqrt2<99/70$.
For the second use $D_6=60/7>(2927/1000)^2$ and the same
upper bound for $\sqrt2$. On the first $i$ coordinates the paid
reference is $a,(s_i-a)/(i-1),\ldots,(s_i-a)/(i-1)$.
Its slopes are positive and decreasing, so it uses nonnegative
multipliers of the actual pressures $x_1>a,S_i>s_i$.
The exact two costs are
\[
 a^2+(s_5-a)^2/4=.459665,\qquad
 a^2+(s_6-a)^2/5=.45518445.
\]
Reflection supplies the same bounds on the last five/six coordinates.
Those coordinate intervals are disjoint, so their costs add on the
same original $A$ projection: $e_A>.9103689>.91$.
This excludes both typed exceptions and exhausts the rank-two remainder.
\end{proof}

\begin{lemma}[The final packet preserves an extra initial gain]
\label{app:tail:lem:prefixgain}
With exactly the hypotheses of Proposition~\ref{prop:prefix}, let
$G_0=3969/800000=.00496125$. Then its same actual prefix energy is
greater than $E_M-H_M/3200+G_0$, for every $6\le M\le16$.
This is a replacement bound, not an independent gain to add to a
different initial-coordinate coupling.
\end{lemma}
\begin{proof}
The weak-packet points and full scalar seals in the cap proof give
$x_1>.424$ for a weak packet at every $5\le i\le21$.
Here we use only $i=M+1\le17$. Its selected actual roots are
$i,1,2$ for native five, $i,1$ for seven/eight, $i$ for the
remaining finite owns, or the actual dotted pair. Its converted
coordinates use old long neighbours $i-1,i+1$, which exist because
$i+1\le r-3$ in $B$ and $\le r-4$ in $D$.
All foreign spectators have already been cleared and all point
coefficients and wall signs are those of that pressure proof.

If five weak, the generic-prefix proof gives
$L_M-E_M>.76-.7-9/640=147/3200>G_0$; at $M=6$ use its
separate direct $L_6>.76,E_6<.7$ comparison. No strong-five
conclusion is needed in that branch.
Now suppose five strong and put $K=M+1$.
If $K$ strong, its reference gains an extra increment
$d_M^2>1/(4M+2)\ge1/66$. At most one extra hole costs
$1/3200$, and $G_0+1/3200<3/500<1/66$.
This also accounts for the possible hole at six when $M=6$.
If $K$ weak, $M$ strong and the actual packet gives $x_1>.424$.
Modify only the initial five reference slopes to
\[
 .424,\quad(f_5-.424)/4,\quad\ldots,\quad(f_5-.424)/4.
\]
They are positive and decreasing at both junctions: $1.8<f_5<1.805$,
and the next slope is $<1/4$ even after omitting its neighbour.
The exact gain is $(5/4)(.424-f_5/5)^2>G_0$.
The old hole intervals lie after five and have disjoint interiors,
so their losses do not overlap this change. These are legal reference
duals on one vector, with no tail or total-norm bootstrap.
\end{proof}

\subsection{The complete forest accounting in dimensions twenty-three and twenty-two}
\label{app:tail:23and22}
\begin{proposition}
\label{app:tail:prop:23and22}
Dimensions twenty-three and twenty-two are impossible.
\end{proposition}
\begin{proof}
At either dimension there is exactly one dominant component, by the
same ordinary/debit and half-cost argument used above. The three
spectator costs are still complete. For $S\ge2$ the typed cost is
larger than the short-rank-discard cost by $13/700$.
For dominant $A$, the half-cost gives $S\le7$, so $r\ge16$ at
dimension twenty-three and $r\ge15$ at twenty-two. The caps give
$r\le20$. The following weaker consequences of already written
component theorems suffice:
\[
\begin{array}{c|cccccc}
r&15&16&17&18&19&20\\ \hline
e_A>&4/7&.86&6/7&.91&6/7&6/7.
\end{array}
\]
At dimension twenty-three, the weakest short-discard totals for
$r=16,17,18,19,20$ are respectively
$ .86+6/14$, $6/7+5/14$, $.91+4/14$, $6/7+3/14$ and
$6/7+2/14$, all at least one and strictly exceeded. Every other
debit case is stronger. At twenty-two, ranks fifteen through
nineteen similarly use spectator ranks seven through three:
$4/7+6/14=1$, $.86+5/14>1$, $6/7+4/14>1$,
$.91+3/14>1$ and $6/7+2/14=1$, strictly exceeded.
The last $A_{20}$ rank-two remainder is exactly
Lemma~\ref{app:tail:lem:A20ranktwo}. Thus no $A$ remains.

For dominant $B/D$, the $9/14$ floor gives $S\le5$.
At twenty-three this means $r\ge18$. The generic $M=r-6$
floors for ranks eighteen/nineteen/twenty, even for the weaker $D$
requirement, are $>.8689525,.88864,.90715$; their spectator ranks
five/four/three give totals above one after either debit.
At rank twenty-one the floor $>.9240775$ meets a rank-two remainder.
A typed block costs $>.09$. An ordinary irreducible rank-two block
costs $>1/7$. Two $A_1$'s have at least one ordinary cost $>.1$,
including the possible short-rank discard; $J$ is stronger. Thus all
these totals exceed one. Rank twenty-two with its floor $>.9401975$
permits only isolated $I$, because an ordinary $A_1$ costs $>.1$.
It is excluded by Proposition~\ref{app:tail:prop:classicalI}.
At rank twenty-three, $B_{23}$ has its cap and whole $D_{23}$ is
excluded by Proposition~\ref{app:tail:prop:whole58}.

At dimension twenty-two, $r\ge17$. Ranks seventeen/eighteen/nineteen
have floors $>.8472225,.8689525,.88864$ and spectator ranks
five/four/three. Their weakest costs $4/14,3/14,2/14$ suffice.
At rank twenty the improved standalone prefix floor is
\[
 .90715+G_0=.91211125.
\]
Every rank-two remainder costs $>.09$: a typed exception does so
directly; an ordinary irreducible costs $>1/7$; two $A_1$'s include
one nonshort cost $>.1$, and $J$ is stronger. All totals exceed one.
Rank twenty-one has only one spectator. An ordinary $A_1$ contradicts
even the original $>.9240775$ floor, while actual $I$ is excluded
by Proposition~\ref{app:tail:prop:classicalI}. At rank twenty-two
only whole $B_{22},D_{22}$ remain, excluded by
Proposition~\ref{app:tail:prop:whole58}. Every own-rank finite forest
has now been covered. The argument never descends to an arbitrary
face or uses a higher-dimensional cover at a different rank.
\end{proof}

\subsection{One rational curve for the two remaining large-$A$ couplings}
\label{app:tail:Atwoends}
We use the same central inventory with a different reference, rather
than starting a new component classification. At $N=20$ the target
will be $e_A<.91$; at $N=21$ no target beyond $e_A<1$ is needed.
The complete pure-helper inventory and its nonadjacency were proved
before the $6/7$ contradiction. Only the removal of near-four must
be revisited. At $N=20$ its mate and opposite anchor have height
sum $>1.63$, correlation $<.4$ and hence sum-direction square
$<2.8$. This contradicts $e_A<.91$, with margin $.1089$.
At $N=21$ the opposite anchor is fifteen or sixteen. At fifteen its
height exceeds $.79$, since $D_{15}=60/7>(61/21)^2$, and its
correlation with the near mate is $\le1/\sqrt8<.354$.
Then $1.65^2>2.708$. At sixteen its height exceeds $.77$ and
the correlation is $\le5/16$, so $1.63^2>2.625$.
An opposite near-four mate has at least the latter height and smaller
correlation. All these comparisons contradict $e_A<1$.
Thus the central weak indices are nonadjacent pure helpers on both
stated domains, with arbitrary finite spectators retained.

Use the symmetric rational heights $H_{N-i}=H_i$, linear from
zero to five and from $N-5$ to $N$:
\[
\begin{array}{c|rrrrrr}
N&H_5&H_6&H_7&H_8&H_9&H_{10}\\\hline
20&2.128&2.288&2.406&2.488&2.536&2.552\\
21&2.150&2.317&2.445&2.537&2.597&2.626 .
\end{array}
\]
At $N=21$, $H_{11}=H_{10}$. Every entry is strictly below
$\sqrt{D_i}-.61$, by positive squaring, and successive differences
decrease. The tight $N=21,i=7$ comparison, for example, has
margin $28/3-3.055^2=925/3000000>0$.
The coordinate reference is $w_j=(H_j-H_{j-1})/\sqrt2$.
Summation by parts uses nonnegative actual strong-floor coefficients
and gives $x_A.w\ge\norm w^2$ at its available bends.

If end five is weak, replace that entire end by the straight line
zero to six; six is strong. Do the reflected replacement at a weak
$N-5$. Delete each remaining weak $j$ by averaging its two neighbours.
The exact energy loss is
$(H_{j-1}-2H_j+H_{j+1})^2/4$. Nonadjacency gives disjoint affected
edge pairs, so these losses add for every permitted weak set.
The following finite table records the full bound, not selected weak
patterns. In the loss column, count is the maximum number of
nonadjacent internal holes and the second factor is the maximum
quarter-square on the remaining interval:
\[
\begin{array}{c|c|c|c|c}
N&\text{weak ends}&\text{initial energy}&\text{internal loss}\le&
 \text{surviving energy}>\\\hline
20&0&.9544848&5(.000441)&.9522798\\
20&1&.9250917333\ldots&4(.000441)&.9233277\\
20&2&.8956986666\ldots&4(.000324)&.8944026\\
21&0&.981678&5(.00038025)&.97977675\\
21&1&.9528575833\ldots&5(.00038025)&.9509563\\
21&2&.9240371666\ldots&4(.000324)&.9227411 .
\end{array}
\]
For an exact rational formula behind every repeating decimal, put
$a=H_5,b=H_6,T=\sum_{j=7}^{\lfloor N/2\rfloor}(H_j-H_{j-1})^2$.
The three initial energies are
\[
 a^2/5+(b-a)^2+T,\quad
 \tfrac12[b^2/6+a^2/5+(b-a)^2+2T],\quad b^2/6+T.
\]
In the two-weak-end rows, six and $N-6$ are strong and the holes
lie in $7..N-7$, giving four possible holes in either rank.
The zero/one-end interval lengths give the other counts. Direct
second differences of the displayed rational heights give the
per-hole maxima. This proves completeness of the finite loss table
without a combinatorial enumeration or an arithmetic packing premise.

An actual endpoint floor $x_1\ge z$ improves a constant initial
block of length $m$ and coordinate mean $b_0$ by replacing its
entries with $z,(mb_0-z)/(m-1),\ldots$. If $z>b_0$, the gain is
\[
 {m\over m-1}(z-b_0)^2.
\]
If $z\le b_0$, leave the block unchanged. All reference sums are
preserved. For every $z\le1$ the remaining means exceed the next
coordinates: at $N=20$, the first-five and first-six comparisons
are $>.126$ versus $<.114$, and $>.123$ versus $<.084$;
at $N=21$, they are $>.130$ versus $<.119$, and $>.127$ versus
$<.091$. Deleting an internal neighbour only decreases that next
coordinate. Thus the new duals are nonnegative at both junctions.
The first and last blocks have disjoint supports, so their changes
add without changing any intervening height.

\begin{lemma}
\label{app:tail:lem:A19typed}
An actual $A_{19}$ cannot coexist with either typed rank-two exception.
\end{lemma}
\begin{proof}
The typed energy $>.09$ gives $e_A<.91$, the stated $N=20$ domain.
Zero or one weak ends already contradict that target by the table.
With two weak ends its reference energy exceeds $.8944026$.
Both typed exceptions give every other actual normalized direction
height square greater than $(11-4\sqrt7)/3>11/80$.
For $N_7$ the stronger $>1/5$ propagation suffices; for $O_7$ use
the actual no-contact/native-height proof in
Appendix~\ref{app:typedpropagation}. The rational seal is
$\sqrt7<847/320$, with squared margin $609/102400$.
At the two actual $A_{19}$ ends this gives
\[
 x_1^2,x_{20}^2>(19/20)(11/80)=209/1600>.36^2.
\]
The first and last six-block means are $<.27$. Their legally
funded disjoint modifications add $>2(6/5)(.36-.27)^2=.01944$.
The same-vector energy is therefore
$>.8944026+.01944+.09=1.0038426>1$.
Both typed roles, rather than a bare dihedral label seven, are covered.
\end{proof}

\begin{lemma}
\label{app:tail:lem:A20A1}
An actual $A_{20}$ has energy $>.92274$ and cannot coexist with
an ordinary nonshort $A_1$ or an actual isolated $I$.
\end{lemma}
\begin{proof}
The $N=21$ window and table apply under $e_A<1$ with arbitrary
finite spectators. Their least energy is $>.9227411>.92274$.
An ordinary nonshort $A_1$ costs $>.1$, already contradictory.

For actual $I$, write its universal parameters as $s,t$:
$e_I\ge s^2$, $t+(3/2)s>.767$, $0<s<1/3$ and $t<1$.
Both actual normalized end heights are at least $t$. Put
$a=\sqrt{20/21}>.975$; then $x_1,-x_{21}\ge at$.
Every endpoint change has $z=at<1$, so the junction checks just
proved apply. In the zero/one-weak-end rows the two absolute block
means are $<.305$; in the two-weak-end row they are $<.274$.
Each block length is five or six, so their two changes add at least
$(12/5)(at-b)_+^2$, using $b=.305$ or $.274$ respectively.
Weighted Cauchy, on these two changes and the \emph{single} $I$ cost,
gives
\[
 s^2+(12/5)(at-b)_+^2>
 {(a\cdot.767-b)^2\over(9/4)a^2+5/12}.
\]
Indeed $(at-b)_++(3/2)as>a\cdot.767-b$.
The denominator is $215/84<2.56$. With $b=.305$ the numerator
exceeds $.442825^2$, so the gain is $>.076$.
With $b=.274$ it exceeds $.473825^2$, so the gain is $>.0875$.
The least final total is $>.9227411+.0875>1.01024$;
the zero/one-end totals are larger. No separate $I$ charge was
added at the second end, and no other component energy was reused.
\end{proof}

\subsection{The proper $D_{19}/D_{20}$ continuation}
\label{app:tail:properD1920}
\begin{lemma}
\label{app:tail:lem:D4proper}
For an actual $D_{19}$ or $D_{20}$ component, the $.61$ pure-helper
window extends through fifteen or sixteen respectively, with arbitrary
finite spectators. This continuation does not assume strong five.
\end{lemma}
\begin{proof}
The valid $M=13/14$ bulk reaches fourteen/fifteen. Its complete
geometry supplies an actual strong eleven/twelve. In either initial-five
branch there is a positive paid reference ending there, with
\[
 q_{11}>.8472225,\quad q_{12}>.8689525,\quad
 F_{11}\le f_{11}<2.886,\quad F_{12}\le f_{12}<3.034,
\]
and remaining individual terminal coefficient $>1/7$.
For weak five use the actual $x_1>.5965,S_6>2.01815$ interpolation;
it has larger energy and terminal slope $1/(\sqrt K+\sqrt6)>1/7$.
No ordinary spectator energy has been used in these references.

At $i=r-4=15,16$, deletion retains $A_{i-1}+D_4$.
Its finite extensions are three at any right leaf, or full endpoint
$D_r$ bridges. A foreign branch at the internal new node of such a
rank-$r>8$ bridge is not finite. Own zero keeps its actual cycle;
positive owns are directly strong. Long-leaf three has its bulk
strong/affine disposition. Fork three has residual $1/2-Q$ and
numerator $g+1/4$; own zero/three/four/six/parallel have the actual
principals, all $Q>0$ cases and native owns at least seven are
directly strong. Both physical forks are retained. The only extra
role is pure own-five/fork-three.

Put $g=\phi/2$, $\alpha^2=\phi+5/8$ and
$\delta_5=(5-3\phi)/8$. The actual left $H_4$ correction is
$g(3n_i/2+n_{i-1}+n_{i-2}/2)$. Its finite projected point has
positive time, all original wall signs, and normalized displacement
flat through $i-3$, with
\[
 A^2=i\alpha^2-1/2,\quad Y^2<i\alpha^2,\quad
 A/\alpha>\sqrt{i-25/112},\quad
 \sqrt{\delta_5}/\alpha<.091.
\]
Here $\alpha^2>56/25$. At fifteen this gives $G>3.75,R=15$;
at sixteen it gives $G>3.88,R=16$. Both are flat through twelve.
Pay $1/7$ at the early reference endpoint before adding the actual
point. The two lower bounds
$q_K+2(G-f_K)/7-(R-K)/49$ exceed $1.0122$ or $1.029$,
respectively, for \emph{both} $K=11,12$. Thus the extra role is
excluded without a strong-five or later-tail assumption.
The actual two-helper Schur calculation then extends nonadjacency
to the newly supplied index; all right and spectator blocks stay positive.
\end{proof}

\begin{lemma}
\label{app:tail:lem:proper21}
The entire forests $D_{19}+N_7$, $D_{19}+O_7$ and $D_{20}+I$
are impossible in dimension twenty-one.
\end{lemma}
\begin{proof}
For either typed defect beside $D_{19}$, weak five has actual
prefix energy
$L_{13}>.7600299945+1.156^2/7>.95$, using
$\sqrt{13}>3.6055,\sqrt6<2.4495$, already too much beside $>.09$.
Hence five strong. Lemma~\ref{app:tail:lem:D4proper} supplies
coverage through fifteen. If fifteen strong, its reference has
energy $>.9240775$. If fifteen weak, fourteen strong and the
actual pure packet gives $x_1>.424$, so the legally modified
initial reference gives $>.90715+.00496125=.91211125$.
In both branches the separate typed energy $>.09$ contradicts one.
This repeats the proved prefix-gain calculation only \emph{after}
supplying its missing D4 geometry; it does not invoke $M=14$
at $D_{19}$ outside that theorem's rank hypothesis.

For $D_{20}+I$, the same lemma first supplies coverage through
sixteen in both initial-five branches. The positive actual long-root
helper points, nonadjacency and initial $I$ parameters are therefore
exactly the geometric hypotheses used in the proof of
Proposition~\ref{app:tail:prop:classicalI}. Its weak-five reference
at a strong fifteen/sixteen has energy $>.9851486$ and remaining
$P_6$ price $>.12533$; the capped single-$I$ graft gives $>.015$.
Thus five strong. Its strong fifteen/sixteen flat-five reference
has $q>.9240775$, actual $P_5$ demand $f_5$ and remaining price
$>.146$. Deleting long one retains actual $D_{19}$, with just the
same isolated/long-two-three head list and only the old $I$ outside.
The selected actual head roots and both genuine forks retain the
local point signs, so the q-free head/$I$ gain $>.07665$ closes it.
We have replaced the missing $M=15$ bulk justification by a proved
actual D4 continuation, without applying its rank-restricted theorem
to $D_{20}$. No D3 or terminal-seventeen floor is needed.
\end{proof}

\subsection{One full weak-nineteen payment}
\label{app:tail:weak19}
\begin{lemma}
\label{app:tail:lem:weak19packet}
Suppose an actual pure helper at long node nineteen has own five,
finite at least seven or dotted, $g<6/5$, with actual long dual
$\sqrt2P_{19}$ and actual selected roots $n_{19},n_1,n_2$.
Let a paid classical vector $W$ be supported through eighteen, have
sum $F$ and retained actual $P_{18}$ demand $F$, constant first
three entries $a$, no foreign coordinate,
effective energy $L>.985613$, and remaining $P_{18}$ coefficient
$>.107$. If $\sqrt2F\le38/7$ and $a/\sqrt2<.267$, its full actual
foot contradicts the norm budget. This is a conditional local
pressure lemma, not a classification of node nineteen.
\end{lemma}
\begin{proof}
The helper time is $A^2=38g^2+2g-1/2$. Its actual finite or
future cores, using the initial and own long roots, give
\[
\begin{array}{c|c|c}
\text{own}&U&T\\\hline
5&g(\sqrt2P_{19}-n_{19})+H_3&H(A,2/3-g^2)\\
7,8&g(\sqrt2P_{19}-n_{19})+H_2&A-\sqrt{g^2-3/4}\\
\ge9\text{ finite}&g(\sqrt2P_{19}-n_{19})+H_1&H(A,1-g^2)\\
\text{dotted}&g\sqrt2P_{19}-g^{-1}n_{19}+H_1&A-\sqrt{1-g^{-2}},
\end{array}
\]
where $H_k=P_k/(k\sqrt2)$ and $x.U\ge T$.
These are the actual correction/finite-projection recipes already
proved, now at their physical index nineteen. They require no $B$
short normal. In type $D$, $\sqrt2P_{20}$ occurring in the
coordinate sum is the sum of the two actual fork duals, and both
fork products remain nonpositive under each correction.

Put $\alpha=\sqrt2g$ and $V=U-\alpha P_{18}$.
The four variances are $g^2+1/6,g^2+1/4,g^2+1/2$ and
$2g^2-3/2+g^{-2}$; all are $<2.1$, the last increasing to
$1867/900$ at $6/5$. The helper mass in the first three
coordinates is exactly $1/\sqrt2$, so the \emph{full} cross is
\[
 W.U=\alpha F+a/\sqrt2,
 \qquad \Delta=T-W.U>T-g(38/7)-.267.
\]
The following complete scalar seals prove $\Delta>.19$.
For own five use $A>5.097,\sqrt\delta<.111,g<.81$;
for seven use $A>5.663$, radical $<.253,g<.902$;
for eight use $A>5.8119$, radical $<.323,g<.9239$.
Each substituted difference exceeds $.19$.
For finite $g\ge.939$, $H-g(38/7)$ increases because
$A'>\sqrt{38}$, $H_A>47/50$ and $H_\delta<0$.
At $.939$, $A>5.906$, the other time radical $>5.916$ and
$\sqrt\delta<.344$ give $H>5.58$, again sufficient.
For dotted owns the deficit decreases on $(1,6/5)$:
$A'-\sqrt{38}<.1$, $\sqrt{38}-38/7<.74$ and the subtracted
radical derivative exceeds one. At its enlarged endpoint,
\[
 \Delta>7.5246-.55278-(6/5)(38/7)-.267>.1905>.19.
\]
These are positive-square and whole-interval derivative bounds,
not sampled native labels or a numerical maximum.

Pay $\mu=1/16$ at the true endpoint-eighteen coefficient.
Its price $\mu\alpha<.107$ is available. Adding the actual $U$
pressure and completing the \emph{same} square gives
\[
 e>L+2\mu\Delta-\mu^2V^2
 >L+.015546875>1.0011598>1.
\]
The nonzero $W.V=a/\sqrt2$ is already included in $\Delta$;
no separate early-block or head energy has been added.
\end{proof}

\subsection{The whole $B_{21}$ terminal}
\label{app:tail:wholeB21}
\begin{proposition}
\label{app:tail:prop:wholeB21}
The entire farthest forest $B_{21}$ is impossible.
\end{proposition}
\begin{proof}
Lemma~\ref{app:tail:lem:whole58local} supplies the complete whole
$.58$ bulk through seventeen, with isolated dotted five retained.
There is an actual strong fifteen/sixteen, so the full paid-reference
contract of Lemma~\ref{app:tail:lem:whole58reference} applies.
At eighteen the right $B_3$ permits long-three or sealed short-three
$F_4$. Own zero has its actual affine $F_4$ principal and every
positive own has ratio at least $1+1/\sqrt2$. Full endpoint bridges
are directly strong, including own zero. Thus only pure helpers
can be weak there.

At nineteen the right $B_2$ permits three at either endpoint.
The long end is the same local row. At the short end, own
zero/four/six/parallel have the actual $F_4$/star/path/null
principals and owns at least seven are strong. Every full bridge
has residual $1/38$, with minimum ratio
$\sqrt{38}(1/38+1/(2\sqrt2))>83/50$, including own zero.
Own five has its actual left $H_4$ foot on nineteen/eighteen/seventeen,
$\alpha=\sqrt2g+1/2$, $A^2=19\alpha^2-1/2$ and $R<19$.
The squares $A/\alpha>4.337$ and
$\sqrt{\delta_5}/\alpha<.083$ give $G>4.2$; it is flat
through sixteen and is excluded by the early paid table.

Own three instead gives the same actual cap
$v=(\omega_{19}+\omega_{21})/2$ at the finite $A_{21}$
replacement of old twenty. Every original wall is retained and
the released wall is strict. Its exact cosine
$(20+19\sqrt2)/\sqrt{22}>466/47$ gives
$x.P_{20}>4.04,R=20$, since $2299(466/47)>22701$.
The early fifteen/sixteen paid formula gives respectively
$>1.00655,>1.00679$. This is an actual edge from the original
vertex, not a farthest assumption at its mate. The complete
helper-or-strong window therefore reaches nineteen.

One of eighteen/nineteen is strong. Weak helper five would give
the earlier interpolation energy $>.74385+.2679>1$.
For isolated dotted five, the same zero-sum residual reference at
a strong $K=18,19$ has these effective lower bounds:
\[
\begin{array}{c|c|c|c}
J&\text{initial base plus gain}&\text{tail}_{18}>&\text{tail}_{19}>\\\hline
6&.720462&.2679&.2804\\7&.757957&.2318&.2445 .
\end{array}
\]
Its endpoint-$J$ payment was legal at fifteen and becomes easier.
At strong eighteen it has $L>.988362$ or $>.989757$;
at nineteen it already exceeds one. Its sum is $F_K$, endpoint
coefficient $> .14$, and its first three coordinates are constant
$a=F_J/J+\mu_0(1-5/J)<.377$.

If nineteen weak, eighteen strong. For strong five its ordinary
concave reference has $L>.97119+.0154-1/1024>.985613$,
sum $F_{18}$, first-five mean $<.366$ and endpoint coefficient
$>1/9$. The isolated-dotted reference just described is stronger
and also has first-three mean $<.377$. Both satisfy
$\sqrt2F_{18}=5.42<38/7$ and $a/\sqrt2<.267$.
Every remaining weak-nineteen cap is an actual pure helper with
the required physical roots. Lemma~\ref{app:tail:lem:weak19packet}
therefore excludes it, keeping the full early residual cross.

It remains strong nineteen; weak helper five and isolated dotted
five were already excluded in that branch. For strong five,
$\mathcal E_{19}(.58)>1.00010$. If an internal helper occurs
at six through nine, its actual $x_1>.45$ gain $.00882$ pays
the full hole loss $<1/1024$. Otherwise every hole is at least
ten; a helper's $x_1>.378$ gain $>.00018$ pays $<1/4096$.
If there are no holes, no correction is needed. In every branch
the legal positive reference has energy greater than one.
\end{proof}

\subsection{Whole $D_{21}$: sum the exceptional feet before bounding their errors}
\label{app:tail:wholeD21}
\begin{proposition}
\label{app:tail:prop:wholeD21}
The entire farthest forest $D_{21}$ is impossible.
\end{proposition}
\begin{proof}
Use actual long roots through twenty and the other actual fork
$n_{21}=(e_{20}+e_{21})/\sqrt2$. The genuine long duals stop at
nineteen; the fork duals are $(P_{20}\mathbin\mp e_{21})/\sqrt2$.
The last coordinate is signed. This proof changes the loss to
$c_1=4/7$ and writes $G_i=\sqrt i-c_1/\sqrt2$.
We reclassify the weak directions at this threshold; no $.58$
classification is transported unchanged.

The sufficient ratio is
\[
 T_1(q)=7/4-{q-4/7\over q(2q-4/7)}.
\]
It increases on $q\ge\sqrt{10}$; positive-denominator comparisons
give $T_1(q)<1.673$ through $q=\sqrt{38}<37/6$,
$T_1(q)<1.63$ through $q=\sqrt{14}<15/4$,
$T_1(4)<1.64$ and $T_1(\sqrt{12})<1.62$.
The actual bulk $5..16$ has retained $A_{i-1}+D_{21-i}$, with
right rank at least five. Its complete finite-white list, zero-own
principals and full rank-twenty-one bridges are already proved
local geometric statements. There is no foreign component.

Near-three has the same increasing least ratio
$[i(\phi+1)-1]/\sqrt{2i(i+1)}$. At five/six native owns at least
seven are strong; at seven the own-five ratio is $>1.63$, at
eight it is $>1.66>T_1(4)$, and from nine it is $>1.68$.
Thus only near-three/own-five at five and six need points.
Next-to-end white-three has least ratio $(i\phi/2+1)/\sqrt{2i}$,
strong from six onward. At five retain its actual floor, rather
than call it strong:
\[
 h_5\ge\sqrt{10}-c_D,\quad c_D<.57549.
\]
Indeed its cosine is at least $1+5\phi/2$, giving
$h_5^2\ge10(5\phi)/(5\phi+4)>6.6915>2.58679^2$;
use $\phi>1.61803$ and $\sqrt{10}<3.16228$.
Exceptional positions have ratio $2g+3/8>1.673$.
The far-four/near-four, right long-three, finite $D$ fork and
bridge ratios or actual principals are unchanged geometric rows;
their lower bounds now exceed the displayed $T_1$ targets. In
particular a fork at retained rank at least five has ratio
$>(.8+.25)/\sqrt{3/8}>1.673$ after its own-zero/three/four/six/null
cases. Full endpoint bridges retain their positive-ratio own-zero
branch or the actual affine cycle of the other orientation.

Isolated finite five and seven have the preceding losses
$<25/76+5/24$ and $<1000/4779+9/25$, both below $4/7$.
For finite owns at least eight, use the exact singleton projection.
With $r=\sqrt{q^2-1}$ and $s=\tan(\pi/m)\le\sqrt2-1$, its
positive finite correction gives
\[
 h_i>r\sqrt{(r-s)/(r+s)}.
\]
To derive this, select the actual cap and $n_i$. Their residual is
$1-g^2>0$, the displacement after correction is $g(\omega_i-n_i)$,
and the finite-point formula divided by $g$ is
$\sqrt{r^2-s^2}/(1+s/r)$, the displayed expression.
Its loss decreases in $q$ and increases in $s$. At $q=\sqrt{10}$,
$r=3,s<.415$, it exceeds $2.6$ by
$9(517/683)>169/25$, giving loss $<3.163-2.6=.563<4/7$.
Owns three/four/six/parallel retain their actual affine/null exclusions.

For isolated dotted own, $u=g^{-2}\in(0,1)$, the positive future
core gives the actual pressure
\[
 x.(\omega_i-u n_i)\ge T_i(u)
 =\sqrt{2i-u}-\sqrt{u(1-u)}.
\]
At $i\ge8$ its loss is at most
$a_i u+\sqrt{u(1-u)}$, $a_i=1/(\sqrt{2i}+\sqrt{2i-1})<.128$,
whose maximum is $<.5681<4/7$. Retain isolated dotted five/six/seven.
Every other weak bulk row is the pure helper, with
$\phi/2\le g<6/5$; the upper cutoff is reproved from
$4(6/5)-2=14/5>(1.673)^2$.

\paragraph{The two near-three points.}
The own-five near role at five was excluded outright by the actual
right-$H_4$/old-floor argument in
Lemma~\ref{app:tail:lem:whole58local}. At six use the actual
right correction $g(3n_6/2+n_7+n_8/2)$, $g=\phi/2$.
Its finite point retains old five/nine and every unknown wall.
For its displacement $Y$,
\[
 A^2=9/2+8\phi,\quad Y^2=(41+61\phi)/8,\quad
 \sum Y_j=\sqrt2(5/2+3\phi),
\]
giving $x.Y>4.04,Y^2<17.463,\sum Y_j<10.401$.
Here $A>4.1766$, $\sqrt{\delta_5}<.136$,
$1.618<\phi<1.61804$ and $\sqrt2<1.414214$ suffice.
The valid old $.61$ bulk supplies a strong twelve/thirteen, with
floors $>3.03,>3.17$. For $Y+(3/4)P_K$, the squared demands are
$39.84765625,41.18430625$, and the norm squares are
$<39.8145,<40.377$. Both contradict Cauchy before a new tail or
strong-five conclusion is used.

New weak helpers remain nonadjacent: their actual mutual product
is $>1-ab>-11/25$, hence zero, and the positive mixed-point loss
is $<7/16+1/(4\sqrt{10})<4/7$. Isolated dotted five/six/seven
need not be nonadjacent. We now retain all their feet together.

\paragraph{The summed-foot identity and its complete error bounds.}
Let $w$ be a concave formal prefix reference, with nonnegative
coefficients $\beta_i$ in the spatial duals $\omega_i$.
At each selected isolated dotted index use its actual pressure,
put $Z=\sum\beta_i u_i n_i$ over those indices and $W=w-Z$.
Since $w.n_i=\beta_i$ and distinct old root products are
nonpositive, $Z^2\le\sum\beta_i^2u_i^2$. Summing the actual
pressures and completing the single square gives effective energy
\[
 L_W\ge w^2+\sum_i
 \{2\beta_i[T_i-(q_i-c_1)]+(2u_i-u_i^2)\beta_i^2\}.
\]
An actual D-position-five floor instead costs
$2\beta_5(c_D-c_1)$. These are pressure certificates for the
same $W$, not separate scalar lower bounds to add later.
All dotted corrections have support through eight and sum zero;
they preserve the late endpoint coefficient and vector sum. Their
mutual cross terms have a favourable sign, even for adjacent dots.

Concavity gives $\beta_5>.1,\beta_6>.012,\beta_7>.0095$;
the latter follow from $(d_5-d_6)/\sqrt2$ and
$(d_6-d_7)/\sqrt2$. Omitting a neighbour increases them.
Write $r(u)=u-u^2/2$. Each signed summand is
$2\beta[c_1-\operatorname{loss}(u)+\beta r(u)]$.
For $u\le3/4$, $r(u)\ge(5/8)u$. At five, $a_5<.163$ and
$\beta\ge.1$ bound the effective loss by
$.1005u+\sqrt{u(1-u)}<.5528<c_1$.
At seven, $a_7<.137,\beta\ge.0095$ similarly give
$.1310625u+\sqrt{u(1-u)}<.56983125<c_1$.
For $u\ge3/4$ the same expressions decrease, so both summands
are nonnegative throughout their actual intervals.

At six the summand increases with $\beta\ge\beta_0=.012$:
its derivative is positive because
$\operatorname{loss}(u)-2\beta_0r(u)$ is bounded by
$.133u+\sqrt{u(1-u)}<.57095<c_1$ on $u\le3/4$,
and decreases thereafter. At $\beta_0$ the effective loss is
$<.1405u+\sqrt{u(1-u)}<.5752$, again with a decreasing tail.
Its entire debit is therefore
\[
 2(.012)(.5752-4/7)=99/1093750<.000091.
\]
All maxima used here are the elementary identity
$\max[au+\sqrt{u(1-u)}]=(a+\sqrt{1+a^2})/2$.
No individual dotted-six floor was assumed.

\paragraph{The early reference and the actual terminal prices.}
First suppose five is not a weak helper. Keep formal five, whether
ordinary, D-position or isolated dotted, and keep dotted six/seven.
Omit only nonadjacent weak-helper floors. Through a strong fifteen/
sixteen their entire loss is bounded by the five decreasing terms
$j=6,8,10,12,14$, whose sum is $<.000422$.
For D-position five, $\beta_5<.11425$: the first mean is
$<.366402$ and the next retained slope, at six or seven, is
$>.20484$. Thus its separate debit is
$<2(.11425)(.57549-4/7)<.000929$.
The common initial-block change from $.61$ adds
\[
 (27/700)\sqrt{2/5}-[(.61)^2-(4/7)^2]/10>.0198377.
\]
The old exact $E_{15}>.92564,E_{16}>.94176$, minus all three
conservative errors, give a paid $W$ supported through $K$, of
sum $G_K$ and effective energy
\[
 L_W>Q_{15}=.944035,\qquad L_W>Q_{16}=.960155.
\]
Subtracting both dotted and D-position errors is conservative even
when their five cases cannot coexist. The true individual $P_{15}$
coefficient is at least $d_{14}>9/70$, and $P_{16}$ at least
$d_{15}>1/8$; the early corrections do not change them.

If five a weak helper, its actual scalar point gives $x_1>.530$:
subtract at most $(.58-4/7)\sqrt2(6/5)$ from $.5455$.
Six is strong or isolated dotted. Interpolate through
$(1,.530),(6,G_6),(K,G_K)$ and use its actual summed dotted-six
foot if necessary. The first-six energy is $>.74018$ because
$G_6>2.0454$; the fifteen/sixteen tails are $>.225,>.2404$.
Even after $.000091$ these exceed both $Q_K$.
The endpoint coefficient $1/(\sqrt K+\sqrt6)$ also exceeds the
same required prices. Thus every initial branch has been covered.

For an actual normalized pressure flat through sixteen, with demand
$G$ and full norm square $\le R$, the positive-dual payment gives
\[
 e>L_W+2\tau(G-G_K)-\tau^2(R-K).
\]
Use $G_{15}<3.468923,G_{16}<3.595940$.
The legal price is $\tau=1/8$ at both endpoints, except native-eight
at the centre, where $\tau=9/70$ is used at fifteen.
Effective energy is never identified with $W^2$.

\paragraph{Actual D4 at seventeen and D3 at eighteen.}
At seventeen, the complete $D_4$ leaf-three/bridge classification
and the old affine principals leave fork-three owns five/seven/eight.
Owns at least nine have ratio
$\sqrt2(\cos(\pi/9)+1/4)>1.68>1.673$.
Own five has its actual left $H_4$ point,
$A^2=17\alpha^2-1/2$, $G>4,R<17$, flat through fourteen.
It uses the old $.61$ early eleven/twelve contract in both initial
branches, not the new sixteen-flat formula. Its two totals exceed
$1.04$. For seven/eight use the actual future core
\[
 c+g^{-1}n_{17}+(2g)^{-1}n_{16},\quad k=1-3/(4g^2)>0.
\]
The cap/sixteen products vanish and the seventeen product is negative.
For $.9<g<.924$, $A/\alpha>4.09$, $\sqrt{k}/\alpha<.21$,
so $G>3.88,R<17$. It is flat through fifteen, and the old
eleven/twelve price $1/7$ gives both totals $>1.0086$.
No incorrect sixteen-flatness has been used.

At eighteen the right $D_3$ finite list is middle-three, end-three/
four/five, with every full bridge retained. Middle-three and end-four
have their actual affine/direct alternatives; end-five is directly
strong even at own zero. A full middle bridge has residual $1/36$
and is directly strong at zero own. An end-three bridge has residual
$1/8+1/36$, zero-own cycle and positive-own ratio $>1.9$.
The no-left end-three case has the old own-zero/three/four/six/null
principals. Own five has its left $H_4$ point, $G>4,R<18$,
flat through fifteen; the old early reference gives both totals
$>1.022$. Owns $g\ge43/40$ are directly strong.
For $.9<g<43/40$ use
$c+g^{-1}n_{18}+(2g)^{-1}n_{17}$, with
\[
 A^2=18\alpha^2-5/8,\quad
 A/\alpha>4.21,\quad
 k/\alpha^2<519200/5193841<.317^2.
\]
This actual core is flat through sixteen, $G>3.893,R<18$.
The new paid formula exceeds $1.0031$ at both early endpoints.
Thus eighteen has only helper or strong as well.

\paragraph{The complete actual centre nineteen.}
Deleting nineteen retains $A_{18}+2A_1$. With own $g$, one left
contact $\eta$ at $j$ and fork contacts $\delta,\epsilon$, the
data are
\[
 \Delta=1-q_L-\delta^2-\epsilon^2,\quad
 \sigma=g+\eta j/19+(\delta+\epsilon)/2,\quad
 C=\sqrt{38}\sigma/\sqrt\Delta.
\]
The complete finite-white list is the same actual high-rank tree
list already proved: left end-three/four or next-to-end-three;
with both forks only the full endpoint $D_{21}$ bridge; with one
fork only left endpoint-three and fork-three/four.
The full and fork-four bridges have residual $1/38$ and minimum
numerators $1/2+1/38$, $\sqrt2/4+1/38$, directly strong.
Fork-three bridges have residual $21/76$; own zero/three give
the cycle/triangle and all larger owns have ratio $>1.8$.
With no fork contact, the rechecked left and isolated inventory
applies. Its helper feet use $\omega_{20}+\omega_{21}=\sqrt2P_{20}$
only as a spatial identity, retaining both actual forks.

With no left contact the white list is one arbitrary finite fork,
or pairs $(3,3),(3,4),(3,5)$. At own zero, fork-three/four/six
have their actual affine principals. Fork-five uses the actual
left $H_4$ point, $R=18$, $A^2/\alpha^2=5+8\phi$ and
$\sqrt{\delta_5}/\alpha<15/56$, giving $G>3.96$.
Its corrected rational premise is $\phi>79/49$, as proved above.

For fork at least seven, put $b=\delta^2$. The actual core
$c+\delta^{-1}n_f+2(\delta^{-1}-\delta)n_{19}$ has selected
products zero/negative, positive time and negative square.
Its full signed data are
\[
 R=29-16/b+6/b^2,\quad
 G=\sqrt{21-2/b}-\sqrt{-8+14/b-6/b^2}.
\]
Monotonicity of $R$ and the quadratic-radical maximum give this
complete native split:
\[
\begin{array}{c|c|c|c}
\text{fork label}&b&G>&R<\text{ or }\le\\\hline
7&(.81,.82)&3.91&18.42\\
8&(2+\sqrt2)/4&3.91128&18.49034\\
9&(.88,.89)&3.925&18.60\\
\ge10&(.9,.95]&3.948&18.81\\
\ge10&[.95,1)&4.04&19 .
\end{array}
\]
For seven the two roots are $>4.304,<.388$; for nine,
$>4.327,<.402$; for the last two ranges, $>13/3,<.385$ and
$>4.346,<.298$. The original native cosine bounds follow from
their triple-angle polynomials and positive squares.
The critical eight row is not replaced by a coarse continuous bound:
\[
 R=109-64\sqrt2,\qquad
 G=\sqrt{13+4\sqrt2}-\sqrt{68\sqrt2-96}.
\]
The rational enclosure $1.41421356<\sqrt2<1.41421357$ gives
roots $>4.319358,<.408078$ and the displayed $R$ upper bound.
At fifteen the price $9/70$ gives energy $>1.00008648$;
at sixteen the price $1/8$ gives $>1.00007843$.
Every other row gives both totals $>1.00085$.

At own three, fork-three gives the triangle. Fork-four/five instead
gives the actual finite $D_{21}$ contacted-fork endpoint,
$C_f=(19+21\delta)/2>16$, signed direction $\sqrt2\omega_f$,
$R=21,G>\sqrt{315/17}>4.30$. Fork at least six is direct.
Own four/six/parallel retain their actual principals.
For own at least five and fork at least four, only own-five/
fork-four is not directly strong; its left $H_4$ point has
$\alpha=\sqrt2g+1/2$, $A^2=19\alpha^2-1/2$, $G>4.25,R<19$.
Own-five/fork-three has $A^2=19\alpha^2-3/4$, $G>4.2,R<19$.
Their lowest selected old root is seventeen, so both are genuinely
flat through sixteen. For remaining finite/dotted own with fork-three,
$g\ge6/5$ is direct; on $.9<g<6/5$ the actual core
$c+g^{-1}n_{19}+(2g)^{-1}n_{18}$ has $R<19$,
$A/\alpha>4.326$ and $\sqrt{k}/\alpha<.34$, hence $G>3.98$.

For two forks, own three has its triangle and larger owns are direct.
At own zero pair-three-three has its square; pair-three-four has
an actual finite $B_{21}$ mate with
$C_f=(38+21\sqrt2)/4>16$, again $G>4.30,R=21$.
Pair-three-five has its actual finite $H_3$ foot with
$R=19,q=2\phi^{-6}<1/8$ and
$G>\sqrt{19-1/8}/(1+\sqrt{1/152})>4$.
Every row is excluded by the stated positive-dual early formula.
The actual list is exhausted and nineteen is helper or strong.

\paragraph{The final summed reference.}
One of eighteen/nineteen is strong. Weak helper five gives effective
energy $>.74018+.2679-.000091>1.0079$, so is impossible.
Otherwise the same formal reference and summed-dot correction apply
through that endpoint. Their finite nonadjacent hole loss is
$<.00044$, by appending the sixteen/eighteen terms as needed.
The same $.000091$ and $.000929$ errors remain valid. Consequently
\[
 L_{18}>.9895677,\qquad L_{19}>1.0030777.
\]
Thus nineteen cannot be strong. In the weak-helper case at nineteen,
the paid eighteen reference has sum $G_{18}$, true retained demand
$G_{18}$, first-three mean $a=G_5/5<.367$ and unchanged endpoint
coefficient $>d_{17}>1/9$. The summed dots only affect coordinates
five through eight. Hence $a/\sqrt2<.260<.267$ and
$\sqrt2G_{18}=38/7$. The actual helper uses the physical long
$n_{19}$ and both signed fork directions; its $P_{20}$ coordinate
sum is not declared a long old dual.
All hypotheses of Lemma~\ref{app:tail:lem:weak19packet} now hold,
with the stronger $L>.9895677$. Its single paid full-foot square
contradicts the norm budget. All summed errors and terminal prices
have been charged once, with no arithmetic or unproved dotted floor.
\end{proof}

\subsection{The complete exact-twenty-one cover and its consumers}
\label{app:tail:exact21}
\begin{proposition}
\label{app:tail:prop:exact21}
Dimension twenty-one is impossible.
\end{proposition}
\begin{proof}
The common ordinary/debit costs exclude a forest without a dominant
block, and two dominant blocks cost more than one. After establishing
uniqueness, write the dominant rank as $r$ and spectator rank as
$S=21-r$. The complete three spectator budgets remain
$S/14,(S-1)/14,.09+(S-2)/14$, with their actual incompatible
defect definitions. A dominant $A$ has $S\le7$ and hence
$14\le r\le20$. Modern written component bounds make its lower
ranks immediate: use $5/7,.826,.86,6/7,.91$ at
$r=14,15,16,17,18$ respectively. With the weakest short-discard
spectator costs $6/14,5/14,4/14,3/14,2/14$, every total is
strictly greater than one. Typed and fully ordinary cases are
stronger. Thus the older detailed rank-fifteen/sixteen spectator
partitions need not be reproduced as separate consumer arguments.

At $A_{19}$, an ordinary irreducible rank-two remainder costs
$>1/7$ beside $>6/7$. Two ordinary $A_1$'s cost $>.2$;
actual $I+A_1>.2$ is a single joint bound, and $J$ is stronger.
Both remaining typed exceptions are exactly
Lemma~\ref{app:tail:lem:A19typed}. At $A_{20}$, the sole
ordinary or actual-$I$ rank-one remainder is excluded by
Lemma~\ref{app:tail:lem:A20A1}; $J$ and typed exceptions need
two ranks. This exhausts the $A$ side at its own rank.

For $B/D$, $S\le5$ and $r\ge16$. The valid D-admissible
$M=r-6$ floors at sixteen/seventeen/eighteen are
$>.823735,.8472225,.8689525$. Their weakest short-discard
spectators cost $4/14,3/14,2/14$, already too much.
At $B_{19}$, its genuinely valid $M=14$ improved floor
$>.91211125$ closes an ordinary rank-two block, a two-$A_1$
remainder with its at least one nonshort cost $>.1$, or either
typed block $>.09$. At $D_{19}$, its valid $M=13$ floor
$>.88864$ closes an ordinary irreducible remainder. A two-$A_1$
remainder costs $>.2$ if ordinary, $>.2$ jointly if it contains
actual $I$, or $>.34$ if it is $J$. The two typed cases are
the separately proved actual D4 continuation in
Lemma~\ref{app:tail:lem:proper21}.

At rank twenty even the weaker common floor $>.90715$ contradicts
an ordinary nonshort $A_1$ cost $>.1$. The remaining entire
$B_{20}+I$ is covered by Proposition~\ref{app:tail:prop:classicalI};
$D_{20}+I$ is covered by Lemma~\ref{app:tail:lem:proper21},
after its own missing D4 layer was supplied. At rank twenty-one
there are no spectators. Propositions~\ref{app:tail:prop:wholeB21}
and~\ref{app:tail:prop:wholeD21} exclude those whole diagrams.
Every finite farthest forest of rank twenty-one is therefore
impossible. There is no appeal to face inheritance or to a
necessary cover computed at another dimension.
\end{proof}

\subsection{The proper classical forests in exact dimension twenty}
\label{app:tail:proper20classical}
\begin{lemma}
\label{app:tail:lem:proper20classical}
The entire forests $B_{18}+N_7/O_7$, $D_{18}+N_7/O_7$,
$B_{19}+I$ and $D_{19}+I$ are impossible in dimension twenty.
\end{lemma}
\begin{proof}
We use the old loss $.61$ and true actual-rank continuations.
For $B_{18}$, the valid bulk reaches fourteen. At fifteen the
right $B_3$ permits long-three or the sealed short-three $F_4$.
The latter has residual $1/4$, numerator $g+1/(2\sqrt2)$:
own zero has its actual affine principal and every positive own
has ratio at least $1+1/\sqrt2$. All full endpoint bridges are
finite $B_{18}$, foreign-sealed and directly strong, including
their zero-own ratio. Thus coverage/nonadjacency reach fifteen.

For $D_{18}$, the valid bulk reaches thirteen. At fourteen the
complete right $D_4$ list and its affine/direct branches leave
only pure own-five/fork-three. Its actual left $H_4$ point has
$\alpha^2=\phi+5/8$, $A^2=14\alpha^2-1/2$, $R<14$ and
$G>3.62$, using $\alpha^2>56/25$ and residual ratio $<.091^2$.
It is flat through \emph{eleven}, so choose an actual strong
ten/eleven, not twelve. The early positive-dual formula and the
separate typed energy $>.09$ give both totals $>1.085$.
Both initial-five branches admit those same lower reference bounds.
The finite white list retains both forks, every foreign cost and
every full bridge; the non-direct role has $Q=0$ by the strong
foreign estimate, not by an assumption about the typed block.

At fifteen the right $D_3$ permits middle-three or end-three/four/five.
Middle and end-four have their actual affine/direct branches;
sealed end-five is strong even at own zero. Full bridges retain
their zero-own cycles and direct positive-own bounds, with no
finite extra branch at the internal cap. The end-three no-left
role has residual $5/8-Q$, numerator $g+1/4$; positive $Q$ and
$g\ge1$ are strong and own zero/three/four/six/parallel have
their actual principals. Own five gives its actual left $H_4$
pressure $G>3.74,R<15$, flat through twelve.
For finite $.9<g<1$, average both actual $A_2$ cores on sides
fourteen/sixteen, retaining their full cross $8g^2/9$.
With $d=4g^2/3-1$ their data are
\[
 A^2=15\alpha^2-5/8,\quad
 Y^2=A^2+1-14g^2/9<15\alpha^2,\quad
 x.Y\ge A-\sqrt d.
\]
The bounds $A/\alpha>\sqrt{15-125/529}>3.842$ and
$\sqrt d/\alpha<1/3$ give $G>3.5,R<15$.
At actual strong eleven/twelve, the pure-point formula plus the
orthogonal typed cost gives both totals $>1.099$ or $>1.030$,
respectively. Each core is a positive combination of actual walls,
future feasible against the entire typed block and every other wall.
Thus through fifteen only pure helpers can be weak.

The prefix-gain \emph{calculation}, now with that actual extended
coverage, gives $e_{B/D}>.90715+.00496125=.91211125$.
If five weak its actual reference is stronger; if fifteen strong
the extra increment pays the gain; otherwise the actual weak-fifteen
packet supplies the legal $x_1>.424$ modification at strong fourteen.
The physical long roots through sixteen exist in both components.
We have not applied the old $M=14$ theorem below its original ranks.
Adding either typed cost $>.09$ excludes all four typed forests.

For entire $B_{19}+I$, the valid bulk reaches fifteen. Its actual
$B_3$ layer at sixteen has the same long-three/sealed-$F_4$/full-bridge
list as just proved, now at its physical rank; all remaining weak
caps are pure helpers. Their actual nonadjacency extends through
sixteen. The scalar argument in
Proposition~\ref{app:tail:prop:classicalI} therefore applies after
this independently supplied geometric window: weak five has the
paid $>.9851486$ reference and unused-$I$ gain $>.015$; strong
five has the flat fifteen/sixteen reference $q>.9240775$, true
$P_5$ demand $f_5$ and individual price $>.146$.
Deleting the head retains actual $B_{18}$, with precisely the
isolated/long-two-three whitelist and only the old $I$ outside.
The q-free local head/$I$ gain $>.07665$ gives the contradiction.
This is a proof at supplied actual geometry, not the generic
$B_r\ge20$ theorem asserted at $r=19$.

Finally consider entire $D_{19}+I$. Lemma~\ref{app:tail:lem:D4proper}
first supplies coverage through fifteen. At sixteen the right $D_3$
has the same complete finite, foreign, affine and bridge list as above.
Only pure own-five or finite $.9<g<1$ in the end-three row remain.
Own five has its actual left $H_4$ point, $G>3.87,R<16$, flat
through thirteen; the bare eleven/twelve formula already exceeds
$1.026$. For finite owns use both actual $A_2$ cores on sides
fifteen/seventeen, with their full cross, $A^2=16\alpha^2-5/8$,
$R<16$ and $G>3.63$, since
$A/\alpha>\sqrt{16-125/529}>3.97$ and the radical ratio is $<1/3$.
If five weak, its actual eleven/twelve references have energies
$>.9104,>.9314$, so their pure-point totals exceed $1.020$.
If five strong, the flat-five reference can first use the single
initial-$I$ adjustment: from $t+(3/2)s>3/4$ its joint gain is
\[
 s^2+(5/4)(t-.361)_+^2>151321/3050000>.0496,
\]
with a legal clip at $.6$ when necessary. Its support first five
plus $I$ is orthogonal to the actual tail residual. It uses only
the available $P_5$ price, leaving the early terminal price intact.
The pure-point formula plus this one joint gain gives both totals
$>1.007$. Thus coverage/nonadjacency reach sixteen in either branch.

The same independently checked scalar bootstrap and flat-reference
argument now prove five strong and give $q>.9240775$.
Deleting long one retains actual $D_{18}$, so the same local head/$I$
kernel, its selected long roots and both signed forks apply and give
$q+.07665>1$. Earlier tail contradictions were used only to remove
their branches; no spent $I$ gain is carried into this new final
reference. This whole-forest argument stops before centre seventeen;
its older mixed-$H_4$ catalogue remains preserved but is unnecessary
for this composition with the later local head theorem.
\end{proof}

\subsection{The two rank-one remainders beside $A_{19}$}
\label{app:tail:A19rankone}
\begin{lemma}
\label{app:tail:lem:A19ordinary}
An actual $A_{19}$ has projection energy $>563/625=.9008$,
with arbitrary finite spectators. In particular it cannot coexist
with an ordinary nonshort $A_1$.
\end{lemma}
\begin{proof}
Suppose $e_A\le.9008<.91$. The complete $N=20$ window and near-four
comparison above apply. Zero or one weak end already has reference
energy $>.9522798$ or $>.9233277$. Thus both five/fifteen are weak.
The paid rational reference $w$ has energy $E>.8944=559/625$,
first six coordinates $H_6/(6\sqrt2)$ and last six their negatives,
where $H_6=2.288$. Its internal weak holes have already been paid.
For the actual weak cap at five, either helper assignment gives
\[
 A^2=(15/2)\gamma^2+C\gamma-21/40,\quad C\in\{1/2,3/2\},
 \qquad K=w.v={143\over750}(10\gamma+1).
\]
The positive corrections below use only old $n_3,n_4,n_5$, so have
zero product with this flat reference. For $Y=v-u$ we therefore
have $w.Y=K$.

For a finite selected correction, $S^2=Y^2=A^2+\delta$,
and its actual finite projected point gives the normalized pressure
$x_A.(Y/S)\ge g_0$, with
$g_0^2=(S-\sqrt\delta)/(S+\sqrt\delta)$.
For a future core of square $-d$, $S^2=A^2-d$ and
$g_0=(A-\sqrt d)/S$.
We prove the full overlap gap $g_0-K/S>.08$ on every own range.

Native five uses the actual left $H_4$ correction
$\gamma(3n_5/2+n_4+n_3/2)$, with
$\delta=(5-3\phi)/8$. The smaller correlation $C=1/2$ gives
$S^2=79/40+7\phi/4>2.19^2$, $\sqrt\delta<.14$ and
$K<13013/7500$. Hence $K/S<.8$ and
$g_0^2>93/107>.93^2$, a gap $>.13$.
Native seven uses $\gamma(4n_5/3+2n_4/3)$,
$d=4\gamma^2/3-1>0$. Its selected products are zero/negative
and its time is positive. From $.9<\gamma<.902$ we have
$A>2.449$, $\sqrt d<.292$, $K<1.911$ and
$S^2=(37/6)\gamma^2+C\gamma+19/40<977/120<9$.
Thus $(A-\sqrt d-K)/S>.246/3>.08$.

Every finite own at least eight uses the actual finite pair correction
$u=\gamma n_5$, $\delta=1-\gamma^2$.
Here $S^2=(13/2)\gamma^2+C\gamma+19/40>2.5^2$,
$\sqrt\delta<.4$, so $g_0^2>21/29>.85^2$.
The exact lower overlap polynomial is
\[
 .77^2S^2-K^2\ge
 {1573\over7200}\gamma^2-{387563\over900000}\gamma
                +{4414927\over18000000}
 >.218(\gamma-1)^2+.005\gamma+.027>0.
\]
Thus $K/S<.77$ and the gap is $>.08$ on the whole finite interval.

For dotted own $1<\gamma<27/25$, use the actual future core
$c+\gamma^{-1}n_5$. It has zero cap product, negative old-five
product and positive time. Now
$S^2=(15/2)\gamma^2+C\gamma-61/40+\gamma^{-2}>7.325>2.7^2$
and $\sqrt d<.38$. Since $A>S$, $g_0^2>58/77>.86^2$.
Multiplying the overlap difference by $\gamma^2$ gives the bound
\[
 \gamma^2(.77^2S^2-K^2)>
 .0319+.069(\gamma-1)+2.632(\gamma-1)^2
       +2.813(\gamma-1)^3+.811(\gamma-1)^4>0.
\]
The rounded quartic coefficients come from
$146047/180000,-387563/900000,-16929473/18000000,.5929$,
so every coefficient comparison is rational. This gives a gap
$>.09$ for the entire dotted interval.

All selected finite points or future cores are feasible against
every actual wall by their positive corrections, and have pure
$A$ displacement because $Q=0$ was proved in the inventory.
For their unit $y=Y/S$, pay $\mu=.08$ in the same vector:
\[
 e_A\ge E+2\mu(g_0-w.y)-\mu^2
 >559/625+4/625=563/625.
\]
This contradicts the contrary target. A separate nonshort $A_1$
costs $>.1$, giving total $>1.0008$.
\end{proof}

\begin{lemma}
\label{app:tail:lem:A19I}
An actual $A_{19}$ cannot coexist with an isolated $I$, with
arbitrary further finite spectators.
\end{lemma}
\begin{proof}
Write the actual $I$ parameters as $s,t$, so $y_I\ge s$,
$t+(3/2)s>.767$ and both normalized end heights are at least $t$.
Put $a=\sqrt{19/20}>.974$, giving $x_1,-x_{20}\ge at$.
We first remove every near-four without assuming $e_A<.91$.
Its actual mate has index $j=4,5$ and height $>.86$; at the
opposite end there is a mate or a strong $k\ge14$ of height $>.77$.
Except $(j,k)=(5,14)$, the exact dual correlation is at most $1/3$.
The positive vector $(543/800)u_j+(435/800)u_k$ has pressure
$>80193/80000$ and norm square at most that same value.
One completed square gives energy greater than one.

At the exceptional pair, the actual heights exceed $2.355,2.231$;
their positive squared margins are $.000975,.002999$.
The concave reference through those heights at five/fourteen has
\[
 E=\tfrac12[2.355^2/5+.124^2/9+2.231^2/6]>.97.
\]
If $s\ge.18$, its orthogonal $I$ energy already gives $>1.0024$.
If $s<.18$, the barrier gives $t>.497$ and $-x_{20}>.484$.
The last-six mean is $<.264$; its legal endpoint modification adds
$>(6/5)(.484-.264)^2=.05808$. The remaining mean is $>.217$
and the middle slope $<.01$, so the new reference stays concave.
This also exceeds one. Reflection treats every orientation. The
complete weak window is therefore pure-helper and nonadjacent.

Use the same full rational curve and omission table at $N=20$.
If zero or one ends are weak, its energy is at least the weaker
bound $>.9233277$. The two endpoint block means are $<.302$ and
each length is five or six. Put $z=(at-.302)_+$.
Their disjoint legal changes add at least $(12/5)z^2$.
Weighted Cauchy with the single $I$ height gives
\[
 s^2+(12/5)z^2>
 {(a\cdot.767-.302)^2\over(9/4)a^2+5/12}
 >{.445058^2\over511/200}>.0775.
\]
The exact denominator is $613/240$; the last cleared margin is
$.000064123364$. Thus both end indices must be weak.

At five write the actual cap as $v=\gamma\omega_5+\omega_e/2$,
$e=1$ or nineteen, with $4/5<\gamma<27/25$ and $Q=0$.
Its time is $A^2=(15/2)\gamma^2+C\gamma-21/40$, with
$C=3/2$ for the near helper and $1/2$ for the far helper.
Native five uses its actual $H_4$ projection with pressure
$H(A,2/3-\gamma^2)>A-.111$.
For finite owns at least seven, the positive actual core
\[
 c+\gamma n_5+2(1-\gamma^2)n_e
\]
has cap/five products zero, helper product $3/2-2\gamma^2<0$,
square $-(1-\gamma^2)(4\gamma^2-3)$ and positive time.
Its radical is at most $1/4$. For dotted own use
$c+\gamma^{-1}n_5$ and radical $\sqrt{1-\gamma^{-2}}$.
Positive old corrections allow the common left pressure
$\gamma h_5+h_e/2$ to retain each actual demand. Every other
wall, including the actual $I$ wall, is still satisfied.

A near helper forces $x_1>.65$ on all three ranges.
For native five use $A>2.365,\gamma<.81$ and
$\sqrt2[2.365-.111-.81(2.129)]>.65$.
For finite owns at least seven, $A-\gamma(q-.61)$ increases
from $\gamma=.9$, where $A>2.6267$; subtracting its radical
$\le.25$ still gives $x_1>.65$.
For dotted own the deficit decreases through $1<\gamma<1.08$:
$A'<3$, the radical derivative exceeds two, and $q-.61>2.12$.
At $1.08$, $A>3.1373$, radical $<.37771$, so
$\sqrt2[3.1373-.37771-1.08(2.129)]>.6508$.
The weak end has six-block mean $<.270$; this endpoint gain is
$>.17328$, impossible beside $>.8944026$. Neither weak cap
can therefore have a near helper.

For either far-helper cap, divide its full actual pressure by
$\gamma>.8$. Native five gives $(2.188-.111)/.81>2.39$.
Finite owns at least seven have increasing $A/\gamma$, initially
$\sqrt{200/27}>68/25$, and radical ratio at most $5/18$.
For dotted own, $A/\gamma-\sqrt{\gamma^2-1}/\gamma^2$ decreases;
its first derivative term is $<.11$ and its subtracted derivative
$>1$. At $1.08$ it exceeds $(2.9602-.37771)/1.08>2.39$.
Since $h_e>0$ and $1/(2\gamma)<5/8$, both actual feet give
\[
 h_5+(5/8)h_{19}>2.39,\qquad
 h_{15}+(5/8)h_1>2.39.
\]
The full finite and continuous dotted alphabets have been covered.

To spend these pressures together, set $A_0=.556,B_0=2.0425$,
$C_0=2.288,r=5/8$; then $B_0+rA_0=2.39$.
Build a symmetric reference with heights $A_0,B_0,C_0$ at
one/five/six, then the old central heights and reflection.
Its initial slopes are $.556,.371625$ repeated four times,
$.2455,.118,\ldots$, decreasing. Delete actual internal weak
knots seven through thirteen; their entire loss is $\le.001296$.
The positive end coefficients in the spatial duals are
\[
 \beta_1={5A_0-B_0\over8}=59/640,\quad
 \beta_5={5B_0-A_0-4C_0\over8}=1009/16000,\quad
 \eta=\beta_1-r\beta_5=1351/25600>0.
\]
All remaining coefficients use actual strong floors. Direct rational
calculation of that same central segment gives energy $\ge.021912$,
so the reference $w_A$ has
\[
 E_A\ge A_0^2+(B_0-A_0)^2/4+(C_0-B_0)^2+.021912
       =.9437388125.
\]
Let $\lambda=\sqrt{19/10}>1.3784$.
The single actual-$I$ barrier gives
$h_1+h_{19}+3\lambda s>2\lambda(.767)$.
Use the one full vector $W=w_A+3\eta\lambda e_I$.
Its two foot coefficients are $\beta_5$, and the remaining end
coefficient is exactly $\eta$ times that one barrier. Consequently
\[
 x.W>E_A+2\eta[\lambda(.767)-A_0],\qquad
 W^2=E_A+(171/10)\eta^2.
\]
The reference values $A_0,B_0$ are used through these positive
combined pressures. No separate endpoint floor is required there.
The rational bounds
$4\eta[\lambda(.767)-A_0]>.105807111375$ and
$(171/10)\eta^2<.047625$ give
\[
 \norm x^2\ge2x.W-W^2>
 .9437388125+.105807111375-.047625
 =1.001920923875>1.
\]
This is one spatial certificate and one $I$ coordinate. Every
internal hole and every foot overlap has been charged in that vector.
\end{proof}

\begin{figure}[ht]
\centering\setlength{\unitlength}{1mm}
\begin{picture}(145,40)
\put(0,28){\framebox(45,10){\scriptsize $h_5+\frac58h_{19}>2.39$}}
\put(50,28){\framebox(45,10){\scriptsize $h_1+h_{19}+3\lambda s>1.534\lambda$}}
\put(100,28){\framebox(45,10){\scriptsize $h_{15}+\frac58h_1>2.39$}}
\put(22.5,28){\vector(2,-1){32}}
\put(72.5,28){\vector(0,-1){16}}
\put(122.5,28){\vector(-2,-1){32}}
\put(32,21){\makebox(0,0){\scriptsize $\beta_5$}}
\put(77,20){\makebox(0,0){\scriptsize $\eta$}}
\put(112,21){\makebox(0,0){\scriptsize $\beta_5$}}
\put(40,2){\framebox(65,10){$W=w_A+3\eta\lambda e_I$}}
\end{picture}
\caption{The two far-helper pressures and the single actual-$I$
barrier fund one vector. The arrows record positive pressure
coefficients, rather than Coxeter bonds.}
\label{fig:tail:A19Ipressure}
\end{figure}

\subsection{The whole twenty-dimensional $A$ row}
\begin{lemma}\label{app:tail:lem:wholeA20}
The entire vertex forest cannot be $A_{20}$.
\end{lemma}
\begin{proof}
Use the $N=21$ notation of Lemma~\ref{app:tail:lem:Awindow2021}.
That lemma and the rational central curve apply before any contrary
energy assumption: the weak internal knots are nonadjacent pure
helpers, and the near-four roles have already been excluded by norm.
Put $C=2.317$ and $B=2.150$. According as the end knots five and
sixteen are both weak, just one is weak, or both are strong, the
concave reference energies, including every actual internal hole, are
\[
 E_2>.92274,\qquad E_1>.9509563,\qquad E_0>.97977675.
\]
These are the $N=21$ curve rows above. No nonexistent strong end
is inserted into the reference.

An actual weak cap at five has own $\gamma$ and a simple helper
at one or twenty, with
\[
 A^2=(160\gamma^2+c\gamma-11)/21,
 \qquad c=32\ (\text{near}),\quad c=10\ (\text{far}).
\]
The native-five correction is
$\gamma n_5+(2/3)n_e+(1/3)n_{e'}$, with finite residual
$2/3-\gamma^2$. For finite owns at least seven use the core
$c+\gamma n_5+2(1-\gamma^2)n_e$, of negative square
$-(1-\gamma^2)(4\gamma^2-3)$; for dotted own use
$c+\gamma^{-1}n_5$. All coefficients are positive.
The finite row uses the projected point, the other rows a future
negative core. Every original wall inequality therefore holds by
acuteness. Their pressures on $U=v-w_{\rm old}$ exceed, respectively,
\[
 A-.111,\quad A-\sqrt{(1-\gamma^2)(4\gamma^2-3)},\quad
 A-\sqrt{1-\gamma^{-2}}.
\]
The finite and continuous dotted domains are those of the window
lemma; the latter has $\gamma<1.08$.

A near helper forces $x_1>.64$. For five use $A>2.365,\gamma<.81$.
For finite owns at least seven, $A-2.151\gamma$ increases from
$\gamma=.9$, where $A>2.649$; its radical is at most $.25$.
For dotted own the deficit decreases through $1.08$, where
$A>3.1636$ and the radical is $<.37771$.
Substituting these bounds into the actual pressure proves the
$.64$ floor throughout. The weak-end six-block mean is $<.274$.
Its endpoint graft gains $>(6/5)(.64-.274)^2=.1607472$, already
impossible beside $E_2$. Every weak end thus has a far helper.

Suppose both ends are weak. The reference is flat on its first
and last six coordinates, with endpoint coefficients
$\ge1549/12000$. The projection of a far foot is
$(5\gamma/6)\omega_6+\omega_{15}/12$.
The residual $Y$ has zero sum on both blocks and vanishes in the
middle, hence $w.Y=0$. Its variance and projected deficit satisfy
\[
\begin{array}{c|c|c}
\text{own}&Y^2&\Delta\text{ lower}\\\hline
5&(2/3)\gamma^2+1/12<1.1&.30\\
\text{finite }\ge7&4\gamma^4-(16/3)\gamma^2+29/12<1.1&.30\\
\text{dotted}&(5/3)\gamma^2-19/12+\gamma^{-2}<1.22&.32
\end{array}
\]
For five use $A>2.2$, radical $<.111$, giving $\Delta>.3319$.
For the other finite owns split at $.91,.92$: lower times
$2.4649,2.4937,2.5225$ and upper radicals $.2318,.2434,.25$
give deficits $.302266,.300158,.30305$.
The dotted deficit decreases through $1.08$, where $A>2.9794$
and radical $<.37771$ give $>.3233$. These are complete intervals.

Add $.13$ times both residuals to one reference. At either anchor
the combined cost is $<.13(59/60)=767/6000<1549/12000$.
Their cross is retained using
$\norm{Y_L+Y_R}^2\le2(Y_L^2+Y_R^2)$.
With zero, one or two dotted ends the gains are
$>.08164,.082784,.083928$. Each added to $E_2$ exceeds one.

If just five is weak, the six-block and opposite five-block
coefficients are at least $1549/12000$ and $.1315$.
The far-foot projection becomes
$(5\gamma/6)\omega_6+\omega_{16}/10$.
Its variance decreases by $1/60$ and its deficit by
$B/10-C/12=263/12000$. Thus nondotted owns have
$\Delta>.278,Y^2<1.1$, and dotted owns have
$\Delta>.30,Y^2<1.22$. Payment $.14$ costs $<.126$ and
$.014$ at the two anchors and gains $>.05628$ or $.060088$.
Both contradict $E_1$.

Both ends are now strong. The reference has first and last
five-block value $b<.305$, and both anchor coefficients $>.1315$.
Releasing the first root retains $A_{19}$. Finite classification
gives isolated, simple-three at $2,3,19,20$, or four at $2,20$.
Two contacts make a cycle; exceptional finite ranks are too small.
The complete nonisolated edge-cosh list is
\[
\begin{array}{c|c|c}
\text{contact}&\text{white residual}&\cosh L_1\\\hline
3\text{ at }2&21/40&(40\gamma+19)/21\\
3\text{ at }3&1/10&(40\gamma+18)/(2\sqrt{21})\\
3\text{ at }19&1/10&(40\gamma+2)/(2\sqrt{21})\\
3\text{ at }20&21/40&(40\gamma+1)/21\\
4\text{ at }2&1/20&(40\gamma+19\sqrt2)/\sqrt{42}\\
4\text{ at }20&1/20&(40\gamma+\sqrt2)/\sqrt{42}
\end{array}
\]
The minimum allowed owns are $1/2,0,1/2,1/\sqrt2,0,1/2$;
smaller owns have time-zero or actual affine/null obstructions.
Every row forces $x_1>.396$. All but far-three have cosh $>1.5$;
that last squared floor is
$200(\sqrt2-1)/[21(10\sqrt2+11)]>.396^2$.
It reduces to $10441790\sqrt2>14764031$, certified by
$\sqrt2>707/500$.
If both ends are nonisolated, their disjoint grafts gain
$2(5/4)(.396-.305)^2=.0207025$, contradicting $E_0$.

Reflect an isolated end to the left. Its time is
$A^2=(40/21)\gamma^2-1$.
Own zero, three and four have nonpositive time square; six has
the affine six--three path; parallel has the positive feasible
null sum. Only five, finite owns at least seven, or dotted remain.

For five take $\gamma(3n_1/2+n_2+n_3/2)$.
Its finite residual is $1-3\gamma^2/2$, and its displacement
is $\gamma\omega_4/4$. Its projected-point pressure exceeds
$.393$, since $A>.4966$, the second time root is $>.5146$,
and the residual root is $<.1352$.
Centering in the first five coordinates leaves deficit $>.0447$
and variance $\gamma^2/10<.06561$. Payment $.6$ costs
$<.0972$ and gains $>.0300204$.

For finite owns at least seven the actual core is
$c+\gamma^{-1}n_1+2(\gamma^{-1}-\gamma)n_2$,
of square $-[7-4\gamma^2-3\gamma^{-2}]$.
Its centered variance is $(28/5)\gamma^2-8+3\gamma^{-2}<.6$.
For seven, $.9009<\gamma<.901$ gives $A>.7388$, radical
$<.2395$, projected deficit $>.11187$.
For larger labels $\gamma>.92$. The radical is strictly
concave, because its squared argument has second derivative
$-8-18\gamma^{-4}<0$. Its derivative at $.92$ is $<.66$,
whereas $A'>1.38$; the deficit derivative is therefore
$>1.38-.66-.43>0$. Its initial value is $>.1222$.
The uniform deficit is $>.111$. Payment $.2$ costs $<.04$
and gains $>.0204$, contradicting $E_0$.

For dotted own, $\gamma>1$ is unbounded. The core
$c+\gamma^{-1}n_1$ has centered variance
$(8/5)\gamma^2-2+\gamma^{-2}$.
Put $u=\gamma^{-2}$, $D=40/21$, $q=\sqrt D$ and
$a=q-\sqrt{D-1}$, so $0<a<.43$.
Concavity gives $\sqrt{D-u}\ge q-au$, while
\[
 au+\sqrt{u(1-u)}\le(a+\sqrt{a^2+1})/2
                         <(.43+1.089)/2=.7595.
\]
Since $q>1.38$, the projected deficit divided by $\gamma$
is $>.1905$, including its $.430$ block payment.
Multiplier $1/(8\gamma)$ costs $1/40$ at the anchor;
the variance divided by $\gamma^2$ is at most $8/5$.
Its gain $>.022625$ contradicts $E_0$.
All corrections have positive coefficients on actual walls;
the finite rows use projected points, the negative rows future
cores. The full zero-, one- and two-weak-end partition is exhausted.
\end{proof}

\subsection{The whole twenty-dimensional $B$ row}
\begin{lemma}\label{app:tail:lem:wholeB20}
The entire vertex forest cannot be $B_{20}$.
\end{lemma}
\begin{proof}
Use the genuine unit $B_{20}$ roots, including the short root
$n_{20}=e_{20}$. Lemma~\ref{app:tail:lem:whole58local} covers
five through sixteen. At seventeen the retained short piece is
$B_3$: short-three own zero has its affine principal, and positive
owns have ratio $\ge1+1/\sqrt2$. At eighteen it is $B_2$:
short-three owns zero, four, six and parallel have their actual
weighted-star, affine and null exclusions; finite owns at least
seven are directly strong. Endpoint bridges have residual $1/36$;
their own-zero cycle and positive-own ratios exclude weakness.

For the remaining own-five/short-three role the left $H_4$
projection is flat through fifteen, has norm $<18$ and pressure
$G>4.1$. Old strong knots eleven/twelve have endpoint coefficient
$>1/7$. Old weak five would give actual $x_1>.5965$, strong six,
and a strong endpoint sixteen/seventeen already supplied by the
ordinary bulk and $B_3$ row. Its positive three-knot interpolation
then has energy
$.7600299945+2.2-(4/5)\sqrt6>1.0004379$.
Thus old five is strong. The old eleven/twelve references exclude
the $G>4.1$ row; even the safe norm bound twenty gives $>1.0102$.

The exceptional $B_2$ own-three role has an actual finite $D_{20}$
mate after releasing the \emph{short} root twenty. Positive
finite-Gram projection identifies the actual endpoint and gives
cosh $9+5\sqrt2>16$. Its pressure obeys
$S_{20}^2>300/17>4.2^2$. Thus $G>4.2,R=20$, flat through
both early endpoints, excludes this row by $>1.042$.
The weaker long-edge mate would not suffice.
We have complete $.58$ coverage through eighteen, with only
nonadjacent left helpers and the isolated dotted-five exception.
The unused last-long-edge catalogue is unnecessary.

Releasing head one retains $B_{19}$, so its finite-white list is
isolated or simple-three at long leaf two. The latter has
$A^2=2\gamma^2+2\gamma$, $C_1=2\gamma+1$.
Own zero has time zero, three the affine triangle; all remaining
owns force $x_1>.643$.
The old bulk gives a strong endpoint fourteen/fifteen. If old
five is strong, its first-five graft gains
$(5/4)(.643-.361)^2=.099405$ beside $q>.90715$.
If old five is weak, six is strong and $S_6>2.01815$.
Interpolation through $(1,.643),(6,2.01815)$ and that endpoint
has decreasing positive slopes, initial energy $>.7916565045$,
and final chord energy $\ge2.5-\sqrt{21}/2>.20871$.
This too exceeds one. The head is therefore isolated.
Time square, affine six--three and positive null sum exclude
owns zero, three, four, six and parallel. Only five, finite
at least seven and dotted remain.

Their actual head pressures, normalized to sum one, are
\[
\begin{array}{c|c|c}
\text{own}&Z&h\text{ with }x.Z\ge h\\\hline
5&P_4/4&>.393\\
\text{finite }\ge7&(1-u/2)(e_1+e_2)+(u-1)e_3&
[\sqrt{2-u}-\sqrt{(u-1)(4-3u)}]/\sqrt2\\
\text{dotted}&(1-u/2)e_1+(u/2)e_2&
[\sqrt{2-u}-\sqrt{u(1-u)}]/\sqrt2
\end{array}
\]
Here $u=\gamma^{-2}$ is respectively greater than one or in
$(0,1)$. The first is an actual $H_4$ projected point; the
others are the positive future cores on actual head roots proved
above. Write $R=Z^2$ in the last two rows. The uniform scalar
\begin{equation}\label{app:tail:eq:B20headscalar}
 .8h-.16R>.2893
\end{equation}
holds throughout both domains. For seven,
$.9009<\gamma<.901$ gives $h>.4315,R<.349$, hence $>.28936$.
For finite labels at least eight, $1<u<4-2\sqrt2<1.172$.
With $D=(u-1)(4-3u)$ the derivative of the left side is
\[
 { .8\over\sqrt2}
 \left[-{1\over2\sqrt{2-u}}-{7-6u\over2\sqrt D}\right]
                   +.64-.48u.
\]
It is negative up to $7/6$; above that point its three
contributions are $<-.28,.032,.08$. At the eight-endpoint
$\sqrt{2-u}>.910,\sqrt D<.289,R<.373$, giving $>.29136$.
For dotted owns use the chords $1.4142-.379u$ on $u\le.5$
and $1.4494-.4494u$ on $u\ge.5$.
Together with $R\le1-.75u$ and $R\le.75-.25u$ they give
\[
\begin{split}
.63987152-.0944u-.5657\sqrt{u(1-u)}&>.30587152,\\
.69978064-.2142u-.5657\sqrt{u(1-u)}&>.29023064.
\end{split}
\]
The circle bounds use
$\sqrt{.0944^2+.5657^2}<.5736$ and
$\sqrt{.2142^2+.5657^2}<.6049$. No sampled minimum is used.

We next prove fresh $.58$-strong five. An actual weak helper
five supplies the curved reference through $(1,.5455)$, six,
and all later strong knots, ending at seventeen/eighteen.
Its initial energy is $>.74385$, compared with the ordinary
$.58$ six-energy $<.7126$. Since $E_{17}>.97231$ and all
holes cost $<1/1024$, its energy is
$>.97231+.74385-.7126-1/1024>1.0025$.
The only other weak-five role is an isolated dotted cap.
A native-five head forbids helper five/six/seven: their actual
wall product would be positive, because
\[
 A_i^2-(\phi^3/2)(2\gamma+1)^2
       <(28\gamma^2-32\gamma-13)/5<0.
\]
Its product with the isolated weak dotted-five cap lies in
$(-.3,0)$, also forbidden. Thus only the other heads need be
considered with the latter cap.

Choose actual strong $J=6$ or seven and $K=17$ or eighteen.
Interpolate their genuine $.58$ floors and every strong knot
between, omitting and charging all weak knots. The resulting data are
\[
\begin{array}{c|c|c|c|c}
J&\text{initial mean }m&w^2&\text{available }P_J&Q_5\\\hline
6&<.3399&>.9530&>.1435&<77/150\\
7&<.31938&>.9354&>.1366&<194/175
\end{array}
\]
The squared-increment lower numerators from six through sixteen,
over $10^6$, are
\[
38461,33333,29411,26315,23809,21739,20000,18518,17241,16129,15151.
\]
Their sums from six and seven are $.260107,.221646$.
Hole costs are $<.000267,.000181$; endpoint eighteen adds
$>1/70$ and pays any seventeenth hole. The terminal slopes give
the displayed available coefficients.

The actual dotted-five foot is
$Z_5=P_5+(v/2)(e_6-e_5)$ with demand $T_5>1.8185,v>.4$;
its centered variance is $Q_5=5-25/J-v+v^2/2$.
Set $Y=Z-P_J/J$, $Y_5=Z_5-5P_J/J$.
The head support lies in the first three coordinates, so
$Z.Z_5=1$ and $Y.Y_5=1-5/J$.
The single vector $W=w+.4Y+.09Y_5$ has total anchor cost
$.85/J$, below both available coefficients. Its effective energy is
\[
 w^2+.2893-.8m+.16/J+.18(T_5-5m)
                     -.0081Q_5-.072(1-5/J).
\]
The two gains are $>.04930,.06699$ and both totals exceed one.
Fresh strong five has been proved before any assumed strong
eighteen floor.

All new references now start at five, with $b<.3652$ and
available $P_5$ coefficient $>.151$.
The head gains are $>.015456$ for five (pay $.56$ on
$P_4/4-P_5/5$), $>.02914$ for the other finite heads
(pay $.4$ and use \eqref{app:tail:eq:B20headscalar}), and
$>.03384$ for dotted heads. For the last, the pure dotted
scalar $>.03624$ at $\sqrt2b<.5106$ loses $<.0024$ when
only its linear block term changes to $\sqrt2b<.5166$.
No typed theorem is imported.
If eighteen is strong, its reference has energy $>.985613$
before the head, and even the least gain contradicts the norm.

Otherwise eighteen is an actual helper and seventeen is strong.
The reference has $q>.97133$, first-five mean $b<.3652$,
and available $P_{17}$ coefficient $>1/9$.
For that actual helper put $V=U-\sqrt2\gamma P_{17}$.
Its variance is $<2.1$, and its pressure deficit, including
the nonzero cross $w.V=b/\sqrt2$, is $>.21$.
For five, seven, eight and larger finite labels the respective
deficit seals are $>.3444,.2853,.2306,.2303$.
The last interval starts at $\gamma>.939$; its derivative is
positive since $A'>6$, $H_A>.94$ and its radical decreases.
The dotted deficit decreases through $6/5$, where
$A>7.3307$ and radical $<.5528$ give $>.2177$.
Thus every continuous parameter is covered.

Payment $\nu=1/16$ costs
$\nu\sqrt2\gamma<.107<1/9$ and gains
$>2\nu(.21)-\nu^2(2.1)=.018046875$.
Combine the head and this foot in one vector. If its initial
support has length $s=3,2,1$, the residual product is
$(\ell_s-.2)/\sqrt2$.
For head five $\ell_s=1/4$, giving total cross cost $<.002476$;
for the other finite heads $\ell_s\le1/2$ gives $<.010608$;
for dotted heads $\ell_s\le1$ gives at most $.02$ exactly.
The final totals are
\[
\begin{split}
.97133+.015456+.018046875-.002476&=1.002356875,\\
.97133+.02914+.018046875-.010608&=1.007908875,\\
.97133+.03384+.018046875-.02&=1.003216875.
\end{split}
\]
Every actual head case is impossible, proving the whole $B_{20}$ row.
\end{proof}

\subsection{The whole twenty-dimensional $D$ row}
\label{app:tail:sec:wholeD20}
The three possible isolated long heads have different paid references,
but share one finite-white inventory at the last centre. We prove
that inventory once, retaining the head-dependent cross products in
a small table. This also keeps the endpoint choice honest: one of
fourteen/fifteen and one of sixteen/seventeen is strong; no proof
below assumes all four floors.

\begin{lemma}[The early $D_{20}$ prerequisites]
\label{app:tail:lem:D20early}
For the entire forest $D_{20}$ the long-end cap is isolated, of
own five, finite at least seven, or dotted. Moreover
$S_5>\sqrt5-.61/\sqrt2$, without a hypothesis on that cap.
\end{lemma}
\begin{proof}
Use $n_i=(e_i-e_{i+1})/\sqrt2$ for $i<20$ and
$n_{20}=(e_{19}+e_{20})/\sqrt2$. Then
$\omega_i=\sqrt2P_i$ through eighteen and
$\omega_{19,20}=(P_{19}\mp e_{20})/\sqrt2$.
Write $s_i=\sqrt i-.61/\sqrt2$.
The old bulk applies literally through fifteen: right retained
rank is at least five, and all walls are actual. Its only weak
rows are nonadjacent pure far helpers.
For the head, retained $D_{19}$ admits isolated or simple-three
at its long leaf two: other leaves give exceptional finite rank
twenty, internal attachments an extra fork, and a four edge is
not a large finite $D$ extension. Two contacts give a cycle.
The nonisolated row has $A^2=2\gamma^2+2\gamma$,
$C_1=2\gamma+1$, exactly as in the preceding proof.
Its $x_1>.643$ contradicts either the old strong-five graft
or the old weak-five three-knot reference there. Those arguments
use only the first fifteen long coordinates, so the signed
forks do not alter a coefficient. The isolated list follows
from $A^2=2\gamma^2-1$, the actual affine six--three path,
and the parallel null sum.

We supply the head-independent strong-five bootstrap before using
any head gain. Suppose $S_5\le s_5$. Its actual far helper gives
$x_1>a=.5965$, $S_6>B=2.01815$.
Some $K=11$ or twelve is old-strong. Interpolate through
$(1,a),(6,B),(K,B+\sqrt K-\sqrt6)$, with zero slope after $K$.
The slopes $a,.28433,1/(\sqrt K+\sqrt6)$ strictly decrease.
Thus all duals are positive, the endpoint coefficient exceeds
$1/7$, and its energy is
\[
 q_K=.7600299945+{K-6\over(\sqrt K+\sqrt6)^2}>.91.
\]
The terminal sum $g_K$ is $<s_K<3.034$.
For any actual pressure $x.V\ge G$ with $V$ equal to one on
the first $K$ entries and $V^2\le R$, paying $1/7$ gives
\begin{equation}\label{app:tail:eq:D20weak5pay}
 \norm x^2>q_K+{2(G-g_K)\over7}-{R-K\over49}.
\end{equation}
This is one square on this reference; no head is paid.

We need a late actual endpoint. Releasing sixteen gives
$A_{15}+D_4$; the right long-leaf-three row has numerator
$\gamma+1/2$, residual $1/2$, with own-zero two-fork affine
principal, own-three triangle, and all other owns ratio
$\ge1+1/\sqrt2$. Only the two short leaves use numerator
$\gamma+1/4$. Releasing seventeen gives $A_{16}+A_3$:
middle-three has its affine star/triangle, positive ratios
$\ge1+1/\sqrt2$; end-four the weighted-star/strong rows;
end-five is a sealed actual $H_4$, strong already at own zero.
Full bridges have residual $1/32$ at sixteen, and
$1/8+1/34$ or $1/34$ at seventeen; own zero is the all-three
cycle, positive owns are strong. Left/isolated roles have the
complete local list of Lemma~\ref{app:tail:lem:whole58local}.
Its proof is local at these actual retained pieces: it does not
assert a rank condition on a larger-row theorem.

Apart from a pure left far helper, only a right short-leaf/end
three remains. Own zero, three, four, six and parallel have
respectively the affine $E_8$, cycle, $F_4$ path, six--three
path and feasible null sum. Own five has the positive left
$H_4$ projection, flat through $i-3$, $R<i$, with
$G>\sqrt{i-125/448}-.091$.
At $i=16,17$ this is $>3.745,3.9$; inserting the safe
$q>.91,g_K<3.034$ in \eqref{app:tail:eq:D20weak5pay} gives
$>1.0111,1.0349$.

For finite owns at least seven or dotted, supposed $.58$
weakness gives $.9<\gamma<.92$ at sixteen, or
$.9<\gamma<1.05$ at seventeen. This follows from the exact
target $2773138/1687945<1.643$ and the residuals $1/2,5/8$.
Use the actual core
\[
 \xi=c+\gamma^{-1}n_i+(2\gamma)^{-1}n_{i-1},\quad
 k=1-3/(4\gamma^2),\quad \alpha=\sqrt2(\gamma+1/4).
\]
Its products at $c,n_{i-1}$ vanish, at $n_i$ are negative;
its square is $-k<0$, future time positive. All other wall
products are nonpositive term by term. Its displacement is
flat through $i-2$, with
\[
 G=(A-\sqrt k)/\alpha,
 \quad A^2=i\alpha^2-(i-12)/8,
 \quad R=i-[(i-12)/8+k]/\alpha^2.
\]
At sixteen, $2.645<\alpha^2<2.738,k<.114$ give
$G>3.768,R<15.82$. At seventeen,
$k/\alpha^2$ increases through $1.05$ to
$2350/24843<.308^2$, while
$\sqrt{17-(5/8)/2.645}>4.094$.
Its derivative sign is
$-4\gamma^4-\gamma^3+6\gamma^2+(9/4)\gamma+3/16>0$
on this interval. Thus $G>3.786,R<16.82$.
Formula~\eqref{app:tail:eq:D20weak5pay} gives $>1.0213,1.0060$.
Only pure left helpers can therefore be weak at sixteen/seventeen.
Their actual mixed replacements forbid adjacency, so at least one
of those two is $.58$-strong, hence old-strong.

Append that actual endpoint $J$ to the same weak-five reference.
Its energy is at least
\[
 .7600299945+{4-\sqrt6\over4+\sqrt6}
 =.7600299945+2.2-(4/5)\sqrt6
 >1.0004379945,
\]
using $\sqrt6<244949/100000$. This contradiction proves old
strong five without any head or centre floor.
\end{proof}

\begin{lemma}\label{app:tail:lem:wholeD20}
The entire vertex forest cannot be $D_{20}$.
\end{lemma}
\begin{proof}
Keep the original $x$ and notation of the preceding lemma. There
are three head branches, denoted $H_5,H_\infty,H_f$ for native
five, dotted, and finite label at least seven. We construct their
paid references before considering the centre.

For $H_5$, put $g=\phi/2$, $A_0^2=(\phi-1)/2$.
The actual left $H_4$ projection has displacement
$g\omega_4/4$, residual $1-3g^2/2$, and exact pressure
$gA_0$. Hence $S_4\ge2\sqrt2A_0>1.572$.
A helper at any of five/six/seven is impossible: its actual
head product is positive since
$A_i^2-(g^2/A_0^2)(2\gamma+1)^2
<(28\gamma^2-32\gamma-13)/5<0$ on $.8<\gamma<1.2$.
An isolated weak dotted-five cap has actual product
$\phi\gamma-A_0\sqrt{10\gamma^2-1}\in(-.3,0)$,
where weakness gives $1<\gamma<2$.
Negativity follows from $(7\phi-11)/2>0$;
the lower bound from $\phi-\sqrt{5(\phi-1)}>-.15$.
Thus fresh $.58$ five/six/seven are genuinely strong.
On any subsequent actual-floor reference pay
$.14T$, $T=P_4-(4/5)P_5$.
Its anchor cost $.112$ is below $.151$, and its effective gain is
\[
 2(.14)(1.572-.8(1.826))-.8(.14)^2=.015456.
\]

For $H_\infty$, $u=g^{-2}\in(0,1)$ and the actual core
$c+g^{-1}n_1$ gives
$Z=(1-u/2)e_1+(u/2)e_2$ with pressure
$h=[\sqrt{2-u}-\sqrt{u(1-u)}]/\sqrt2$.
On any old reference, flat on the first five with $b<.361$,
pay $\lambda(Z-P_5/5)$, $\lambda=\sqrt2/5$.
Its cost $\lambda/5<.057$ fits the available coefficient
$>.146$. The pure scalar proved above gives gain $>.03624$;
if $u\le.5$ it gives $>.04644$.
The paid coordinates satisfy
\[
 w.P_3<1.1962,\quad w.P_2<.8918,\quad
 w_1<.5166\ (u\ge.5),\quad w_1<.5873\ (\text{all }u).
\]
These nonflat entries will be used, not replaced by a flat mean.

For $H_f$, write $g^{-2}=1+v$. The actual three-wall core is
\[
 c+g^{-1}n_1+2(g^{-1}-g)n_2;
\]
its square is $-[7-4g^2-3g^{-2}]<0$ and its displacement is
the same $Z=(1-u/2)(e_1+e_2)+(u-1)e_3$ as the finite head
in the preceding lemma. This identity uses no short normal.
Here $0<v<.234$; the elementary cosine lower bound
$\cos(\pi/7)>.9005$ gives $u<617/500$.
Pay $.4(Z-P_5/5)$. The q-free scalar
\eqref{app:tail:eq:B20headscalar}, at identical actual head
hypotheses, gives gain
$>.2893-.8(.361)+.032=.0325>.031$.
For $v\le.1$, directly
\[
 \sqrt{1-v}>.9486,\quad \sqrt{v-3v^2}<.2646
 \quad\Longrightarrow\quad
 G>.56568(.9486)-.56569(.2646)-.3368>.05.
\]
The paid coordinates obey
$w.P_3<1.243,w.P_2<.962,w_1<.461$ for $v\ge.1$,
and $w_1<.481$ for all $v$.
Only this common local head calculation is reused; the whole
$B_{20}$ theorem and its short endpoint are not inputs.

We justify the required late references. The early eleven/twelve
old reference, now with old strong five, has
$q_{11}>.8472225,q_{12}>.8689525$ and endpoint coefficient
$>1/7$. Pay the appropriate head as just constructed.
All own-five right-fork rows at fifteen/sixteen/seventeen are
excluded already without that payment: their $H_4$ pressure
is $>\sqrt{i-125/448}-.091$, flat through $i-3$, $R<i$,
giving $>1.01101,1.01087$ at fifteen, and increasing thereafter.
The additional finite/dotted rows at sixteen/seventeen have
$G>3.768,R<15.82$ or $G>3.786,R<16.82$ as proved above.
Adding even the least head gain $.015456$ to the old reference
excludes them by $>1.01616$ and $>1.00089$.
The remaining weak rows through seventeen are precisely pure
left helpers. Their actual nonadjacency gives an actual strong
endpoint $K_A\in\{14,15\}$ and $K_B\in\{16,17\}$.
For $H_5$ use fresh floors $f_i=\sqrt i-.58/\sqrt2$;
for the other heads use old $s_i$ floors. No individual
fifteen/sixteen/seventeen floor has been asserted.

Each missing knot loses
$(2\sqrt j-\sqrt{j-1}-\sqrt{j+1})^2/2
<1/[32(j-1)^3]$. In $H_5$, missing knots start at eight:
the total losses at fourteen, fifteen/sixteen and seventeen
are $<.000158,.000172,.000182$.
In the other heads they start at six; at most $4,5,5,6$
holes cost $1/3200$ each. The old rational increment list gives
$E_{14}>.9084,E_{15}>.92564,E_{16}>.94176$,
$E_{17}>.94176+1/66$.
Consequently the same legally paid references have data
\[
\begin{array}{c|rrrr|rrrr}
 &L_{14}&L_{15}&L_{16}&L_{17}
 &a_{14}&a_{15}&a_{16}&a_{17}\\\hline
H_5&.939098&.956324&.972444&.9875&3.332&3.463&3.590&3.714\\
H_\infty&.94339&.9603175&.9764375&.991276&3.311&3.442&3.569&3.692\\
H_f&.93815&.9550775&.9711975&.986036&3.311&3.442&3.569&3.692
\end{array}
\]
Here $L_K$ is a strict lower bound on twice certified pressure
minus squared norm, not on that norm itself. The actual coordinate
sum is $f_K$ or $s_K$, strictly below $a_K$.
For small dotted parameter $u\le.5$ add $.0102$ to its four
$L$ entries; for finite $v\le.1$ add $.019$.
Remaining $P_{14},P_{15}$ coefficients exceed $1/8$;
those at sixteen/seventeen exceed $.12$. The head changes
only the first-five coefficients.

If an actual $V$ is flat one through $K_A$, $x.V\ge G$,
$V^2\le R$, paying $1/8$ gives
\begin{equation}\label{app:tail:eq:D20flatpay}
 \norm x^2>L_{K_A}+(G-a_{K_A})/4-(R-K_A)/64.
\end{equation}
Both actual fork directions are $P_{19}\pm e_{20}$, of square
twenty. Cosh $>41/5$ gives $G>\sqrt{360/23}>3.955$.
Equation~\eqref{app:tail:eq:D20flatpay} exceeds one in every
head/endpoint row (the least is $1.001098$).
Thus both actual fork coshes are at most $41/5$, independently
of a centre floor. Reciprocal uses below concern actual endpoints:
for a positive finite replacement Gram, its inverse is entrywise
nonnegative. Projecting $p$ adds positive multiples of actual
normals, preserves every unselected strict inequality and makes
the released old one strict. The projected point is therefore
the actual neighboring endpoint, with no unseen truncating wall.

\paragraph{The centre inventory.}
Delete eighteen, retaining $A_{17}$ and two orthogonal forks.
If the actual centre wall has own $\gamma$, left contact
$\beta$ at $j$, and fork contacts $\delta,\epsilon$, then
\[
 \Delta=1-q_{\rm left}-\delta^2-\epsilon^2>0,
 \quad \sigma=\gamma+\beta j/18+(\delta+\epsilon)/2,
 \quad C_{18}=6\sigma/\sqrt\Delta.
\]
For simple three, $q_{\rm left}=j(18-j)/36$; endpoint four
doubles it. The complete finite-white list is: one left
endpoint-three/four at one/seventeen or next-three at two/sixteen;
one fork with a left endpoint-three and fork-three/four;
both forks with a left endpoint only in the all-three bridge;
or, without a left contact, the two-fork pairs
$(3,3),(3,4),(3,5)$. More left contacts create a cycle;
other branches would require exceptional finite rank twenty.
This is classification of the actual white set, not of a face.
The sufficient $.58$ strength target is $<83/50$.

Mixed left/fork bridges have residual $1/36$ (both forks or
fork four), with ratios $\ge19/6$ or $1/6+3/\sqrt2$.
The one-fork-three bridge has residual $5/18$: own zero is
the all-three cycle, own three the triangle; other owns give
ratio at least $(1/\sqrt2+1/36+1/4)/\sqrt{5/18}>83/50$.

With no fork contact the isolated roles have the established
$.58$ feet and affine/null exclusions. This remains the genuine
centre: the identity
$\omega_{19}+\omega_{20}=\sqrt2P_{19}$ is used only in
$n_{18}=\omega_{18}-(\omega_{17}+\omega_{19}+\omega_{20})/2$;
no extra root is introduced. For far left-four and next-three,
own zero has time zero and positive owns are strong.
Near left-four is strong already at own zero. Near next-three
own zero has the affine two-fork principal; positive owns have
ratio $\ge5/3$. Near endpoint-three owns zero/three/four/six
have respectively the two-fork, triangle, weighted-star and
six--three affine principals. All other positive owns have
ratio $(36\gamma+17)/(6\sqrt{19})>83/50$.
Far endpoint-three own zero has $A^2=-1/2$, three its cycle,
four/six the right affine principals, and parallel a null sum.
Thus the complete possible weak no-fork list is
\[
 v=\gamma\omega_{18}+\omega_1/2,\quad
 A^2=36\gamma^2+2\gamma-1/2,
 \quad \phi/2\le\gamma<6/5,
\]
with own five, finite at least seven, or dotted.

For these helpers the positive actual corrections are
\[
\begin{array}{c|c|c|c}
\text{own}&w_{\rm old}&Y\text{ in }V=P_{17}+Y&G\text{ lower},\ Y^2\text{ upper}\\\hline
5&\gamma n_{18}+(2/3)n_1+(1/3)n_2&aP_3+(e_{18}+e_{19})/2&4.185,\ .631\\
7,8&\gamma n_{18}+n_1/2&aP_2+(e_{18}+e_{19})/2&4.08,\ .655\\
\text{finite }\ge9&\gamma n_{18}&aP_1+(e_{18}+e_{19})/2&4.067,\ .785
\end{array}
\]
Here $V=(v-w_{\rm old})/(\sqrt2\gamma)$ and
$a=1/(6\gamma),1/(4\gamma),1/(2\gamma)$ respectively.
The finite five and at-least-nine rows use projected points
with residuals $2/3-\gamma^2,1-\gamma^2$; the seven/eight
row is a negative core of norm $-(\gamma^2-3/4)$.
Their common time term is
$\sqrt{18+1/\gamma-1/(4\gamma^2)}>4.33$.
The normalized residual bounds are $.145,.25,.263$;
$Y^2=ka^2+1/2$ gives the displayed variances.
All products against every unselected actual wall are nonpositive
by the positive corrections, including the actual head.

For dotted helper put $a=1/(2\gamma)\in(5/12,1/2)$ and
$t=4a^2$. The future core $c+\gamma^{-1}n_{18}$ gives
\[
 Y=aP_1+(1-t/2)e_{18}+(t/2)e_{19},\quad Y^2\le.75,
 \quad G=\sqrt{18+2a-a^2}-a\sqrt{2-8a^2}.
\]
For each constant $c\le.5873$, $(G-ca)'>7/9-c>0$.
At $a=5/12$, positive square bounds
$\sqrt{2687}>51.83$, $\sqrt{22}<4.691$ give
$G-.394a>3.829$. Hence all head/helper crosses have the bounds
\[
\begin{array}{c|rrrr}
\text{head parameter}&w.Y_5&w.Y_{7,8}&w.Y_{\ge9}&G-w.Y_{\rm dot}\\\hline
H_5&<.24625&<.219&<.211&>3.829\\
H_\infty, u\ge.5&<.2493&<.2478&<.2754&>3.7779\\
H_\infty, u\le.5&<.2493&<.2478&<.3131&>3.7484\\
H_f, v\ge.1&<.259&<.2673&<.2458&>3.8010\\
H_f, v\le.1&<.259&<.2673&<.2564&>3.7927
\end{array}
\]
The finite columns follow by multiplying the actual $w.P_k$
bounds above by $a$; the dotted column uses the increasing
$G-ca$ at the common endpoint $5/12$. No cross is set to zero.

At $K_B=17$ pay $.12$ for $Y$; at $K_B=16$ pay for
$e_{17}+Y$. The latter has an extra unit square and no extra
cross, since the reference stops at sixteen and $e_{17}.Y=0$.
One combined square gives
\begin{equation}\label{app:tail:eq:D20helperpay}
 \norm x^2>L_{K_B}+.24(G-a_{K_B}-w.Y)
                 -.0144(Y^2+\varepsilon),
 \quad\varepsilon=0\ (17),\ 1\ (16).
\end{equation}
The minimum over the four helper rows is $>1.0043$ in $H_5$,
$>1.001092$ in large-$u$ dotted, and $>1.001396$ in large-$v$
finite heads. For the small parameter cases, at $K_B=17$
the reference alone exceeds one; at sixteen the minimum is
$>1.0044935$ or $>1.0186855$ respectively. For example the
tight dotted test is
$.991276+.24(3.7779-3.692)-.0144(.75)=1.001092$.
Every no-fork weak centre role is excluded.

\paragraph{The fork rows.}
For one fork and no left contact put
$\alpha=\sqrt2(\gamma+\delta/2)$,
$A^2=18\alpha^2-(1-\delta^2)$.
At own zero, fork labels three/four/six have their actual affine
$E_8$, weighted-star and six--three principals.
Fork five has the positive left $H_4$ projection with
$G=\sqrt{4+8\phi}-\phi^{-3}>3.87$, $R=17<18$.
For finite fork labels at least seven use the positive core
\[
 c+\delta^{-1}n_f+2(\delta^{-1}-\delta)n_{18}.
\]
Its square is $-(1-b)(4b-3)/b<0$, $b=\delta^2$,
and its products at $c,n_f$ vanish and at $n_{18}$ are
negative. The displacement is flat through seventeen, with
\[
 R=28-16/b+6/b^2,\qquad
 G=\sqrt{20-2/b}-\sqrt{-8+14/b-6/b^2}.
\]
Every finite label is in the following table:
\[
\begin{array}{c|cc}
\text{fork label/domain}&G\text{ lower}&R\text{ upper}\\\hline
5&3.87&18\\
7, .81<b<.82&3.798&17.412\\
8&3.79392&17.49034\\
\ge9, .88<b\le.9&3.808&17.63\\
\ge9, .9\le b<1&3.831&18
\end{array}
\]
For eight the exact correlated values are
$R=108-64\sqrt2$ and
$G=\sqrt{12+4\sqrt2}-\sqrt{68\sqrt2-96}$.
The rational seal $1.414213562<\sqrt2<1.414213563$
gives the table. In the other rows $R$ increases for $b>.75$,
and the second radicand increases up to $6/7$ and decreases
thereafter, proving the interval bounds. Labels seven and at
least nine satisfy the indicated disjoint native ranges.
Equation~\eqref{app:tail:eq:D20flatpay} excludes every row
with either actual early endpoint and each head: the least
test is the native-five/eight/fourteen value
$.939098+(3.79392-3.332)/4-3.49034/64=1.0000414375$.
For finite heads its least test is $>1.0041459$; for dotted
heads $>1.0093859$. No head cross occurs because these
displacements are flat throughout the head support.

For positive own three, release the contacted fork. The same
actual centre wall completes the retained $A_{19}$ to finite
$D_{20}$, so the reciprocal endpoint has
$C_f=9+10\delta\ge14>41/5$, impossible.
Own four has the affine star on $c,18,17$ and the other fork.
Own six/fork three has the six--three principal.
Own six/fork at least four and own at least seven/fork at
least four are directly strong by ratio
$\ge\sqrt2\gamma+1/2>83/50$; own five/fork at least five
is also directly strong.
The two exceptions five/three and five/four have the positive
left $H_4$ correction $\gamma(3n_{18}/2+n_{17}+n_{16}/2)$,
finite residual $(5-3\phi)/8$ and
$\alpha^2\ge\phi+5/8$.
Their displacement is flat through fifteen, $R<18$, and
$G>\sqrt{18-75/224}-.091>4.1$, dominating the common row.
For single-fork three with finite at least seven or dotted
own, $\gamma\ge6/5$ is directly strong. Throughout
$.9<\gamma<6/5$ the actual core
$c+\gamma^{-1}n_{18}+(2\gamma)^{-1}n_{17}$ has
$k=1-3/(4\gamma^2)>0$, flat displacement through sixteen,
$R<18$, $A/\alpha>4.208$ and $\sqrt k/\alpha<.338$.
Thus $G>3.87$ and the same common row excludes it.
Parallel own is independently null; including it in that
continuous core interval would also contradict the reference.

For two forks the complete pairs are three/three, three/four,
three/five. Own three has the affine triangle, own at least
four ratio $\ge1+1/\sqrt2$, and parallel the null sum.
At own zero, three/three is the affine square.
Three/four has, after releasing the three-contacted fork, the
actual finite $B_{20}$ reciprocal endpoint
$C_f=9+5\sqrt2>41/5$, impossible.
Three/five has a positive finite $H_3$ projection, displacement
$\alpha P_{18}$, $\alpha=\phi^2/(2\sqrt2)$ and
$q=2\phi^{-6}<1/8$; its pressure is
$\sqrt{18-q}/(1+\sqrt{q/18})>4.22/(13/12)>3.87$, $R=18$.
It too is covered by the common row. All original wall
inequalities in these finite projections follow from the
positive inverse Gram coefficients.

The complete centre domain has been exhausted. Eighteen is
therefore genuinely $.58$-strong in every head branch.
Append its fresh floor for $H_5$, its old floor otherwise,
to the reference whose actual endpoint is $K_B$.
At seventeen the added energy is $d_{17}^2>1/70$;
at sixteen the chord across the possibly weak seventeenth
adds $17-12\sqrt2>.0294$. Both slopes are below the available
endpoint coefficient and leave the head payment unchanged.
The least resulting totals are the finite-head values
$.986036+1/70>1.0003217$ and
$.9711975+.0294=1.0005975$.
This contradicts the same original $\norm x<1$ in every branch.
\end{proof}

\subsection{The complete dimension-twenty conjunction}
\begin{proposition}\label{app:tail:prop:exact20}
There is no compact hyperbolic Coxeter polytope of dimension twenty.
\end{proposition}
\begin{proof}
We exhaust the forest of an actual farthest vertex in this dimension;
no descent from a higher-dimensional face is involved.
An irreducible $A_r$ with $r\ge7$, or $B_r,D_r$ with $r\ge9$,
is called dominant. Each costs more than one half, so at most
one occurs. If none occurs, the ordinary rank-fourteen rate,
even discarding the single possible isolated $I$ rank, gives
$19/14>1$. A typed rank-two defect instead gives
$.09+18/14>1$; the funded $J$ pair is stronger than its ordinary
rank rate. Thus exactly one dominant occurs.

For an $A$ dominant, rank seven is excluded by its half bound
and the ordinary spectator rate. Ranks eight through twelve
are excluded by $4/7$ and that same rate: even with $I$
there are seven ordinary ranks at twelve, or with a typed
defect six ranks beside its $.09$ cost.
Ranks thirteen/fourteen have energy $>5/7$ and at least five
ordinary ranks after discarding $I$; a typed defect leaves at
least four, so both alternatives exceed one.
At ranks fifteen, sixteen and seventeen use respectively
$>.826,>.86,>6/7$.
The complete minimum spectator comparisons are
\[
\begin{array}{c|c|c|c}
r&\text{ordinary or }J& I\text{ present}&\text{typed present}\\\hline
15&>5/14&>2/7&>.09+17/70\\
16&>2/7&>2/7&>.09+1/7\\
17&>17/70&>.21&>.09+.1
\end{array}
\]
Each sum with its actual dominant floor exceeds one.
For the middle $I$ entry the ordinary remainder has rank three;
the joint $I$ theorem, rather than discarding that rank, supplies
$>2/7$. The rank-two typed remainders are either an ordinary
dihedral or two ordinary $A_1$ factors, both at least the shown
charge. There is at most one defect.
Rank eighteen uses $>.91$: an ordinary rank-two remainder is
$>1/7$, $I+A_1>1/5$, either typed defect $>.09$, and $J>.34$.
The typed sum is strictly greater than one as well.
At rank nineteen its remaining $A_1$ is either ordinary or $I$;
Lemmas~\ref{app:tail:lem:A19ordinary} and
\ref{app:tail:lem:A19I} exclude both. Rank twenty is whole and
is excluded by Lemma~\ref{app:tail:lem:wholeA20}.

For a $B/D$ dominant its $9/14$ bound forces rank at least
fifteen: at rank fourteen even an $I$ leaves five ordinary
ranks, while a typed defect leaves four beside its $.09$.
At fifteen and sixteen the valid weaker $D$ prefix floors
are $>.797425,>.823735$. The smallest remainders are
$I+$ordinary rank four $>2/7$, respectively a typed defect
plus an ordinary dihedral $>.09+1/7$; all other partitions
are stronger. Both rows exceed one.
At seventeen the weaker floor is $>.8472225$.
Ordinary total rank three costs $>17/70$, $I+$dihedral costs
$>.21$, $I+$two ordinary $A_1$ costs $>.3$, and a typed
defect plus ordinary $A_1$ costs $>.19$. Each closes the row.
At eighteen the weaker floor $>.8689525$ excludes ordinary
rank two, $I+A_1$, and $J$. Both typed remainders are excluded
by Lemma~\ref{app:tail:lem:proper20classical}.
At nineteen a nonshort $A_1$ is excluded by the actual complete
coverage through fifteen and its prefix gain $>.91211125$;
that same lemma supplies the needed local $D_4$ continuation
before using the prefix proof. Its $I$ alternative is also
excluded there. Whole $B_{20}$ and $D_{20}$ are excluded by
Lemmas~\ref{app:tail:lem:wholeB20} and
\ref{app:tail:lem:wholeD20}.
Finite classification and the complete spectator partitions above
exhaust the actual rank-twenty forest. Orthogonal component
energies use the same original $x$ throughout.
\end{proof}

\subsection{The two proper $A$ rows in exact dimension nineteen}
\begin{lemma}\label{app:tail:lem:A17typed}
An $A_{17}$ component cannot coexist with either actual typed-seven
defect as its entire rank-two remainder.
\end{lemma}
\begin{proof}
The actual defect costs $>.09$, so $e_A<.91$. Use the complete
$N=18$ window of Appendix~\ref{app:Acentralwindow}, before its
weaker contrary-energy conclusion. A near-four actual mate has
normalized demand $>.86$; the opposite strong anchor or mate
has demand $>.77$ and correlation $<.44$.
Cauchy would give $1.63^2<.91(2.88)$, false.
All weak knots five through thirteen are thus nonadjacent pure
helpers, with $Q=0$, own five, finite at least seven or dotted,
$\phi/2\le\gamma<17/16$.
The rational reference of Appendix~\ref{app:Aendpointcores}
has end heights $H_5=2.077,H_6=2.218$ and central heights
$2.314,2.371,2.389$, reflected. For zero, one or two weak end
sites its complete all-hole energies are
\[
 E_0\ge.8934308,\quad E_1>.86256,\quad
 E_2\ge9978827/12000000.
\]
Strong ends have five-block mean $.29<b_0<.294$;
weak ends six-block mean $1/4<b_1<.2614$.

Both actual defects force the endpoint floors
$x_1>.392,-x_{18}>.392$.
For old $A_2$/new seven, propagation gives endpoint squared
half-tangent $>1/5$, hence $x_1^2>17/90$.
For old $I_2(7)$/new three let its isolated opposite have
time square $h=(1/4)\csc^2(\pi/7)-1$.
The actual endpoint opposite retains that whole old block.
Any finite contact to it would cost at least
$1/[4\sin^2(\pi/7)]>1$, so the endpoint cap has no such contact.
Releasing old one leaves $A_{16}$. Its complete nonisolated
contact list is three at $2,3,16,17$, or four at $2,17$.
The corresponding coshes and minimum allowable owns are
\[
\begin{array}{c|c|c}
3:2&(17\gamma+8)/9&1/2\\
3:3&(34\gamma+15)/(6\sqrt2)&0\\
3:16&(17\gamma+1)/(3\sqrt2)&1/2\\
3:17&(34\gamma+1)/18&1/\sqrt2\\
4:2&(17\gamma+8\sqrt2)/3&0\\
4:17&(17\gamma+1/\sqrt2)/3&1/2
\end{array}
\]
All but the far-three row have cosh $>3/2$, giving
$x_1^2>17/90$. The remaining row gives
$289(\sqrt2-1)/[18(17\sqrt2+19)]>(49/125)^2$;
clearing denominators leaves
$3780919\sqrt2>5336767$, certified by $\sqrt2>707/500$.

For an isolated endpoint $A^2=(17/9)\gamma^2-1$ and its
product with the actual defect wall is $-A\sqrt h$.
For finite own, $hA^2<8/27<1/2$; the native gap forces
$hA^2=1/4$. The literal seven algebra gives
$7h^3+7h^2-1=0$, $.327<h<.329$, with degree three over
$\mathbb Q$. Consequently $.93<\gamma^2<.94$.
Only label twelve can lie in this interval: for labels at most
eleven, $\sin(\pi/11)>717/2662$ gives squared cosine $<.93$;
for labels at least thirteen, $\pi/13<22/91$ gives squared
cosine $>7797/8281>.94$.
At twelve equality would imply $h=9/(17\sqrt3-2)$, impossible
for its cubic degree. This is literal-angle algebra, not an
arithmetic ground-field restriction. Parallel own is excluded
by the same gap and interval. For dotted own $hA^2>1/4$;
the next native possibility gives $A^2\ge1/(2h)>3/2$,
so its cosh squared exceeds $5/2$ and again supplies the floor.
Reflection covers the other endpoint.

Every weak cap gives a stronger floor at its helper end.
The complete raw-pressure seals in Appendix~\ref{app:Arawends}
are $.36082,.327322,.3132358,.31446$ for five, seven, eight
and larger finite owns. For dotted own its exact weak floor
gives $>.310125+1/216>.3132358$.
Thus the helper endpoint exceeds $.4429$, since
$1.4142(.3132358)=.44297806836$.
A near helper has larger time square by $8\gamma/9$.
As the two times are $<3,4$, its time increases by
$>32/315>.1$, with the same residual. Its endpoint exceeds
$.58$ for every own parameter.

On a constant block of length $m$ and mean $b$, grafting a
signed endpoint value $z>b$ gains $m(z-b)^2/(m-1)$.
All the following grafts preserve positive ordered slopes:
the other coordinates stay $>.2517$ in a five-block or
$>.184$ in a six-block, above the next reference coordinate.
Opposite end blocks are disjoint, so they are changed in one
reference, including all previous hole charges.
With no weak ends use $.392$ twice:
$E_A>.8934308+2(5/4)(.392-.294)^2=.9174408$.
With one weak end a near helper alone gives $>.98$.
Otherwise graft $.4429$ in the opposite strong block and
$.392$ in its weak block, giving
$>.86256+(5/4)(.4429-.294)^2+(6/5)(.392-.2614)^2
=.9107416445$.
With two weak ends a near helper alone gives $>.95$;
otherwise their far helper endpoints are distinct, giving
$>9978827/12000000+2(6/5)(.4429-.2614)^2>.91063$.
All cases exceed $.91$, contradicting the orthogonal defect cost.
\end{proof}

\begin{lemma}\label{app:tail:lem:A18I}
An actual $A_{18}$ component cannot coexist with an isolated $I$;
further orthogonal finite spectators are allowed.
\end{lemma}
\begin{proof}
Use $N=19$ and the one actual $I$ coordinate $y\ge s$.
Its rank-free parameters give both endpoints at least $\eta t$,
where $\eta=\sqrt{18/19}>.973$ and $t+(3/2)s>.767$.
We first strengthen the near-four removal to the norm-only
domain needed here. Its actual finite mate, replacing $j=i-1$
at $i=5,6$, has diagram $B_{20-i}+A_{i-2}+S$ and cosh
\[
 \sqrt{2(i-1)/19}\,[(20-i)/2+(19-i)/\sqrt2]>11.
\]
Finite positive projection retains every spectator and recognizes
this actual endpoint. Its normalized height exceeds $.91$.
The opposite mate or strong knot thirteen/fourteen has height
$>.774$ and normalized correlation $<.41$.
Thus $1.684^2=2.835856>2.82$ contradicts $e_A<1$.
The complete shared $N=19$ window now has only nonadjacent pure
helpers, before any other energy assumption.

The shared rational heights are
$H_5=2.104,H_6=2.255,H_7=2.363,H_8=2.433,H_9=H_{10}=2.467$,
reflected. With zero, one or two weak end sites the complete
energies are $E_0\ge.9240352,E_1>.8937,E_2>.86414$.
The strong five-block mean is $b_0<.298$, weak six-block
mean $b_1<.266$, each $>.29$ or $>1/4$ respectively.
Available dual coefficients are
$\beta_6\ge1607/12000$ and $\beta_{14}\ge.1349$;
deleting the next internal knot only increases them.

A weak-five far helper has
$A^2=(140\gamma^2+10\gamma-10)/19$.
Use the full positive actual feet of Appendix~\ref{app:Aendpointcores},
not their raw normal sums. Their corrections are
$\gamma n_5+(2/3)n_e+(1/3)n_{e'}$ for five,
$\gamma n_5+2(1-\gamma^2)n_e$ for finite owns at least seven,
and $\gamma^{-1}n_5$ for dotted own.
Their actual pressures are bounded below by $A$ minus the
respective finite/core radical; the full $I$ block remains
an unselected acute block.
When both ends are weak, projecting the spatial foot gives
\[
 PU=(5\gamma/6)\omega_6+\omega_{13}/12,
 \quad Y=U-PU,\quad w.Y=0.
\]
Its variance $V_2$ is
$(2/3)\gamma^2+1/12$ for five,
$4\gamma^4-(16/3)\gamma^2+29/12$ for other finite owns,
and $(5/3)\gamma^2+\gamma^{-2}-19/12$ for dotted own.
The exact cross between two far feet is
$Y_L.Y_R=(\gamma_L+\gamma_R)/6$: each head residual
pairs with the opposite helper's mass $1/\sqrt2$, and all
same-side block sums vanish.

Pay $\mu=.138$ for each far foot. At either anchor the total
cost is $<.138(31/32)<1607/12000$.
For $\Delta_2=H-2.255(5\gamma/6+1/12)$ the gain per foot,
including half the exact cross, is
$F=2\mu\Delta_2-\mu^2(V_2+\gamma/3)$.
The full finite/dotted seals are
\[
\begin{array}{c|c|c|c}
\text{own}&\Delta_2\text{ lower}&V_2+\gamma/3\text{ upper}&F\text{ lower}\\\hline
5&67/192&.791&.081248696\\
7&.3371541666&1.026&.0735154059\\
8&.32852125&1.087&.069971037\\
\ge9&1591/4800&1.107&.070400792\\
\text{dotted}&.36109375&1.54&.070334115
\end{array}
\]
The first three use actual time/radical bounds
$(2.17,.111),(2.434,.2158),(2.499,.2464)$ and upper native
cosines $.81,.901,.9239$. For larger finite owns evaluate
at $.93$, where $A>2.517$, then use $A'>2.714$ and
$V_2'\le16/3$ to see that the lower gain increases throughout
the remaining finite interval. For dotted owns use
$1<\gamma<17/16$; its deficit decreases to an endpoint
$>.36109375$, while $V_2<1.185,\gamma/3<.355$.
This proves every row continuously, not by sampling.
Both far weak ends therefore give
$E_A>.86414+2(.0698)>1.00374$, without spending $I$.

For exactly one weak end, a near helper already gives
$E_A>1.0120152$ by the shared raw endpoint floor $.58$.
For a far helper the projection becomes
$(5\gamma/6)\omega_6+\omega_{14}/10$.
Its deficit and variance are
$\Delta_1=\Delta_2-1349/60000$, $V_1=V_2-1/60$.
Payment $\mu=.11$ gives a gain $>.058$ in every row:
the conservative five/seven/eight/larger/dotted gains are
$.065714,.06064868,.05809604,.0585244,.06034939$.
The same finite derivative and dotted endpoint bounds prove
the complete intervals.
Retain the propagated left endpoint at the same time.
Let $R=e_1-P_6/6$, so $R^2=5/6$ and
$Y.R=\gamma\sqrt2/6$.
Put $B=b_1+\mu\gamma\sqrt2/6<.294$,
$z=(\eta t-B)_+$ and $\nu=\min((6/5)z,1/4)$.
The full vector is $w+\mu Y+\nu R$.
Its coefficient-six cost is
$\mu5\gamma/6+\nu/(6\sqrt2)<.128<1607/12000$;
the opposite cost is $.011<.1349$.
Its full mixed square therefore gives
\[
 E_A+y^2>E_1+.058+s^2+2\nu z-(5/6)\nu^2.
\]
If not clipped, $z+(3\eta/2)s>.452$ and weighted Cauchy
gives $s^2+(6/5)z^2>.452^2(57/169)>.068$.
If clipped the graft alone is $>5/96$.
Both totals exceed $.8937+.058+5/96>1$.

With no weak ends graft both propagated endpoint floors in
the two five-blocks. Their gain is $(5/2)z^2$, with
$z=(\eta t-b_0)_+$ and
$z+(3\eta/2)s>.448$.
Its sum with $s^2$ is $>.448^2(190/481)>.0792$,
contradicting $E_0$. The grafts stay ordered even for
$\eta t<1$, because their other block entries exceed $.1125$
and the next reference entry is $<.108$.
Finally if both ends are weak with a near helper, its $.58$
graft gives $E_A>.9824552$. If the opposite propagated floor
exceeds $.58$ the other $.58$ graft already contradicts the norm.
Otherwise graft its actual propagated floor there; with
$z=(\eta t-b_1)_+$ one has
$z+(3\eta/2)s>.480$, so its gain with $s^2$ is
$>.480^2(57/169)>.077$.
Both block changes are legal and disjoint. All end patterns,
internal weak sets and actual own/foreign possibilities are
covered, spending the one $I$ energy only once.
\end{proof}

\subsection{The proper classical rows in exact dimension nineteen}
\begin{lemma}\label{app:tail:lem:B18I}
The entire forest $B_{18}+I$ is impossible.
\end{lemma}
\begin{proof}
Use the actual short root $n_{18}=e_{18}$ and long roots before it.
The genuine generic bulk covers five through fourteen.
At fifteen retained $B_3$ admits its long-three and short-three
$F_4$ extensions. The latter has residual $1/4$, numerator
$\gamma+1/(2\sqrt2)$, own-zero affine $F_4$, and positive-own
ratio $\ge1+1/\sqrt2$. Its full white $F_4$ cannot acquire
an $I$ arm; full $B_{18}$ bridges have internal new node and
are also foreign-sealed. Left/isolated/foreign estimates use
only genuine long roots and retain their full old-$I$ clearing.

At sixteen the right piece is $B_2$. A second four is not a
finite extension, so the only new contacts are a three at either
end. Long-three has numerator $\gamma+1/2$, residual
$1/2-Q$, affine own zero/three and otherwise strong ratio.
Short-three has numerator $\gamma+1/(2\sqrt2)$, same residual;
owns zero/four/six have the affine $F_4$/star/six--three
principals, parallel the actual null sum. Own at least five
apart from affine six is strong; at own three a positive
$Q\ge1/4$ is strong. Thus the sole extra row is
$Q=0,v=(\omega_{16}+\omega_{18})/2$.
Full endpoint bridges have internal new node, no foreign
extension, and direct strong ratio already at own zero.

For the extra row release old seventeen. The same actual wall
completes finite $A_{18}$, with the old $I$ retained.
Positive finite projection identifies the actual endpoint and
gives $C_{17}=(17+16\sqrt2)/\sqrt{19}>9.09$, hence
$S_{17}^2>13753/1009>3.6919^2$, stronger than $f_{17}$.
If five is weak the actual interpolation through one, six and seventeen already
has energy $>1.0004379945$.
If five is strong, extend its actual fourteen/fifteen reference
by the chord to $f_{17}$. The energy is $>.9550775$, the first
five entries remain $b<.361$ and their coefficient $>.146$.
The single actual-$I$ initial graft adds $>.0496$: for a
nonclipped five-block, weighted Cauchy with
$t+(3/2)s>.767$ gives
$s^2+(5/4)(t-b)_+^2>(.767-.361)^2/(61/20)>.0496$;
clipping the endpoint at $.6$ already gives $>.0714$.
This rules out the extra row before using sixteen in a bootstrap.

Coverage through sixteen is now precisely nonadjacent pure left
helpers. If five is weak, one of fifteen/sixteen is strong.
At sixteen its weak-five reference alone exceeds one; at fifteen
it has effective energy $>.9851486$, first coordinate $.5965$
and remaining $P_6$ coefficient $>1/60$.
Lemma~\ref{lem:I-capped-initial-graft} adds $>.015$ using
the unused actual $I$. Thus five is genuinely strong.

Choose the actual strong endpoint fifteen or sixteen.
The ordinary old-floor reference has $q>.9240775$, first
five entries $b=f_5/5<.361$, individual demand $S_5=f_5$,
and available $P_5$ coefficient $>.146$.
We now check the local head lemma's actual whitelist, rather
than importing a $B_{17}$ row theorem. Deleting one retains
$B_{17}+I$. Its only nonisolated finite extension is simple
three at old two, with optional $I$ contact only three.
The isolated-within-$B$ row may meet $I$ by any finite label.
These are exactly the two raw equations
\[
 A^2=2\gamma^2+Q-1,\qquad
 A^2=2\gamma^2+2\gamma+Q,
\]
with $Q$ arbitrary finite in the first row and $Q=0,1/4$ in
the second. All head points use only genuine long normals
one/two/three; old short root eighteen remains an unselected
acute wall. The whole head, including positive foreign contact,
therefore satisfies Lemma~\ref{lem:I-head-flat-reference}.
Its q-free gain $>.07665$, with $I$ once, gives
$q+\text{gain}>1.0007275$.
There is no need to prove sixteen individually strong or to
repeat a terminal head/I scalar catalogue.
\end{proof}

\begin{lemma}\label{app:tail:lem:D17typed}
The entire forest $D_{17}$ with either actual typed-seven defect
is impossible.
\end{lemma}
\begin{proof}
Use the genuine signed forks sixteen/seventeen, centre fifteen,
and the one defect projection of energy $e_7>.09$.
For old $I_2(7)$, white positivity prohibits any foreign contact,
since its smallest cost is $1/[4\sin^2(\pi/7)]>1$.
For old $A_2$ both old nodes cannot be contacted (the triangle
is affine); a single contact gives precisely
$Q\in\{0,1/3,2/3,\phi^2/3\}$.
In particular every positive foreign cost is at least $1/3$.
This is not the $I$ cost $1/4$.

The bulk covers five through twelve, with only nonadjacent pure
helpers. Choose actual strong $K=10$ or eleven, and separately
nine or ten. In both initial-five branches a positive reference
has endpoint no greater than $f_K$, energy
\[
 q_9>.797425,\ q_{10}>.823735,\ q_{11}>.8472225,
 \quad f_9<2.569,f_{10}<2.732,f_{11}<2.886.
\]
The endpoint price exceeds $.16$ at nine/ten or $1/7$ at
ten/eleven. For every actual pure pressure $x.V>G$ flat one
through $K$, $V^2\le R$, its one-square payment gives
\begin{equation}\label{app:tail:eq:D17typedpay}
 \norm x^2+e_7>q_K+.09+2\mu(G-f_K)-\mu^2(R-K).
\end{equation}
For $K=10,11,\mu=1/7$, it suffices that $14G-R>32.6$;
the two resulting margins above one exceed
$25003/9800000,47961/19600000$.

At thirteen retained $D_4$ permits three at a leaf.
The long leaf fourteen uses numerator $\gamma+1/2$, not
the fork numerator $\gamma+1/4$. Own zero/three is affine,
otherwise the long-leaf row is strong. In the fork row
owns zero/three/four/six and parallel have their actual
$E_8$/cycle/weighted-star/six--three/null obstructions;
positive $Q$ and pure finite at least seven or dotted are
strong. Its only extra is pure own five.
The left $H_4$ correction on thirteen/twelve/eleven gives
$\alpha^2=\phi+5/8$, $A^2=13\alpha^2-1/2$, flat
displacement through ten, $R<13$, $G>3.48$.
Pay $.16$ at actual nine/ten; the totals in
\eqref{app:tail:eq:D17typedpay} exceed $1.076545,1.076295$.
Full bridges are finite $D_{17}$ with internal new node,
so cannot acquire a finite foreign arm; residual $1/26$
and the own-zero cycle cover them.

At fourteen right $A_3$ adds middle-three/end-three/four/five.
Middle-three/end-four have their affine or direct strong rows,
end-five is a sealed actual $H_4$ strong already at own zero.
Full bridges are finite $D_{17}$ or $A_{17}$, with the new
node too far from the ends to add an $A_2$ arm.
The lone end-three has residual $5/8-Q$, numerator
$\gamma+1/4$. Positive $Q$ is strong; zero/three/four/six
and parallel have the actual affine/null obstructions.
Pure dotted owns are strong. Own five has the left $H_4$
foot $G>3.61,R<14$, flat through eleven, hence $14G-R>32.6$.
For the remaining $.9<\gamma<1$ select both actual $A_2$
cores with old corrections
$\gamma(4n_{14}/3+2n_j/3)$, $j=13,15$.
Their averaged pressure is pure, flat through twelve, with
\[
 \alpha=\sqrt2(\gamma+1/4),\quad
 G>\sqrt{14-125/529}-.327>3.38,\quad
 R=14-{14\gamma^2/9-3/8\over2(\gamma+1/4)^2}<13.67.
\]
The subtracted quotient increases from $177/529> .33$;
the radical quotient increases to $8/75<.327^2$.
Thus $14G-R>33.65$ on the complete interval.
All cores have positive coefficients, negative square, future
time and nonpositive products at every original wall, including
the two defect walls. Coverage now reaches fourteen.

Delete centre fifteen. Retained whites are $A_{14}$ plus
the actual two forks and the defect. With left contact $\eta$
at $j$, forks $\delta,\epsilon$ and total foreign cost $Q$,
\[
 \Delta=1-q_{\rm left}-\delta^2-\epsilon^2-Q>0,
 \quad\sigma=\gamma+\eta j/15+(\delta+\epsilon)/2.
\]
The sufficient $.61$ ratio target is $<25/16$.
The complete old-piece list is left endpoint-three/four,
next-three, neither; left plus one fork only endpoint-three
with fork-three/four; left plus both forks only the all-three
bridge; without left contact, one fork or the pairs
$(3,3),(3,4),(3,5)$. Additional cycles, long arms and interior
four bonds are not finite. Full bridges are foreign-sealed;
the one-fork-three bridge cannot acquire a foreign $A_2$ arm
of lengths fourteen/one/two. Their residuals $1/30,17/60$
give the affine or direct strong rows.
In the no-fork list zero-own far-three would need
$1/2<Q<8/15$, absent from the exact cost list;
far-four/next-three require $Q<1/15,2/15$, so have time zero.
The other zero-own rows are actual affine principals.
The isolated and foreign-cleared feet use genuine old
fifteen/fourteen/thirteen, and only the pure left helper remains
weak; the sum of the two fork duals is solely a spatial identity.

The complete hard no-left centre rows have the following
pure displacements, all flat through at least eleven:
\[
\begin{array}{c|c|cc}
\text{row}&\text{actual construction}&G\text{ lower}&R\text{ upper}\\\hline
\text{own }0,\text{ fork }5&\text{left }H_4&3.49&14\\
\text{own }0,\text{ fork }\ge7&\text{three-wall core}&3.4&15\\
\text{own }3,\text{ fork }4,5&\text{actual reciprocal }D_{17}&3.81&17\\
\text{own }5,\text{ fork }3&\text{left }H_4&3.73&15\\
\text{own }\ge7/\text{dot},\text{ fork }3&\text{special core}&3.49&15\\
\text{own }0,\text{ forks }3,4&\text{actual reciprocal }B_{17}&3.81&17\\
\text{own }0,\text{ forks }3,5&\text{finite }H_3&3.53&15
\end{array}
\]
Every row has $Q=0$: a positive cost makes the five/own-three
rows strong and is prohibited by white positivity in the others.
For the fork-five row, its normalized time and finite residual
are $1+8\phi,13-8\phi$, giving $G>3.49$.
For fork at least seven let $z=1/\delta^2\in(1,100/81)$.
The positive core is
$c+\delta^{-1}n_f+2(\delta^{-1}-\delta)n_{15}$, with
\[
 G=\sqrt{17-2z}-\sqrt{-8+14z-6z^2}>3.81-.41=3.4,
 \quad R=25-16z+6z^2<15.
\]
The second radicand is $1/6-6(z-7/6)^2$.
For own three and fork four/five the same actual centre
wall completes a finite $D_{17}$ endpoint with cosh
$15/2+(17/2)\delta>13$; positive inverse projection retains
the defect and recognizes that endpoint.
For the three/four pair the finite mate is $B_{17}$ with
cosh $15/2+(17\sqrt2)/4>13$.
Both give the stated $G>\sqrt{17(6/7)}>3.81$.
Own-five/fork-three has the left $H_4$ correction and
$G>\sqrt{15-75/224}-.091>3.73$.
For the remaining $.9<\gamma<6/5$ single-fork-three interval
use $c+\gamma^{-1}n_{15}+(2\gamma)^{-1}n_{14}$,
of norm $-[1-3/(4\gamma^2)]$.
Its normalized radical is $<.34$ and time
$>\sqrt{15-150/529}$, giving $G>3.49$.
Larger owns are directly strong. The three/five pair uses the
positive $H_3$ projection with $q=2\phi^{-6}<1/8$ and
pressure $\sqrt{15-q}/(1+\sqrt{q/15})>3.53$.
Own zero fork-three/four/six and the two three forks have
their affine principals; positive own four/six and parallel
their star/path/null exclusions. Remaining positive-own/fork
at least four rows have direct ratio $\ge\phi/\sqrt2+1/2$.
Thus every own/finite-white/foreign case is either listed,
strong, or affine/null impossible.
Each listed row satisfies $14G-R>32.6$, so
\eqref{app:tail:eq:D17typedpay} excludes it.

Only nonadjacent pure helpers remain through fifteen.
Weak five would give a strong fourteen/fifteen endpoint and
energy $>.968$ from its actual first-six reference; adding
the one $.09$ defect contradicts the norm. Thus five is strong.
If fifteen is strong its reference $q>.9240775$ suffices.
Otherwise fourteen is strong, $q>.90715$.
The actual weak-fifteen feet use only $n_{15},n_1,n_2$ and
positive corrections, with both fork products separately
negative; the spatial sum $\omega_{16}+\omega_{17}$ is not
an added wall. More explicitly, put
$\Omega_{16}=\omega_{16}+\omega_{17}$ solely as a spatial sum.
Their displacements are
$\gamma(\omega_{14}+\Omega_{16})/2+\omega_3/6$ for five,
the same with $\omega_2/4$ for seven/eight, or $\omega_1/2$
for finite owns at least nine. For dotted own the displacement is
$(\gamma-\gamma^{-1})\omega_{15}
+(\omega_{14}+\Omega_{16})/(2\gamma)+\omega_1/2$.
Their respective pressures are
$H(A,2/3-\gamma^2)$, $A-\sqrt{\gamma^2-3/4}$,
$H(A,1-\gamma^2)$, $A-\sqrt{1-\gamma^{-2}}$,
with $A^2=30\gamma^2+2\gamma-1/2$.
Concavity of the original heights bounds the dual sum by
$2h_{15}$ and $h_2\le2h_1,h_3\le3h_1$.
The full native/dotted scalar in the proof of
Lemma~\ref{app:tail:lem:prefixgain}, valid for
$\sqrt{10}\le q\le\sqrt{42}$, therefore gives
$x_1>\sqrt2[T-\gamma(\sqrt{30}-.61)]>.424$.
Grafting only the first five slopes gains
$(5/4)(.424-.361)^2=3969/800000$.
The same reference and one defect therefore give
$\norm x^2+e_7>.90715+3969/800000+.09
=1+1689/800000>1$. Both typed defects are excluded.
\end{proof}

\begin{lemma}\label{app:tail:lem:D18I}
The entire forest $D_{18}+I$ is impossible.
\end{lemma}
\begin{proof}
Use the actual centre sixteen and signed forks seventeen/eighteen;
$x_1>\cdots>x_{17}>|x_{18}|$, with no assumed sign of $x_{18}$.
Let $s,t$ be the same rank-free actual-$I$ parameters.
We will need the slightly stronger elementary barrier
\begin{equation}\label{app:tail:eq:I788}
 t+(8/5)s>.788.
\end{equation}
To prove it, substitute $t_0=.788-1.6s>0$ into
$(1-s^2)(1-t^2)-8st$.
At $z=s-1/5\in(-1/5,2/15)$ its value is
\[
 .00093696-.0586944z+8.962384z^2-.4736z^3+2.56z^4
 >.0009-.06|z|+8z^2>0.
\]
The cubic is bounded below by $-.064z^2$ and the last
quadratic has negative discriminant. Since the original
expression decreases with positive $t$, the barrier follows.
On any strong-five reference with $b<.361$, a legal first-five
graft and the one $I$ energy consequently add
\[
 s^2+(5/4)(t-b)_+^2>
 (.788-.361)^2/[4/5+(8/5)^2]>.05426.
\]
If $t\ge.6$, clip that graft at $.6$; it alone gains $>.0714$.
All other block entries stay above the next slope, so this is
an actual positive-dual payment, not an independent norm estimate.

The genuine bulk is complete through thirteen. Choose an actual
strong $K=10$ or eleven before classifying either tail.
In the strong-five branch the sharper ordinary references have
$q_{10}>.8244,q_{11}>.8482$, endpoints
$f_{10}<2.731,f_{11}<2.8854$, and available price $>1/7$.
Indeed their unomitted energies exceed $.824619,.848442$;
nonadjacent holes at six/eight/ten lose respectively at most
half the squares of $.0172,.01111,.008$.
The total losses are $<.000209637,.000241637$.
In the weak-five branch the actual first/six references give
$L_{10}>.8870,L_{11}>.9104$, without any $I$ charge.
For every actual pure pressure $x.V\ge G$, flat one through
$K$, $V^2\le R$, the two branches therefore give
\begin{equation}\label{app:tail:eq:D18Ipurepay}
 \norm x^2+y^2>
 \begin{cases}q_K+.05426,&5\text{ strong},\\L_K,&5\text{ weak},\end{cases}
 \quad +{2(G-f_K)\over7}-{R-K\over49}.
\end{equation}
The pure residual vanishes on the initial support, so its
cross with the first-five/$I$ adjustment is zero.

At fourteen retained $D_4$ admits three at a leaf. Distinguish
its long-leaf numerator $\gamma+1/2$ from short-leaf
$\gamma+1/4$. The former is affine or strong; in the latter
positive foreign cost, finite at least seven and dotted owns
are directly strong. Zero/three/four/six and parallel have
their actual affine/null obstructions. The remaining pure
five has a positive left $H_4$ foot, flat through eleven,
$G>3.62,R<14$, excluded in both branches of
\eqref{app:tail:eq:D18Ipurepay} by $>1.05$.
Full bridges have internal new node and are foreign-sealed.

At fifteen the right $A_3$ admits middle-three, end-three/four/five.
Middle-three/end-four are affine or strong; end-five is the
actual sealed $H_4$, strong even at own zero.
Their full bridges are finite $D_{18}$ or $A_{18}$, internal
too far from the ends for a foreign $I$ arm.
The lone end-three has residual $5/8-Q$, numerator
$\gamma+1/4$. Positive $Q\ge1/4$ is strong, as are pure
owns $\gamma\ge1$. Own zero/three/four/six and parallel
are affine/null. Pure five has $G>3.74,R<15$, flat through
twelve, and is excluded by \eqref{app:tail:eq:D18Ipurepay}.
For $.9<\gamma<1$ use both actual $A_2$ cores with corrections
$\gamma(4n_{15}/3+2n_j/3)$, $j=14,16$.
Their average has
\[
 \alpha=\sqrt2(\gamma+1/4),\quad
 G>\sqrt{15-125/529}-.327>3.515,\quad
 R=15-{14\gamma^2/9-3/8\over2(\gamma+1/4)^2}<14.67.
\]
The saving increases from $177/529>.33$ and the radical
quotient increases to $8/75$. Inserting the full saving gives
$>1.0065$ in the strong-five branch and $>1.0152$ in the
weak branch. All feet are pure D and feasible against the
entire actual $I$ block by positive corrections.
Thus coverage through fifteen consists only of nonadjacent
pure left helpers, before any centre or strong-five floor.

Delete centre sixteen, leaving $A_{15}$, both forks and $I$.
The exact Schur data are
\[
 \Delta=1-q_{\rm left}-\delta^2-\epsilon^2-Q>0,
 \quad\sigma=\gamma+\eta j/16+(\delta+\epsilon)/2,
 \quad C_{16}=\sqrt{32}\,\sigma/\sqrt\Delta,
\]
and the sufficient ratio target is $<1.577$.
Finite classification gives the same complete structural centre
list as in Lemma~\ref{app:tail:lem:D17typed}, now with
$A_{15}$ and the actual $I$ foreign costs.
No-left fork pairs are again $(3,3),(3,4),(3,5)$.
A left-plus-both-forks bridge has residual $1/32$ and is
foreign-sealed; left-plus-fork-four also has residual $1/32$
and is strong at own zero. Left-plus-fork-three has residual
$9/32$ at $Q=0$, and may have the retained $I$-three extension
with $Q=1/4$; its smaller residual $1/32$ only strengthens
the affine/direct ratios. It is not falsely declared sealed.
With no fork the zero-own far-three would need
$1/2<Q<17/32$, containing no native $I$ square;
far-four/next-three need $Q<1/16,1/8$, so have time zero.
The other zero-own/isolated/foreign-cleared rows have the
actual local affine/foot proofs. Only the pure left helper
can remain weak without a fork.

For own zero and one fork, labels three/four/six have the
actual affine $E_8$/weighted-star/six--three principals.
Fork five has $Q=0$ or $1/4$. Its positive left $H_4$
projection, with $\alpha=\phi/(2\sqrt2)$, has
\[
 Y_D/\alpha=(1^{14},1/2,1/2,1/2,\pm1/2),\quad
 Y_D^2/\alpha^2=15,\quad \delta_5/\alpha^2=13-8\phi.
\]
At $Q=0$ its time ratio is $\sqrt{2+8\phi}$, giving
$G>3.62,R=15$, enough in \eqref{app:tail:eq:D18Ipurepay}.
At $Q=1/4$ retain the same point's $I$ coordinate rather
than clearing and reusing it. Its ambient normalized data are
\[
 G={\sqrt{6+6\phi}\over1+\sqrt{(13-8\phi)/(19-2\phi)}}>3.74,
 \quad R=19-2\phi<15.764,\quad d=\sqrt2/\phi<.875,
\]
where the pressure is $x.V_D+d y\ge G$, $V_D$ flat through
both early endpoints. If five is weak the bare ambient graft
already exceeds $1.057$. If strong, its bare value exceeds
$.995$ using the sharper $q_K,f_K$ above.
For this one mixed row form the ambient vector with classical
part $w+(V_D-P_K)/7+\lambda(e_1-P_5/5)$ and $I$ coordinate
$d/7+\nu$. Optimizing the unused nonnegative multipliers
gives the legal additional gain
\[
 (5/4)(t-b)_+^2+(s-d/7)_+^2
 >{[.788-b-(8/5)d/7]^2\over84/25}
 >{.227^2\over84/25}>.01533.
\]
This uses \eqref{app:tail:eq:I788} and includes the full
shifted $I$ cross. If $t\ge.6$ its clipped initial gain
already suffices. The mixed row is therefore excluded.

For own zero and any fork label at least seven, white
positivity forces $Q=0$. The positive actual core
$c+\delta^{-1}n_f+2(\delta^{-1}-\delta)n_{16}$ gives,
with $b=\delta^2\in(.81,1)$,
\[
 G=\sqrt{18-2/b}-\sqrt{-8+14/b-6/b^2},\quad
 R=26-16/b+6/b^2.
\]
We need the correlated continuous bound $14G-R>34.1904$,
not independent coarse bounds for $G,R$.
Set $z=1/b-7/6$. Then the second radicand is
$1/6-6z^2$ and $R=15.5-2z+6z^2$.
Concavity gives $\sqrt{18-2/b}\ge4.252-.252/b$;
the endpoint at $100/81$ has squared margin $1311/5062500$.
The tangent bound on the other root gives
\[
14G-R\ge39.912-14/\sqrt6-1.528z+(42\sqrt6-6)z^2
>34.1965-1.528z+96.8748z^2>34.1904.
\]
The last step completes a square, using
$1.528^2<.0061(4\cdot96.8748)$.
Equation~\eqref{app:tail:eq:D18Ipurepay} then exceeds
$1.000219,1.000313$ for strong five, and $1.008$ for weak
five. Every native label in this continuous interval is covered.

The remaining fork rows have the following actual pure feet:
\[
\begin{array}{c|c|cc}
\text{row}&\text{construction}&G\text{ lower}&R\text{ upper}\\\hline
\text{own }3,\text{ fork }4,5&\text{reciprocal }D_{18}&3.94&18\\
\text{own }5,\text{ fork }3&\text{left }H_4\text{ or cleared core}&3.64&16\\
\text{own }\ge7/\text{dot},\text{ fork }3&\text{special core}&3.62&16\\
\text{own }0,\text{ forks }3,4&\text{reciprocal }B_{18}&3.94&18\\
\text{own }0,\text{ forks }3,5&\text{finite }H_3&3.66&16
\end{array}
\]
For the first row, positive $Q$ is directly strong; at zero
the same cap has actual reciprocal cosh $8+9\delta>14$,
giving $G>\sqrt{18(13/15)}>3.94$.
For five/three, $Q=0$ or $1/4$ remains, with
$A^2=16\alpha^2-3/4+Q$, $\alpha^2=\phi+5/8$.
At $Q=1/4$ the actual core is
$c+\gamma(3n_{16}/2+n_{15}+n_{14}/2)+e/2$,
of square $\delta_5-1/4<0$; it is pure D after clearing.
Its time ratio exceeds $3.971$ and radical ratio is $<.322$,
so $G>3.64$; the zero-cost projected point is stronger.
For single-fork-three and $.9<\gamma<6/5$, $Q=0$ remains;
the actual special core has norm $-[1-3/(4\gamma^2)]$,
time ratio $>3.964$ and radical ratio $<.34$, hence $G>3.62$.
Positive $Q$ and larger owns are strong.
The two-fork three/four actual mate has cosh
$8+(9/2)\sqrt2>14$. The three/five $H_3$ projection has
$q=2\phi^{-6}<1/8$ and pressure
$\sqrt{16-q}/(1+\sqrt{q/16})>3.66$.
The affine triangles, stars, six--three paths, squares and
parallel nulls exclude all remaining small-own rows; positive
own with fork at least four has direct ratio
$\ge\phi/\sqrt2+1/2$.
Each displayed pure foot is flat through at least thirteen,
has positive actual correction coefficients and is feasible
at both actual signed forks and every unknown wall.
Its substitution in \eqref{app:tail:eq:D18Ipurepay} exceeds
one in both initial-five branches. This completes every
finite-white/foreign/own case at centre sixteen.

Coverage through sixteen is now exactly the pure left-helper
list, with actual nonadjacency. The prefix inverse entries
$2\min(i,j)$ hold through the centre; both actual forks are
retained positive old blocks in the mixed projections.
No spatial dual sum is used as a new root.
The same weak-five bootstrap as in Lemma~\ref{app:tail:lem:B18I}
therefore proves five strong: strong sixteen gives bare energy
$>1$, while strong fifteen gives $>.9851486$ and the capped
initial-$I$ gain $>.015$.
Choose the actual strong endpoint fifteen or sixteen; it
supplies $q>.9240775$, first-five value $b<.361$ and remaining
$P_5$ coefficient $>.146$.
Finally delete head one. Retained $D_{17}+I$ has isolated
or simple-three at old two, the latter optionally meeting
$I$ only by three. Other leaves give exceptional arms too
large, and four or a further fork is not finite.
Its raw equations are exactly those printed in the preceding
lemma, with every pure head point using actual D long roots
one/two/three. The actual signed forks remain unselected
acute walls. All conditions of the q-free
Lemma~\ref{lem:I-head-flat-reference} hold.
Its gain $>.07665$ with the one unused $I$ gives a total
$>1.0007275$. We do not need the older separate proof that
sixteen is individually strong.
\end{proof}

\begin{lemma}[The remaining ordinary and $B$ proper rows]
\label{app:tail:lem:proper19standard}
The forests $D_{17}+T_2$, $B_{17}+N_7/O_7$, and
$B_{18}/D_{18}+$ordinary $A_1$ are impossible. Here $T_2$
is ordinary; both actual typed defects are distinguished.
\end{lemma}
\begin{proof}
For $D_{17}+T_2$, the genuine bulk reaches twelve.
At thirteen the retained $D_4$ whitelist is exactly that in
Lemma~\ref{app:tail:lem:D17typed}, with a general ordinary
dihedral retained. Any positive foreign cost is at least $1/4$
and makes the short-leaf row strong. Pure finite owns at least
seven or dotted have ratio $>1.15\sqrt2>25/16$, above the
old target at $\sqrt{26}$.
The sole extra pure five has the actual left $H_4$ foot
$G>3.48,R<13$, flat through ten.
Pay $1/7$ at actual nine/ten, in either initial-five branch;
the pure $D$ energy exceeds $.9757$, incompatible with
$e_{T_2}>1/7$. Full bridges have internal new node and are
foreign-sealed. Hence only nonadjacent pure helpers remain
through thirteen. The \emph{proof} of the prefix-gain lemma
now applies with $M=12$: its actual weak-thirteen foot uses
only old one/two/twelve/thirteen/fourteen, all genuine long
roots. It gives $e_D>.8689525+.00496125=.87391375$;
the ordinary dihedral again excludes it. No D3 floor is used.

For $B_{17}$ with either typed defect, the bulk reaches thirteen
and the right $B_3$ layer fourteen is completed by its actual
long-three or sealed $F_4$ list. Thus pure-helper coverage
reaches fourteen before any bootstrap.
Weak five would give an actual strong thirteen/fourteen
endpoint and energy
$>.7600299945+(19-2\sqrt{78})/7>.95$;
the defect's $.09$ excludes it. Five is strong.
At fifteen the new $B_2$ short-three row is affine at
owns zero/four/six, strong for the other owns at least five,
and strong at own three with positive foreign cost.
Its sole extra pure own-three cap is
$v=(\omega_{15}+\omega_{17})/2$.
Release old sixteen: this same cap, with the entire defect
retained, completes finite $A_{17}$ and is the actual endpoint
by positive projection. Its cosh is
$(16+15\sqrt2)/\sqrt{18}>8.7$, giving $S_{16}>3.56$.
Choose actual strong $K=13$ or fourteen; their reference
energies are $>.88864,>.90715$, endpoint upper bounds
$3.175,3.311$ and available coefficients $>.14,>.13$.
Append the chord to $3.56$ at sixteen, skipping fifteen.
The slopes $<.13,1/8$ fit those individual coefficients.
The complete totals with the one defect exceed
$ .88864+.385^2/3+.09>1.028$ and
$ .90715+.249^2/2+.09>1.028$.
The extra cap is impossible; coverage now reaches fifteen.
The actual prefix-gain proof at $M=14$ gives
$e_B>.91211125$, again incompatible with $.09$.

For $B_{18}$ with a nonshort $A_1$, the $B_3$ layer fifteen
is the same sealed affine/strong list as in
Lemma~\ref{app:tail:lem:B18I}; it does not use a short-$I$
parameter. For $D_{18}$ distinguish the D4 long leaf at
fourteen, complete its bridges and isolate the pure own-five
fork row. Its actual foot has $G>3.62,R<14$, flat through
eleven. The old ten/eleven references give $e_D>.9957$,
so the ordinary $A_1$ cost $>.1$ excludes this extra.
At fifteen the right $A_3$ middle-three/end-four/five and
full bridge lists are those printed above. An ordinary-$A_1$
contact has cost at least $1/4$, strengthening the direct
ratios; full bridges have internal new node with long arms
and cannot acquire that foreign arm.
The pure end-three own-five has $G>3.74,R<15$; other
remaining $.9<\gamma<1$ owns use the two actual $A_2$
cores, flat through twelve, $G>3.5,R<15$.
The old eleven/twelve references give respectively $e_D>1.0094$
and $e_D>.9408$, so both are excluded by the ordinary $.1$.
Thus both $B_{18}$ and $D_{18}$ have complete pure-helper
coverage through fifteen in this spectator setting.
The prefix-gain proof at $M=14$, using actual existing long
neighbors rather than a false generic rank invocation, gives
$e_{B/D}>.91211125$. Adding the nonshort $A_1$ is impossible.
All pressures in this proof are pure in the classical block;
the named orthogonal spectator energy is counted once.
\end{proof}

\subsection{The complete dimension-nineteen forest cover}
\begin{proposition}\label{app:tail:prop:exact19cover}
If an actual farthest vertex in compact dimension nineteen exists,
its entire forest is one of $A_{19},B_{19},D_{19}$.
\end{proposition}
\begin{proof}
We supply the own-rank exhaustion, using current component bounds
to avoid an unnecessary succession of older weaker frontier tables.
As in dimension twenty a dominant component is $A_r$ with
$r\ge7$ or $B_r,D_r$ with $r\ge9$.
Each costs more than one half, so two cannot occur. Without
one, discard at most the single $I$ rank: the ordinary cost
is $>18/14$. A typed defect instead gives
$.09+17/14>1$; the funded $J$ pair is stronger than its
ordinary rank baseline. Hence exactly one dominant occurs.

An $A_7$ dominant is excluded by its half bound and remaining
ordinary ranks. At eight through twelve use $4/7$:
even discarding $I$ leaves six ordinary ranks at twelve;
a typed defect leaves five besides its $.09$. Both costs
strictly exceed the complementary budget.
At thirteen/fourteen the actual $5/7$ theorem and four
ordinary ranks suffice even after discarding $I$.
A typed defect leaves total ordinary rank at least three;
its cost $>17/70$ together with $.09$ closes this alternative.
For rank fifteen use $>.826$. Its ordinary rank-four
remainder is $>2/7$, an $I$ plus ordinary rank-three is
$>2/7$, and a typed defect plus an ordinary dihedral or
two ordinary $A_1$ costs $>.09+1/7$.
Every sum exceeds one. Rank sixteen has $>.86$;
the complete rank-three remainder possibilities are an
ordinary irreducible, ordinary dihedral plus $A_1$, three
$A_1$, $I+$dihedral, $I+$two $A_1$, $J+A_1$, or a typed
defect plus ordinary $A_1$.
Their weakest cost is $>.19$, so every sum again exceeds one.
This includes the four older residual A16 cases in one line.

At seventeen the ordinary irreducible dihedral costs $>1/7$,
which with $e_{A_{17}}>6/7$ suffices strictly. Two ordinary
$A_1$, $I+A_1$ and $J$ are stronger. Either actual typed
rank-two defect is excluded by Lemma~\ref{app:tail:lem:A17typed}.
At eighteen an ordinary $A_1$ is excluded by $.91+.1>1$;
the actual $I$ alternative by Lemma~\ref{app:tail:lem:A18I}.
Rank nineteen is the whole $A_{19}$ row.

For a $B/D$ dominant its $9/14$ theorem forces rank at least
fourteen: at thirteen the worst $I$ remainder leaves five
ordinary ranks, or a typed defect leaves four beside $.09$.
At fourteen/fifteen/sixteen the valid weaker $D$ prefix
floors are respectively
$>.7683275,>.797425,>.823735$.
The smallest sufficient spectator charges are respectively
$>2/7$, $>.09+1/7$, and $>.19$.
These account for all ordinary partitions, $I$, $J$, and
both distinct typed defects: the rank-five or rank-four
$I$ remainders use the proved joint $I+$rank-three or
ordinary-rank rates, not a fictitious ordinary cost for $I$.
Every row exceeds one.

At seventeen a two-$A_1$ or $I+A_1$ remainder costs $>1/5$,
already sufficient with the weaker prefix $>.8472225$;
$J$ is stronger. An ordinary dihedral is excluded for
$B_{17}$ by its genuine $M=12$ floor $>.8689525$, and for
$D_{17}$ by Lemma~\ref{app:tail:lem:proper19standard}.
Both typed defects are excluded for $B_{17}$ by that lemma
and for $D_{17}$ by Lemma~\ref{app:tail:lem:D17typed}.
At eighteen the remaining rank-one factor is ordinary or
$I$. The ordinary rows are excluded by
Lemma~\ref{app:tail:lem:proper19standard}; the actual $I$
rows by Lemmas~\ref{app:tail:lem:B18I} and
\ref{app:tail:lem:D18I}. Rank nineteen is whole.

Finite irreducible classification supplies no additional
exceptional types beyond the small components already charged
in the ordinary rate. Defect incompatibility forbids a second
defect; every surviving old rank-two seven without its specified
isolated opposite is ordinary. The rank partitions above are
therefore exhaustive and necessary, not assertions that the
three retained whole forests are realized.
\end{proof}

The final exact-nineteen conjunction uses this complete cover and
the three separate whole-case propositions
\ref{app:exact19:prop:wholeA19},
\ref{app:exact19:prop:wholeB19}, and
\ref{app:exact19:prop:wholeD19}.
Each concerns an actual whole forest at this same rank, with its
own complete head/end classification. No twenty-dimensional theorem
is decremented and no Coxeter structure is inherited by a face.
% The wholeD19 proposition is a separately authored PENDING proof home.
% This conjunction earns no acceptance until that complete body is reviewed.
% All source additions here remain uncompiled and require outside/root review.
% Exact19 wholeA/B/D consumers are owned by the separate consolidation author.
% All earlier coherent source freezes remain unchanged in their archives.

